\documentclass[11pt]{amsart}
\usepackage{amsmath,amssymb,amsthm,mathtools}
\usepackage{enumitem}
\usepackage[hidelinks,pdfauthor={Knowledge Sense},pdftitle={Exact counterexamples to the Barker--Larman problem in every dimension at least four},pdfsubject={Convex bodies with constant noncentral hyperplane sections},pdfkeywords={Barker--Larman problem, section volume, convex body of revolution, small divisors, Newton scheme}]{hyperref}
\usepackage{xurl}
\usepackage[capitalize,nameinlink]{cleveref}
\setlist{leftmargin=2em,itemsep=2pt,topsep=3pt}
\numberwithin{equation}{section}
\renewcommand{\thepart}{\Roman{part}}

\newtheorem{theorem}{Theorem}[section]
\newtheorem{proposition}[theorem]{Proposition}
\newtheorem{lemma}[theorem]{Lemma}
\newtheorem{corollary}[theorem]{Corollary}
\theoremstyle{definition}
\newtheorem{definition}[theorem]{Definition}
\theoremstyle{remark}
\newtheorem{remark}[theorem]{Remark}

\crefname{part}{Part}{Parts}

\newcommand{\R}{\mathbb{R}}
\newcommand{\C}{\mathbb{C}}
\newcommand{\Z}{\mathbb{Z}}
\newcommand{\N}{\mathbb{N}}
\newcommand{\Q}{\mathbb{Q}}
\newcommand{\eps}{\varepsilon}
\newcommand{\abs}[1]{\lvert #1\rvert}
\newcommand{\norm}[1]{\lVert #1\rVert}
\newcommand{\mean}{\operatorname{mean}}
\newcommand{\Lip}{\operatorname{Lip}}
\newcommand{\vol}{\operatorname{vol}}
\newcommand{\dist}{\operatorname{dist}}
\newcommand{\id}{\mathrm{id}}
\newcommand{\Leb}{\operatorname{Leb}}
\newcommand{\sgn}{\operatorname{sgn}}
\newcommand{\Log}{\operatorname{Log}}
\newcommand{\lb}{\lambda_{\mathrm b}}
\newcommand{\Wn}{\mathcal A}
\newcommand{\Disc}{\mathfrak D}
\newcommand{\cA}{\mathfrak A}
\newcommand{\cB}{\mathcal B}
\newcommand{\cG}{\mathcal G}
\newcommand{\cP}{\mathcal P}
\newcommand{\cJ}{\mathcal J}
\newcommand{\cC}{\mathcal C}
\newcommand{\cR}{\mathcal R}
\newcommand{\cF}{\mathcal F}
\newcommand{\cV}{\mathcal V}
\newcommand{\Dh}{\mathcal D_h}

\title[Counterexamples to the Barker--Larman problem]{Exact counterexamples to the Barker--Larman problem in every dimension at least four}
\author{Knowledge Sense}
\email{math@knowledgesense.jp}
\date{12 October 2026}
\subjclass[2020]{Primary 52A20; Secondary 52A38, 37E10, 11J86, 33C45}
\keywords{Barker--Larman problem, section volume, convex body of revolution, Gegenbauer polynomials, small divisors, Newton scheme, linear forms in logarithms}

\setlength{\emergencystretch}{3em}
\hyphenation{manuscript}
\begin{document}

\begin{abstract}
Let $K\subset\R^d$ be a convex body containing the unit ball $B^d$ in its interior such that every hyperplane tangent to $B^d$ cuts $K$ in a section of the same $(d-1)$-dimensional volume. The Barker--Larman problem asks whether $K$ must be a ball. We show that the answer is negative in every dimension $d\ge 4$. For each fixed $d\ge4$ we construct origin-symmetric, strictly convex bodies of revolution with real-analytic boundary which strictly contain $B^d$, are not balls, and have the same volume on all sections at distance one from the origin. In dimension four they form a sequence tending to the ball of radius two, and the set of attained volumes has positive measure near $4\sqrt3\,\pi$. In dimension $d\ge5$ they are parametrized by a compact set of frequencies of length $h-o_d(h)$ inside $[\omega_{0,d}+h,\omega_{0,d}+2h]$, as the restriction of a $C^1$ family of analytic profiles, and the common volumes fill a set of positive Lebesgue measure. The proofs combine finite-order formal solutions of the section equation with a Newton scheme in analytic Wiener algebras. Small divisors are controlled by circle arithmetic based on linear forms in logarithms, by a smooth diagonal elimination for real nonsymmetric matrices, and by the exclusion of a set of frequencies of measure $o_d(h)$.
\end{abstract}

\maketitle

\section{Introduction}\label{sec:intro}

\subsection{The problem}
Let $d\ge2$ and let $B^d$ be the closed Euclidean unit ball. Barker and Larman \cite{BL} asked the following question. Let $K$ and $L$ be convex bodies in $\R^d$ containing $B^d$ in their interiors and assume that
\[
\vol_{d-1}(K\cap H)=\vol_{d-1}(L\cap H)\qquad\text{for every hyperplane $H$ tangent to $B^d$.}
\]
Does it follow that $K=L$? The problem and several modifications are recorded in \cite{YZ}. In the plane the answer is affirmative when $L$ is a centered disk, a result of Santal\'o from 1951 (see \cite{HR}). For central sections the analogous statement for origin-symmetric bodies is a classical consequence of the injectivity of the spherical Radon transform; no such tool is available for sections at distance one. Taking for $L$ a ball centered at the origin one obtains the special case in which $K$ is required to have constant section volume over all tangent hyperplanes of $B^d$, and one asks whether $K$ is a ball.

The case $d=4$ of bodies of revolution was studied by Haddad and Ryabogin \cite{HR}, who proved rigidity under arithmetic conditions on the continued fraction expansion of a rotation number computed from the volume. Haddad \cite{Had} then constructed, for every $N$, approximate bodies $K_{\lambda,N}$ in $\R^4$ whose tangent sections have volume within $c(\abs{\lambda-\lb}/r_N)^{N+1}$ of $\lambda$ while the difference between outradius and inradius is at least $c\abs{\lambda-\lb}$, where $\lb=4\sqrt3\,\pi$ \cite[Theorem~1]{Had}. Convergence of the underlying formal series as $N\to\infty$ would give a negative answer to the problem in dimension four; this question is posed in \cite[Section~6.1]{Had}.

\subsection{Main results}
Throughout, $\kappa_k$ is the volume of the unit ball of $\R^k$ and $\abs{S^{k}}$ the area of the unit sphere $S^k\subset\R^{k+1}$. A body of revolution is a body invariant under all rotations fixing the first coordinate axis. For a body $K$ with $B^d\subset\operatorname{int}K$ and a number $V>0$ we say that $K$ has \emph{constant distance-one sections of volume $V$} if $\vol_{d-1}(K\cap H)=V$ for every hyperplane $H$ tangent to $B^d$.

\begin{theorem}[Exact sequence in dimension four]\label{thm:d4seq}
Put $\lb=4\sqrt3\,\pi$. There are an integer $m_0$ and, for every integer $m\ge m_0$, an origin-symmetric strictly convex body $K_m\subset\R^4$ of revolution with real-analytic boundary and a number $\lambda_m>0$ such that
\begin{enumerate}[label=\textup{(\roman*)}]
\item $B^4$ lies in the interior of $K_m$;
\item $K_m$ has constant distance-one sections of volume $\lambda_m$;
\item $K_m$ is not a Euclidean ball;
\item $K_m$ tends to the ball of radius $2$ in every fixed $C^j$ norm of the radial function, and $\lambda_m\to\lb$ as $m\to\infty$.
\end{enumerate}
\end{theorem}

\begin{theorem}[Positive measure of volumes in dimension four]\label{thm:d4meas}
There is $\eps_0>0$ such that for each sign $\pm$ and each $0<\eps<\eps_0$ there is a compact set $V_\eps^{\pm}\subset\R$ with
\[
\Leb(V_\eps^{\pm})\ \ge\ 18\pi^2\eps,\qquad \pm(\lambda-\lb)\in[\,40\pi^2\eps,\,100\pi^2\eps\,]\quad(\lambda\in V_\eps^\pm),
\]
such that every $\lambda\in V_\eps^\pm$ is the common distance-one section volume of a body $K$ with properties \textup{(i)}--\textup{(iii)} of \cref{thm:d4seq}. The bodies tend to the ball of radius two in every fixed $C^j$ norm of the radial function as $\eps\to0$, uniformly in $\lambda\in V_\eps^\pm$.
\end{theorem}

For $d\ge5$ put $\omega_{0,d}=2\arccos(1/\sqrt d)$, $I_{d,h}=[\omega_{0,d}+h,\omega_{0,d}+2h]$ and $q_d=32(d+1)$. For $\rho>0$ let $\Wn^{q}_{\rho}$ be the Banach space of real, odd, $\pi$-periodic functions $v(u)=\sum_kv_ke^{iku}$ with $\norm v_{\rho,q}=\sum_k(1+\abs k)^qe^{\rho\abs k}\abs{v_k}<\infty$. Every $v$ with $\norm v_{\rho,1}$ small defines, for every $\omega=2\theta$ near $\omega_{0,d}$, an origin-symmetric strictly convex body of revolution $K(v,\omega)\subset\R^d$ by the centered parametrization of \cref{def:centered} below; its radial function is real analytic and close to $\sec\theta$.

\begin{theorem}[Positive measure of volumes in dimension $d\ge5$]\label{thm:dge5}
Fix an integer $d\ge5$. There is $\rho_\infty=\rho_\infty(d)>0$ such that for every sufficiently small $h>0$ there are a compact set $E_{d,h}\subset I_{d,h}$ with $\Leb(E_{d,h})\ge h-o_d(h)$ and a $C^1$ map $\omega\mapsto v_\omega$ from $I_{d,h}$ to $\Wn^{q_d}_{\rho_\infty}$ with the following properties, where $K_\omega=K(v_\omega,\omega)$.
\begin{enumerate}[label=\textup{(\roman*)}]
\item For every $\omega\in I_{d,h}$, $K_\omega$ is an origin-symmetric, strictly convex body of revolution with real-analytic boundary, strictly contains $B^d$ and is not a ball.
\item For every $\omega\in E_{d,h}$, $K_\omega$ has constant distance-one sections, of volume $V(\omega)$.
\item The radial profiles of $K_\omega$ converge to the constant $\sqrt d$, uniformly in $\omega\in I_{d,h}$, on a fixed complex neighborhood of the real axis as $h\to0$.
\item Let $V(\omega)$ denote, for every $\omega\in I_{d,h}$, the mean of the distance-one section volume of $K_\omega$ over all unit normals. Then $V\in C^1(I_{d,h})$ and
\[
V'(\omega)=-2\,\abs{S^{d-2}}\,(d-1)^{(d-2)/2}+O_d(h).
\]
\end{enumerate}
Consequently $V(E_{d,h})$ has positive Lebesgue measure.
\end{theorem}

In particular $\omega\mapsto v_\omega$ restricted to $E_{d,h}$ is a Whitney $C^1$ family in $\Wn^{q_d}_{\rho_\infty}$. The section equation is asserted only on $E_{d,h}$; the extension to $I_{d,h}$ is a family of bodies with a $C^1$ mean volume, not a family of solutions. A set of exact volumes of positive measure is obtained; an interval of exact volumes is not asserted. The constants and smallness thresholds depend on the dimension.

\begin{corollary}[Negative answer to the Barker--Larman problem for $d\ge4$]\label{cor:BL}
Let $d\ge4$ and let $K$ be one of the bodies of \cref{thm:d4seq}, \cref{thm:d4meas} or \cref{thm:dge5}\textup{(ii)}, with common volume $V$. Let $R_V=\bigl(1+(V/\kappa_{d-1})^{2/(d-1)}\bigr)^{1/2}$. Then $K$ and the ball $R_VB^d$ are distinct convex bodies containing $B^d$ in their interiors which have equal volume on every hyperplane tangent to $B^d$.
\end{corollary}

\begin{proof}
A hyperplane at distance one from the center cuts $R_VB^d$ in a $(d-1)$-ball of radius $\sqrt{R_V^2-1}$, whose volume is $\kappa_{d-1}(R_V^2-1)^{(d-1)/2}=V$. Since $K$ is not a ball it differs from $R_VB^d$.
\end{proof}

Dimension three is not covered by these results. The bodies depend on the parameter and the distance is fixed; no statement is made about two fixed bodies and several distances.

The exact solutions constructed here do not arise by summing Haddad's formal series. In a companion paper \cite{KSdiv} we show that this series, in the normalization recalled in \cref{sec:formal}, has radius of convergence zero:

\begin{theorem}[{\cite{KSdiv}}]\label{thm:div}
Let $\lb=4\sqrt3\,\pi$, $z=(\lambda-\lb)/(\sqrt3\,\pi)$ and let $h(z,t)=t-\pi/3+\sqrt3\sum_{n\ge1}R_n(t)z^n$ be the normalized formal solution of the dimension-four section equation with $R_1=\tfrac1{48}+\tfrac1{12}\cos2t$. Then $\limsup_{n\to\infty}\norm{R_n}_{C^0}^{1/n}=\infty$. The same holds in every fixed $C^j$ and Sobolev $H^s$ norm, in the Fourier $\ell^1$ norm and in every positive-width analytic Wiener norm.
\end{theorem}

\subsection{Relation to earlier work}
We list what is used and what is new.
\begin{itemize}
\item From \cite{Had} we use the section equation in dimension four and the recursion producing finite trigonometric solutions to arbitrary finite order; we fix the order $N=12$ in \cref{sec:d4init}.
\item The rigidity theorem of \cite{HR} (Theorem~1, Definition~8 and Theorem~9 there) holds on a particular arithmetic set of volumes, with $A>8\kappa_3$ in their Definition~8. Our base volume $\lb=4\sqrt3\,\pi$ lies below $8\kappa_3=32\pi/3$. Moreover the rotation attached to the volume in \cite{HR} differs from the rotation $\omega$ of the endpoint map used here.
\item In \cite{FRSY} one fixed pair of bodies is considered with infinitely many distances having an interior accumulation point (Theorem~6.2 there), and Theorems~4.2--4.4 there need information over a distance interval. Here the body varies and the distance is fixed.
\item Lazutkin \cite{Laz} proved that a sufficiently smooth strictly convex planar billiard has a family of caustics near the boundary whose rotation numbers form a set of positive measure; his KAM-type method is a general ancestor of the small-divisor arguments below. The reduction of a variational equation to two successive differences and a scalar compatibility condition, the quasi-Newton terms proportional to the residual, and the use of finite-order initializers after fixing the analytic losses have direct antecedents in \cite{OAI1,OAI2}. The hypotheses of \cite{OAI1,OAI2} concern a full caustic collar of a single planar billiard, which is not supplied by the single-distance section equation for varying bodies. The geometry-specific forced-zero division, the scalar calculations and the retained-family inverse estimates are proved here.
\item The conjugacy of analytic circle maps with the quantitative a posteriori theorems of \cite{FHL} is used in dimension four.
\item Lower bounds for $\abs{\sin k\theta_0}$ and for the even-dimensional round symbols come from the coefficient-height consequence of the theorem of Baker and W\"ustholz \cite{BW}, recalled in \cref{lem:BW}.
\end{itemize}

\subsection{Outline of the proofs}
A body of revolution is described by its meridian profile. Parametrizing the profile by the angles of the tangent lines to $B^d$ through its points (the \emph{centered parametrization}, \cref{sec:geometry}) turns the constancy of the section volume into a functional equation for an odd lift $p$ of the circle, in which the unknown enters through two translates by $\pm\omega/2$, where $\omega$ is the rotation number of the endpoint map. This holds in every dimension.

In dimension four the base ball of radius two has rational rotation $\omega=2\pi/3$, and the linearization at the ball has infinitely many resonances. We use a finite-order formal solution, adjust the rotation to a nearby Diophantine number with continued fraction $[0;3,m,1,1,\dots]$ by an a posteriori conjugacy theorem, and solve the linearized equation by two successive cohomological equations with a forced zero. The scalar compatibility condition is solved by the volume and is nonzero because of a contribution from the third resonance. A Newton scheme with fixed width loss then gives exact solutions.

In dimension $d\ge5$ the ball of radius $\sqrt d$ has irrational rotation and a one-dimensional bifurcation in the degree-two zonal harmonic. A finite-order bifurcation expansion in the amplitude $\eps\sim h$ gives initializers to every finite order. The linearization, suitably normalized, has the form $D+A$ with a diagonal $D$ built from Gegenbauer values (odd $d$) or from cotangents (even $d$), and an off-diagonal part $A=O(h)$. A smooth diagonal elimination followed by a three-class analysis of the effective diagonals gives inverses on a set of frequencies of measure $h-o(h)$, with costs controlled in two regimes; this makes a Newton iteration with infinitely many steps converge. The surviving set of frequencies is nested and has a positive limiting threshold.

\subsection{Results used without proof}
Besides standard facts about Gegenbauer polynomials \cite{DLMF}, the proofs use the following results which are not reproved: the finite-order solutions of \cite{Had}; Theorems~3.1 and~4.1 of \cite{FHL} together with their proofs, with the parameter budget discussed in \cref{rem:fhl}; and the coefficient-height consequence of \cite{BW} recalled in \cref{lem:BW}. The companion result \cref{thm:div} is not used in the proofs of the existence theorems.

\subsection{Organization}
\Cref{sec:geometry} collects the geometry and function spaces common to all dimensions. \Cref{part:d4} treats $d=4$: Sections~\ref{sec:d4eq}--\ref{sec:d4proof} prove \cref{thm:d4seq} and \cref{sec:d4meas} proves \cref{thm:d4meas}. \Cref{part:d5} treats $d\ge5$: \cref{sec:sections} sets up the section operators, \cref{sec:resonances} treats resonances, arithmetic and finite initializers, \cref{sec:normalization} constructs the normalized operators for odd and even $d$, \cref{sec:schur} proves the high-mode inverse theorem, \cref{sec:center} treats the finite center and \cref{sec:newton} carries out the Newton iteration and proves \cref{thm:dge5}.

\section{Conventions and common geometry}\label{sec:geometry}

\subsection{Function spaces}
For a $2\pi$-periodic function $f=\sum_n\hat f_ne^{inu}$ and $\rho\ge0$ put
\[
\abs{f}_{\rho}=\sum_{n\in\Z}\abs{\hat f_n}e^{\rho\abs n},\qquad \abs{f}_{\rho,1}=\abs f_\rho+\abs{f'}_\rho .
\]
This is a Banach algebra norm which is invariant under real translations, and $\abs{f'}_{\rho-\eta}\le\eta^{-1}\abs f_\rho$ for $0<\eta\le\rho$. More generally $\norm f_{\rho,q}=\sum_n(1+\abs n)^q e^{\rho\abs n}\abs{\hat f_n}$; since $1+\abs{n+m}\le(1+\abs n)(1+\abs m)$, these norms are submultiplicative as well. Spaces of functions holomorphic in the interior of a strip (or a Bernstein ellipse) and continuous on its closure are equipped with the supremum norm. When we speak of the parity of a function on the circle (even, odd, $\pi$-periodic, $\pi$-antiperiodic) we mean its parity as a function of the angular variable in which it is written.

\subsection{Bodies of revolution}
Let $K\subset\R^d$ be a star body of revolution about the $x_1$-axis with radial function $r$. We write $r$ both for the function on $S^{d-1}$ and for the function $t\mapsto r(\cos t,\sin t,0,\dots,0)$ of the angle $t$ from $e_1$; $K$ is origin symmetric exactly when the latter is even and $\pi$-periodic.

\begin{lemma}[Analyticity at the poles]\label{lem:poles}
Let $r$ be even, $\pi$-periodic and holomorphic on a strip $\abs{\operatorname{Im}t}<\rho$. Then $r(u)=F(2u_1^2-1)$ for a function $F$ holomorphic on a Bernstein ellipse around $[-1,1]$; in particular the radial function is real analytic on all of $S^{d-1}$, including $u_1=\pm1$.
\end{lemma}

\begin{proof}
Write $r(t)=\sum_{n\ge0}a_n\cos 2nt$. Holomorphy on the strip gives $\abs{a_n}\le C_{\rho'}e^{-2n\rho'}$ for every $\rho'<\rho$. Since $\cos2nt=T_n(\cos2t)$ with the Chebyshev polynomials $T_n$, and $\abs{T_n}\le e^{2n\rho''}$ on the Bernstein ellipse with parameter $e^{2\rho''}$, the series $F(x)=\sum a_nT_n(x)$ converges uniformly on that ellipse for $\rho''<\rho'$. With $\cos2t=2u_1^2-1$ the claim follows.
\end{proof}

\begin{lemma}[Global strict convexity]\label{lem:convex}
For every $d\ge2$ and $R>1$ there is $\delta_1(d,R)>0$ with the following property. Let $r$ be a $C^2$ radial function on $S^{d-1}$ with $\norm{r-R}_{C^2(S^{d-1})}\le\delta_1$. Then $K=\{su:\,u\in S^{d-1},\ 0\le s\le r(u)\}$ is strictly convex, its boundary is a $C^2$ hypersurface \textup{(}real analytic if $r$ is\textup{)}, and $B^d\subset\operatorname{int}K$.
\end{lemma}

\begin{proof}
Put $\Phi(x)=\abs x^2/r(x/\abs x)^2$ for $x\ne0$ and $\Phi(0)=0$. On the unit sphere the Hessian of $\Phi$ depends continuously on the $2$-jet of $r$ and equals $2R^{-2}I$ for $r\equiv R$; hence for small $\delta_1$ it is at least $R^{-2}I$ there, and by homogeneity of degree two (so that the Hessian is homogeneous of degree zero) at every $x\ne0$. Also $\abs{\Phi(x)}\le C\abs x^2$, so $\Phi$ is $C^1$ at $0$ with zero gradient. On a line $s\mapsto a+sb$ the function $\phi(s)=\Phi(a+sb)$ is $C^1$, and $C^2$ with $\phi''\ge R^{-2}\abs b^2$ except at most at one point; hence $\phi'$ is strictly increasing and $\phi$ is strictly convex. Thus $\Phi$ is strictly convex, $K=\{\Phi\le1\}$ is strictly convex, and since $x\cdot\nabla\Phi=2$ on $\{\Phi=1\}$ the boundary is a $C^2$ (resp.\ analytic) hypersurface. Finally $\min r\ge R-\delta_1>1$.
\end{proof}

\subsection{The centered parametrization}
Suppose $r>1$ and write $r(t)=\sec\psi(t)$ with $0<\psi<\pi/2$. The tangent lines from the point $r(t)(\cos t,\sin t)$ of the meridian curve to the unit circle touch it at the points with normal angles $t\mp\psi(t)$.

\begin{definition}[Centered parametrization]\label{def:centered}
Fix $\theta\in(0,\pi/2)$ and put $\omega=2\theta$. Let $p(u)=u+v(u)$ with $v$ real analytic, odd and $\pi$-periodic, and $p'>0$. Define
\begin{gather*}
\chi(y)=\tfrac12\bigl[p(y-\theta)+p(y+\theta)\bigr],\qquad \alpha(y)=\tfrac12\bigl[p(y+\theta)-p(y-\theta)\bigr],\\
\hat r(y)=\sec\alpha(y),
\end{gather*}
and let $r$ be the function with $r(\chi(y))=\hat r(y)$. We write $K(v,\omega)$ for the body of revolution with radial function $r$.
\end{definition}

If $v=0$ then $\chi(y)=y$, $\alpha=\theta$ and $r\equiv\sec\theta=:R$ is the ball of radius $R$ centered at the origin. For general $p$ the pair $(\chi,\alpha)$ describes the profile point of polar angle $\chi(y)$ whose two tangent lines to the unit circle have normal angles $\chi-\alpha=p(y-\theta)$ and $\chi+\alpha=p(y+\theta)$. Consequently the chord cut from the meridian curve by the tangent line with normal angle $p(u)$ has end points with parameters $y=u-\theta$ and $y=u+\theta$: in the parameter $y$ the endpoint map is the rotation $y\mapsto y+\omega$. This holds at every $\theta$ with no arithmetic assumption.

\begin{lemma}\label{lem:centered}
Let $\mathcal J\subset(0,\pi/2)$ be a compact interval, $d\ge2$ and $0<\rho\le1$. There is $\delta_0=\delta_0(d,\mathcal J,\rho)>0$ such that the following holds for $\theta\in\mathcal J$ and $\abs v_{\rho,1}\le\delta_0$. The map $\chi$ is an analytic diffeomorphism of $\R/\pi\Z$ commuting with $y\mapsto-y$, it maps the strip $\abs{\operatorname{Im}y}<\rho/2$ biholomorphically onto a domain containing $\abs{\operatorname{Im}t}<\rho/4$, and $\alpha$ takes values in a compact subset of $(0,\pi/2)$ on $\abs{\operatorname{Im}y}\le\rho/2$, so that $\sec\alpha$ is unambiguous. The radial function $r$ is even, $\pi$-periodic, holomorphic on $\abs{\operatorname{Im}t}<\rho/4$ and satisfies $\norm{r-\sec\theta}_{C^2}\le C(\mathcal J,\rho)\abs v_{\rho,1}$; the body $K(v,2\theta)$ is origin symmetric, strictly convex, has analytic boundary and strictly contains $B^d$.
\end{lemma}

\begin{proof}
On $\abs{\operatorname{Im}y}\le\rho$ we have $\abs{\chi(y)-y}+\abs{\chi'(y)-1}\le2\abs v_{\rho,1}$ and $\abs{\alpha-\theta}\le\abs v_{\rho}$. Injectivity of $\chi$ on the convex strip of width $\rho/2$ follows from $\abs{\chi'-1}<\frac12$ by integration along segments, and the inverse on the strip of width $\rho/4$ is obtained from the contraction $y=t-(\chi(y)-y)$ on disks of radius $\rho/4$ contained in the larger strip; the Cauchy estimates on the remaining buffer bound its derivatives. Oddness and $\pi$-equivariance of $\chi$, and evenness and $\pi$-periodicity of $\alpha$, follow from those of $v$. The bound on $r-\sec\theta=\sec\alpha\circ\chi^{-1}-\sec\theta$ follows from these estimates and Cauchy's estimate on the buffer. The geometric statements follow from \cref{lem:poles} and \cref{lem:convex} with $R=\sec\theta$, after reducing $\delta_0$.
\end{proof}

\part{Dimension four}\label{part:d4}

\section{The section equation in \texorpdfstring{$\R^4$}{R4} and its finite-order solutions}\label{sec:d4eq}

Throughout \cref{part:d4} put $\alpha=\pi/3$, $\lb=4\sqrt3\,\pi$ and $w=1/16$.

\subsection{The equation and its converse}
A real increasing lift $h$ of the circle with $h(t+\pi)=h(t)+\pi$ and $h(t)-h(-t)=2t$ can be written $h(t)=t-\psi(t)$ with $\psi$ even and $\pi$-periodic. Assume that $\psi$ is close to $\alpha$ in an analytic norm. The planar curve with polar equation $t\mapsto\sec\psi(t)$ encloses the unit disk and is close to the circle of radius $2$; revolving it about the first axis gives a body $K\subset\R^4$. Put $a(x)=\tan\bigl(h^{-1}(x)-x\bigr)$ and
\begin{equation}\label{eq:W}
W(h)(x)=\pi\sin x\,\bigl[a(x)^2-a(-x)^2\bigr]+\tfrac{2\pi}{3}\cos x\,\bigl[a(x)^3+a(-x)^3\bigr].
\end{equation}
Let $A_K(x)$ be the three-dimensional volume of the section of $K$ by the hyperplane at distance one with unit normal $(\cos x,\sin x,0,0)$. In \cite[Section~2.2]{Had}, $A_K(x)$ is written as an explicit integral over the chord cut from the meridian curve by the tangent line, using the inverse-angle parametrization of \cite[Lemma~1]{Had}; our $\psi$ satisfies the assumptions of that lemma. Differentiating this integral gives \cite[equation~(6)]{Had}: for $0<x<\pi$,
\begin{equation}\label{eq:sinA}
\frac{d}{dx}\bigl[\sin x\,A_K(x)\bigr]=W(h)(x).
\end{equation}

\begin{proposition}[Equation and converse]\label{prop:d4eq}
Let $K$ be as above, with $\psi$ analytic and so close to $\alpha$ that $K$ is strictly convex, has analytic boundary and contains $B^4$ in its interior. Then $A_K\equiv\lambda$ on all hyperplanes tangent to $B^4$ if and only if $W(h)=\lambda\cos$ on $\R$.
\end{proposition}

\begin{proof}
If $A_K\equiv\lambda$ then \eqref{eq:sinA} gives $W(h)=\lambda\cos$ on $(0,\pi)$, and on $\R$ by analyticity. Conversely, if $W(h)=\lambda\cos$ then $\sin x\,[A_K(x)-\lambda]$ has zero derivative on $(0,\pi)$, so it is constant. Section volumes are bounded and continuous near $x=0$: for a smooth strictly convex body containing $B^4$ in its interior, nearby affine sections are uniformly transverse compact sections, so their radial functions in these hyperplanes depend continuously on the normal. Hence the constant is zero (let $x\downarrow0$), so $A_K=\lambda$ on $(0,\pi)$ and, by continuity, at both endpoints. Revolution symmetry covers all normals $u\in S^3$.
\end{proof}

\subsection{The formal branch and a finite initializer}\label{sec:formal}\label{sec:d4init}
Put $\lambda_z=\lb+\sqrt3\,\pi z$. The section equation $W(h)=\lambda_z\cos$ has a formal solution
\[
h_z(t)=t-\alpha+\sqrt3\sum_{j\ge1}R_j(t)z^j,
\]
where each $R_j$ is a real even $\pi$-periodic trigonometric polynomial, normalized by $R_1=\tfrac1{48}+\tfrac1{12}\cos2t$; this is the formal branch of \cite[Sections~3--5]{Had}. The first coefficients are
\begin{equation}\label{eq:R12}
R_1=\tfrac1{48}+\tfrac{\cos2t}{12},\qquad R_2=\tfrac3{256}-\tfrac{\cos2t}{144}+\tfrac{\cos4t}{72}.
\end{equation}
We use only a fixed finite order. Let
\begin{equation}\label{eq:hz}
h_z(t)=t-\alpha+f_z(t),\qquad f_z(t)=\sqrt3\sum_{j=1}^{12}R_j(t)z^j .
\end{equation}
By the construction of the $R_j$, the coefficients of $z^1,\dots,z^{12}$ in the expansion of $W(h_z)(x)-\lambda_z\cos x$ vanish; this is an identity between finite trigonometric polynomials for real $x$ and extends by analytic continuation to complex $x$.

\begin{lemma}[Common complex domain]\label{lem:d4domain}
There is $r_*>0$ such that for $\abs z<r_*$ and $\abs{\operatorname{Im}t}<4w$ one has $\abs{f_z(t)}\le w/2$ and $\abs{\partial_tf_z(t)}\le1/4$. For $\abs{\operatorname{Im}x}<3w$ the equation $t=x+\alpha-f_z(t)$ is a contraction on the closed disk $\abs{t-(x+\alpha)}\le w$, whose unique fixed point $t=h_z^{-1}(x)$ is holomorphic in $(x,z)$. After reducing $r_*$ the function $W(h_z)$ is jointly holomorphic on $\abs{\operatorname{Im}x}<2w$, and
\begin{equation}\label{eq:d4res}
W(h_z)-\lambda_z\cos=O(z^{13})
\end{equation}
uniformly on $\abs{\operatorname{Im}x}\le w$; every fixed finite number of derivatives satisfies the same bound on smaller fixed strips.
\end{lemma}

\begin{proof}
Since $f_z$ is a polynomial in $z$ with trigonometric polynomial coefficients vanishing at $z=0$, the bounds on $f_z$ hold for small $z$. The contraction maps the disk into the subdisk of radius $w/2$ and has Lipschitz constant at most $1/4$; the holomorphic implicit function theorem gives holomorphy in $(x,z)$, and periodicity identifies the local inverses as one inverse lift. This holds for complex $z$. Holomorphy of $W(h_z)$ requires the two tangent arguments to stay away from their poles, which holds after reducing $r_*$ since $h_z^{-1}(x)-x$ is close to $\alpha$. The first twelve Taylor coefficients in $z$ of $W(h_z)-\lambda_z\cos$ vanish, so Cauchy's formula on a smaller $z$-disk gives \eqref{eq:d4res}. The radius $r_*$ depends on the fixed order $12$.
\end{proof}

\section{Circle arithmetic and the approximate conjugacy}\label{sec:d4conj}

\subsection{Explicit arithmetic}
For every integer $m\ge3$ put
\begin{equation}\label{eq:thetam}
\theta_m=[0;3,m,1,1,1,\dots],\qquad \nu_m=-2\pi\theta_m,\qquad \gamma_m=\frac1{m+2}.
\end{equation}
These numbers are irrational. If $t_m=m+1/\phi$ with $\phi=(1+\sqrt5)/2$ then $\theta_m=t_m/(3t_m+1)$, so $\nu_m+\frac{2\pi}3=2\pi/[3(3t_m+1)]>0$ is comparable to $1/m$.

\begin{lemma}[Diophantine bound]\label{lem:dioph4}
For every nonzero integer $q$, $\norm{q\theta_m}\ge1/((m+2)\abs q)$, and consequently
\begin{equation}\label{eq:4gamma}
\abs{1-e^{iq\nu_m}}\ \ge\ \frac{4\gamma_m}{\abs q}.
\end{equation}
\end{lemma}

\begin{proof}
All partial quotients of $\theta_m$ are at most $m$, so the convergent denominators satisfy $q_{k+1}\le(m+1)q_k$. By best approximation and the convergent-error inequality, for $q_k\le\abs q<q_{k+1}$ one has $\norm{q\theta_m}\ge\norm{q_k\theta_m}>1/(q_k+q_{k+1})\ge1/[(m+2)q_k]\ge1/[(m+2)\abs q]$; the case $q=1$ follows directly. Then $\abs{1-e^{iq\nu_m}}=2\sin(\pi\norm{q\theta_m})\ge4\norm{q\theta_m}$ since $2\sin\pi x\ge4x$ for $0\le x\le\frac12$.
\end{proof}

In particular $\theta_m$ and $-\theta_m$ satisfy the Diophantine condition $D(\gamma_m,1)$ of \cite{FHL}. For the centered difference used below a Fourier mode $n$ has divisor $2i\sin(n\nu_m/2)$, whose modulus equals that in \eqref{eq:4gamma}; on $\pi$-periodic functions only even $n$ occur.

\subsection{The circle lift of the approximate body}
Put
\begin{equation}\label{eq:sigmaz}
\mathcal S_z(t)=-h_z^{-1}\bigl(-h_z(t)\bigr).
\end{equation}
It is analytic on $\abs{\operatorname{Im}t}<2w$, because $h_z$ maps this strip into $\abs{\operatorname{Im}x}<5w/2<3w$; it is real and increasing for real $z$, $\pi$-equivariant, and reversible under $t\mapsto-t$. Uniform bounds for its derivatives up to total order two on the closed strip of width $2w$ follow from the extra buffer in \cref{lem:d4domain} after reducing $r_*$. A direct substitution using only \eqref{eq:R12} gives, with
\begin{equation}\label{eq:kz}
\begin{gathered}
k_z(t)=t+\tfrac z{12}\sin2t-\tfrac{z^2}{48}(\sin2t+\sin4t),\\ \bar\nu(z)=-\tfrac{2\pi}3+\tfrac{\sqrt3}{24}z+\tfrac{11\sqrt3}{1152}z^2,\\
\mathcal S_z(k_z(t))-k_z(t+\bar\nu(z))=O(z^3).
\end{gathered}
\end{equation}
Explicitly $\mathcal S_z(t)=t-\tfrac{2\pi}3+zs_1(t)+z^2s_2(t)+O(z^3)$ with
\begin{align*}
s_1&=\tfrac{\sqrt3}{24}+\tfrac{\sqrt3}{24}\cos2t-\tfrac18\sin2t,\\
s_2&=\tfrac{5\sqrt3}{384}-\tfrac{\sqrt3}{72}\cos2t+\tfrac1{48}\sin2t+\tfrac{\sqrt3}{144}\cos4t+\tfrac1{24}\sin4t,
\end{align*}
and \eqref{eq:kz} is a finite algebraic identity after expanding every sine and cosine.

Since $\bar\nu'(0)=\sqrt3/24>0$, for all large $m$ there is a unique small $z_m>0$ with $\bar\nu(z_m)=\nu_m$. Then $z_m$, $\nu_m+2\pi/3$ and $\gamma_m$ are pairwise comparable with constants independent of $m$, and all later $O(z_m^j)$ bounds have constants independent of $m$.

\subsection{Preparation for the a posteriori conjugacy theorem}
We use the unit-circle conventions of \cite{FHL}:
\[
F_z(x)=\frac{\mathcal S_z(2\pi x)}{2\pi},\qquad H_z(x)=\frac{k_z(2\pi x)}{2\pi},\qquad \vartheta_m=\frac{\nu_m}{2\pi}.
\]
Take the fixed parameter interval $\mathcal A=(-r_*/2,r_*/2)$ and widths
\begin{equation}\label{eq:fhlwidths}
\hat\rho=w/\pi,\quad\rho=w/(2\pi),\quad\delta=\rho/8,\quad\rho_\infty=\rho/2 .
\end{equation}
Condition G1 of \cite[Theorem~3.1]{FHL} holds by the mixed derivative bounds above and G2 holds by \eqref{eq:4gamma}. For small $r_*$ the map $H_z$ has zero mean displacement, derivative within $1/8$ of $1$ and displacement at most $w/(8\pi)$; hence its image of the strip of width $\rho$ is at distance at least $3w/(8\pi)$ from the boundary of the strip of width $\hat\rho$, and the derivative and reciprocal-derivative bounds of H1 hold with $\sigma_1=\sigma_2=2$. The function of H2, which we denote by $\beta$, is
\[
\beta(x)=\frac{(\partial_zF_z)(H_z(x))}{H_z'(x+\vartheta_m)},\qquad \mean \beta\ \text{at }(z,\vartheta)=(0,-\tfrac13)\ \text{equals}\ \beta_0=\frac{\sqrt3}{48\pi},
\]
because the angular mean of $\partial_z\mathcal S_z$ at $z=0$ is $\sqrt3/24$. By continuity $\abs{\mean \beta-\beta_0}\le\beta_0/4$ after a fixed restriction, and we take $\sigma_b=2/\beta_0$. Also $z_m<r_*/4$ eventually, so $z_m$ has distance at least $r_*/4$ from $\partial\mathcal A$.

\begin{remark}[Uniformity as $\gamma_m\to0$ and the parameter budget]\label{rem:fhl}
Equation numbers in this remark refer to the 2019 author manuscript of \cite{FHL}. With the widths \eqref{eq:fhlwidths} the constants $a_1=2$, $a_2=2$, $a_3=8$ of \cite[(3.37)--(3.38)]{FHL} are fixed, and in \cite[(3.21), (3.22), (3.26)]{FHL}, with $\tau=1$, all terms containing $\gamma$ other than the displayed inverse powers carry nonnegative powers of $\gamma$. The iteration ratio $\kappa$ of \cite[(3.40)]{FHL} is $O(z_m^3/\gamma_m^2)=O(z_m)$, hence eventually at most $\frac12$, and the sums $\Sigma_{\kappa,p}$ of \cite[(3.41)]{FHL}, $p=0,1,2$, are then bounded by $2$. The positive denominators in \cite[(3.48)]{FHL} have the fixed margins above. In the parameter budget \cite[(3.43)]{FHL} and in the corresponding entry of \cite[(3.48)]{FHL} the numerator is written $\sigma_2\sigma_1$, while the single-step bound \cite[(3.23)]{FHL} is $\abs{\Delta_\alpha}<\sigma_b\sigma_2\norm e_\rho$. We therefore use the numerator $\sigma_2\max(\sigma_1,\sigma_b)$ in these two places; equivalently we sum the step bounds of \cite[Lemma~3.2]{FHL}:
\[
\sum_j\abs{\Delta z_j}\le\sigma_b\sigma_2\,\Sigma_{\kappa,2}\,\norm{e_0}_{\rho}\le2\sigma_b\sigma_2\norm{e_0}_{\rho}=O(z_m^3),
\]
which is eventually below $r_*/4$, so the parameter stays in $\mathcal A$ at every step.
\end{remark}

\begin{proposition}[Approximate conjugacy]\label{prop:d4conj}
For all sufficiently large $m$ there are a real number $z_m^*$ and an analytic, odd, $\pi$-equivariant diffeomorphism $k_m$ of the circle, holomorphic on the angular strip of width $w/2$, such that
\begin{equation}\label{eq:d4conj}
\begin{gathered}
z_m^*=z_m+O(z_m^3),\qquad k_m=k_{z_m}+O(z_m^2),\\
\mathcal S_{z_m^*}\circ k_m=k_m\circ T_{\nu_m},\qquad \mean(k_m-\id)=0,
\end{gathered}
\end{equation}
where $T_\nu(s)=s+\nu$. The parameter estimate is of the order of the initial error, and the conjugacy estimate of the order of the initial error divided by $\gamma_m$.
\end{proposition}

\begin{proof}
The initial error $e=\mathcal S_{z_m}\circ k_{z_m}-k_{z_m}\circ T_{\nu_m}$ is $O(z_m^3)$ by \eqref{eq:kz}. The hypotheses of \cite[Theorem~3.1]{FHL} and of the iteration in its proof (with the parameter budget of \cref{rem:fhl}) hold by the preceding paragraphs; part T1 of that theorem gives existence and part T2 gives $\abs{z^*_m-z_m}\le C_3\norm e=O(z_m^3)$ and $\norm{k_m-k_{z_m}}\le C_2\norm e/(\gamma\rho^\tau)=O(z_m^2)$.

A normalized conjugacy to an irrational rotation is unique for its fixed circle map: the quotient of two such conjugacies commutes with the irrational rotation and is therefore a translation by density of orbits, and a phase shift $c$ adds $c$ to the mean displacement, so the normalization fixes it. Applying this to $-k_m(-s)$, using the reversibility of $\mathcal S_z$, proves oddness; applying it to $k_m(s+\pi)-\pi$ proves $\pi$-equivariance.
\end{proof}

Define
\begin{equation}\label{eq:Gm}
G_m(y)=h_{z_m^*}\bigl(k_m(y-\nu_m/2)\bigr).
\end{equation}
The conjugacy equation and $h_z(-t)=h_z(t)-2t$ give
\begin{equation}\label{eq:Gsym}
\begin{gathered}
G_m(-y)=-G_m(y),\qquad G_m(y+\pi)=G_m(y)+\pi,\\
k_m(s)=\tfrac12\bigl[G_m(s-\tfrac{\nu_m}2)+G_m(s+\tfrac{\nu_m}2)\bigr].
\end{gathered}
\end{equation}
For instance $h_z(-k_m(s+\nu_m))=-h_z(k_m(s))$ gives oddness after setting $y=s+\nu_m/2$. Expanding \eqref{eq:Gm} with \eqref{eq:R12}, \eqref{eq:kz} and \eqref{eq:d4conj},
\begin{equation}\label{eq:Gexp}
G_m(y)=y-\tfrac{z_m}6\sin2y+O(z_m^2).
\end{equation}
All these maps are holomorphic with uniform estimates on the fixed strip of width $w/4$, because $k_m$ is $O(z_m)$-close to the identity on the strip of width $w/2$.

\subsection{Transfer to Wiener norms}
Set
\begin{equation}\label{eq:rho0}
\rho_0=w/8=\tfrac1{128}.
\end{equation}
If a $2\pi$-periodic function has supremum norm $M_0$ on the strip of width $w/4$, its Fourier coefficients obey $\abs{\hat f_n}\le M_0e^{-\abs nw/4}$, so its Wiener norm at $\rho_0$ is at most $M_0\sum_ne^{-\abs nw/8}$; multiplying by $\abs n^j$ gives the same transfer for every fixed $j$. Consequently \eqref{eq:Gexp}, including any fixed number of derivatives, holds in the norm $\abs\cdot_{\rho_0}$. The residual \eqref{eq:d4res}, composed with $G_m$, is $O(z_m^{13})$ on the strip of width $w/4$ after a fixed reduction of $r_*$; the parameter $z^*_m$ is comparable to $z_m$, and the volume parameter is $\lambda_{z_m^*}$. The prepared residual is therefore $O(z_m^{13})$ in the norm $\abs\cdot_{\rho_0}$.

\section{The fixed-shift Newton scheme}\label{sec:d4newton}

In this section $\varsigma=10^{-5}$, $\rho_0=1/128$ as in \eqref{eq:rho0}, and $b_*=1/(96\pi)$.

\subsection{The equation and a bounded class}
Let $D_\nu Y=Y-Y(\cdot+\nu)$ and $\Delta_\nu v=v(\cdot+\nu/2)-v(\cdot-\nu/2)$. We consider real parameters with
\begin{equation}\label{eq:class}
\begin{gathered}
0<\gamma\le1,\qquad \abs{\nu+\tfrac{2\pi}3}\le\varsigma,\qquad \abs{\lambda-\lb}\le1,\qquad \tfrac{\rho_0}2\le\rho\le\rho_0,\\
G(y)=y+g(y),\qquad g\ \text{real, odd, $\pi$-periodic},\qquad \abs g_{\rho,1}\le\varsigma,
\end{gathered}
\end{equation}
and assume \eqref{eq:4gamma}, i.e.\ $\abs{1-e^{iq\nu}}\ge4\gamma/\abs q$ for $q\ne0$. Put
\begin{align}
A_-(y)&=\tfrac12[G(y-\nu)-G(y)],\qquad A_+(y)=\tfrac12[G(y)-G(y+\nu)],\notag\\
a_\pm&=\tan A_\pm,\notag\\
E(G,\lambda)&=\pi\sin G\,(a_-^2-a_+^2)+\tfrac{2\pi}3\cos G\,(a_-^3+a_+^3)-\lambda\cos G.\label{eq:E}
\end{align}
The differences of lifts in $A_\pm$ are periodic with constant term $-\nu/2$, within $\varsigma/2$ of $\alpha$, and remaining Wiener norm at most $\varsigma$, so their distance in the algebra from $\alpha$ is at most $3\varsigma/2$. Hence $\cos A_\pm$ differs from $1/2$ by at most $3\varsigma$ and is invertible by a Neumann series with inverse of norm $<3$; no tangent pole occurs on the complex strip. The map $E$ is analytic in $(g,\lambda)$ with uniformly bounded second derivative as a map from the Wiener algebra to itself throughout \eqref{eq:class}, and on any straight segment inside the class; its value is even and $\pi$-antiperiodic. The constants are independent of $\gamma,\nu,\rho$.

For $G$ in the class put
\begin{equation}\label{eq:khG}
k(s)=\tfrac12\bigl[G(s-\tfrac\nu2)+G(s+\tfrac\nu2)\bigr],\qquad h(k(s))=G(s+\tfrac\nu2).
\end{equation}
Real derivatives of $G$ and $k$ are positive, $k$ is odd and $\pi$-equivariant, $h(t)-h(-t)=2t$ and $-h^{-1}(-h(k(s)))=k(s+\nu)$. At $x=G(y)$ one has $h^{-1}(x)-x=A_-(y)$ and $h^{-1}(-x)-(-x)=A_+(y)$, so the two tangent arguments in \eqref{eq:W} are exactly $A_-(y)$ and $A_+(y)$. This proves:

\begin{lemma}\label{lem:Esubst}
For $G$ in the class \eqref{eq:class}, $E(G,\lambda)(y)=\bigl[W(h)-\lambda\cos\bigr](G(y))$.
\end{lemma}

The substitution is exact, and every shift in \eqref{eq:E} is fixed rather than depending on the unknown.

\subsection{The forced-zero multiplier and the derivative identity}
Put
\begin{equation}\label{eq:P}
\begin{gathered}
q(s)=\tfrac12[G(s-\tfrac\nu2)-G(s+\tfrac\nu2)],\\
P(s)=\tfrac12G'(s-\tfrac\nu2)G'(s+\tfrac\nu2)\,\tan q(s)\sec^3q(s)\,\sin k(s).
\end{gathered}
\end{equation}
This fixes the holomorphic expression without ambiguity of a fractional power. $P$ is odd and $\pi$-antiperiodic. If $k=s+\mathfrak k$, then $\mathfrak k$ is odd and $\pi$-periodic with $\abs{\mathfrak k}_{\rho,1}\le\varsigma$.

\begin{lemma}[Division by $\sin$]\label{lem:sinquot}
For an odd sine series $f$ one has $\abs{f/\sin}_\rho\le\abs{f'}_\rho$. Moreover $P(s)=\sin s\,B(s)$ with $\abs{B-4\sqrt3}_\rho<\tfrac1{10}$, $B^{-1}$ exists with norm $<1$, and for odd $\pi$-antiperiodic $Y$
\begin{equation}\label{eq:YP}
\abs{Y/P}_{\rho-\eta}\le\eta^{-1}\abs Y_\rho\qquad(0<\eta\le1).
\end{equation}
\end{lemma}

\begin{proof}
For $n\ge1$ the Laurent polynomial $\sin(ns)/\sin s$ has $n$ coefficients of modulus one with exponents between $-(n-1)$ and $n-1$; this gives the first claim. By $\sin(s+\mathfrak k)=\sin s\cos\mathfrak k+\cos s\sin\mathfrak k$,
\[
\frac{\sin k(s)}{\sin s}=\cos\mathfrak k+\cos s\,\frac{\mathfrak k}{\sin s}\,\operatorname{sinc}\mathfrak k ,
\]
which is close to $1$ in the same algebra without loss of width: $\abs{\cos s}_\rho=\cosh\rho<2$, $\abs{\cos\mathfrak k-1}\le\varsigma^2$, $\abs{\operatorname{sinc}\mathfrak k}\le2$, so its difference from $1$ is at most $5\varsigma$. Next $\abs{\tan q\sec^3q-8\sqrt3}<1/20$: the quantities $\sin q-\sin\alpha$ and $\cos q-\cos\alpha$ are at most $3\varsigma$ in norm, the inverse of $\cos q$ has norm $<3$, and along the joining segment the derivative of $w\mapsto w^{-4}$ is bounded by $4\cdot3^5$, whence the difference is at most $3\varsigma\cdot3^4+4\cdot3^5\cdot3\varsigma<1/20$. Expanding $G'_-G'_+$ as $1$ plus a term of norm $\le2\varsigma+\varsigma^2$ and multiplying these bounds gives $\abs{B-4\sqrt3}<1/10$. The Neumann series gives $B^{-1}$. Thus $P$ has exactly the zeros of $\sin s$ on the working strip, all simple, and \eqref{eq:YP} follows from the first claim and the Cauchy estimate.
\end{proof}

\begin{lemma}[Derivative identity]\label{lem:DE}
Let $v$ be odd and $\pi$-periodic and $Y=P\,\Delta_\nu v$. Then
\begin{multline}\label{eq:DE}
DE(G,\lambda)[G'v,\delta\lambda](y)\\
=-\frac{2\pi}{G'(y)}\bigl[Y(y-\tfrac\nu2)-Y(y+\tfrac\nu2)\bigr]+E'(y)v(y)-\delta\lambda\cos G(y).
\end{multline}
\end{lemma}

\begin{proof}
Direct differentiation of \eqref{eq:E}. With $X=G(y)$, $A=A_-(y)$, $D=A_+(y)$, the partial derivatives of $E$ with respect to the two neighboring values $G(y\mp\nu)$ are $\pi\tan A\sec^3A\sin(X+A)$ and $\pi\tan D\sec^3D\sin(X-D)$. Subtracting $E'(y)v(y)$ the central derivative cancels, and the neighboring coefficients times $v(y\pm\nu)-v(y)$ are exactly the two terms built from \eqref{eq:P}.
\end{proof}

\subsection{Two differences, one scalar}
For the current $E$ define
\begin{equation}\label{eq:Q}
Q_0(s)=\frac{G'(s+\frac\nu2)E(s+\frac\nu2)}{2\pi},\qquad Q_\lambda(s)=-\frac{G'(s+\frac\nu2)\cos G(s+\frac\nu2)}{2\pi}.
\end{equation}
Both are $\pi$-antiperiodic, hence of zero mean, and satisfy $Q(-s-\nu)=Q(s)$. Their unique zero-mean solutions $Y_0,Y_\lambda$ of $D_\nu Y=Q$ are odd and $\pi$-antiperiodic (apply $s\mapsto-s$ and change sign; the reflection identity preserves the equation and uniqueness gives oddness). Define the scalar
\begin{equation}\label{eq:Lambda}
\Lambda(G,\nu)=\mean(Y_\lambda/P),
\end{equation}
and assume $\abs\Lambda\ge b_*$. Put
\begin{equation}\label{eq:corr}
\begin{gathered}
\delta\lambda=-\frac{\mean(Y_0/P)}{\Lambda},\qquad H=\frac{Y_0+\delta\lambda\,Y_\lambda}{P},\\
\Delta_\nu v=H,\qquad\mean v=0,\qquad\delta G=G'v.
\end{gathered}
\end{equation}
The numerator of $H$ is odd and $\pi$-antiperiodic and vanishes at the zeros of $\sin$, so $H$ is even and $\pi$-periodic and by construction of zero mean; its centered-difference inverse is odd and $\pi$-periodic. Everything preserves reality. Substituting \eqref{eq:Q}--\eqref{eq:corr} in \eqref{eq:DE} gives
\begin{equation}\label{eq:newtoneq}
DE[\delta G,\delta\lambda]=-E+E'v,
\end{equation}
and there is no other compatibility condition in these symmetry spaces.

\subsection{Estimates}
Put $\rho_*=\rho_0/4$, let $\rho$ range over $[\rho_0/2,\rho_0]$, and put $e=\abs E_\rho$. The zero-mean inverses of $D_\nu$ and $\Delta_\nu$ cost $C\gamma^{-1}\eta^{-1}$ in passing from $\rho$ to $\rho-\eta$: divide each Fourier coefficient by $1-e^{in\nu}$, resp.\ $2i\sin(n\nu/2)$, and use \eqref{eq:4gamma} and $\sup_{n\ge1}ne^{-\eta n}\le\eta^{-1}$. Using in succession the cohomological inverse from $\rho$ to $\rho-\eta$, \eqref{eq:YP}, and the second cohomological inverse, there are constants $C_p,C_v,C_D,C_B,C_L\ge1$ independent of $m,\gamma$ and the iterate such that for $0<\eta\le\min(1,\rho/4)$:
\begin{align}
\abs{\delta\lambda}&\le C_p\gamma^{-1}e,\label{eq:est1}\\
\abs v_{\rho-3\eta}+\abs{\delta G}_{\rho-3\eta}&\le C_v\gamma^{-3}\eta^{-3}e,\label{eq:est2}\\
\abs{\delta G}_{\rho-4\eta,1}&\le C_D\gamma^{-3}\eta^{-4}e,\label{eq:est2b}\\
\abs{E(G+\delta G,\lambda+\delta\lambda)}_{\rho-4\eta}&\le C_B\gamma^{-6}\eta^{-6}e^2,\label{eq:est3}
\end{align}
whenever the correction segment stays in \eqref{eq:class}. In \eqref{eq:est2b} the whole product $G'v$ is differentiated on one further strip, so no bound on $G''$ on the original boundary is needed. For \eqref{eq:est3}, Taylor's formula for the analytic map \eqref{eq:E}, \eqref{eq:newtoneq} and the Cauchy estimate $\abs{E'}_{\rho-\eta}\le\eta^{-1}e$ show that the Taylor term is bounded by a constant times $(\abs{\delta G}+\abs{\delta\lambda})^2$, while the retained term $E'v$ is at most $C\gamma^{-3}\eta^{-4}e^2$; both are bounded by the right side of \eqref{eq:est3}.

\emph{Scalar stability.} For two maps $G$ in the class on any strip of width $\ge\rho_0/2$, compare $Q_\lambda$ on $\rho_*$ in the Wiener norm; the difference is $\le C\abs{G_{\mathrm{new}}-G}_{\rho_*,1}$. Apply the cohomological inverse from $\rho_*$ to $\rho_*/2$ and division by $\sin$ from $\rho_*/2$ to $\rho_*/4$. The multiplier $B$ of \cref{lem:sinquot} changes by at most $C\abs{G_{\mathrm{new}}-G}_{\rho_*,1}$, since the quotient of the variation of $\mathfrak k$ by $\sin$ is controlled by the first part of \cref{lem:sinquot} and all other factors are analytic functions of $g,g'$. Expanding the difference of $(Y_\lambda/\sin)B^{-1}$ gives
\begin{equation}\label{eq:est4}
\abs{\Lambda(G_{\mathrm{new}},\nu)-\Lambda(G,\nu)}\le C_L\gamma^{-1}\abs{G_{\mathrm{new}}-G}_{\rho_*,1},
\end{equation}
since $Y_\lambda/\sin$ is itself $O(\gamma^{-1})$ on fixed smaller strips. All widths here are fixed independently of the iteration and strictly below $\rho_0/2$.

\subsection{The Newton lemma}
Put $A_D=C_D(16/\rho_0)^4/(1-2^{-6})$ and $A_p=C_p/(1-2^{-10})$, and let $c_*>0$ satisfy
\begin{equation}\label{eq:c}
c_*\le\frac{2^{-10}}{C_B(16/\rho_0)^6},\qquad c_*\le\frac\varsigma{2A_D},\qquad c_*\le\frac {b_*}{C_LA_D},\qquad c_*\le\frac1{2A_p}.
\end{equation}

\begin{lemma}[Quantified Newton lemma]\label{lem:newton4}
Assume \eqref{eq:4gamma} with $0<\gamma\le1$ and $\abs{\nu+2\pi/3}\le\varsigma$. Let $g_0$ be real, odd, $\pi$-periodic and $\lambda_0$ real with
\begin{equation}\label{eq:hyp4}
\begin{gathered}
\abs{g_0}_{\rho_0,1}\le\varsigma/2,\qquad\abs{\lambda_0-\lb}\le\tfrac12,\\
\abs{\Lambda(G_0,\nu)}\ge2b_*,\qquad e_0=\abs{E(G_0,\lambda_0)}_{\rho_0}\le c_*\gamma^6 .
\end{gathered}
\end{equation}
Then the corrections \eqref{eq:corr} converge on $\rho_0/2$ to real $G_\infty=\id+g_\infty$ and $\lambda_\infty$ with $E(G_\infty,\lambda_\infty)=0$, with preserved symmetries, $\abs{g_\infty}_{\rho_0/2,1}\le\varsigma$ and
\begin{equation}\label{eq:newton4concl}
\abs{G_\infty-G_0}_{\rho_0/2,1}\le A_D\gamma^{-3}e_0,\qquad\abs{\lambda_\infty-\lambda_0}\le A_p\gamma^{-1}e_0.
\end{equation}
\end{lemma}

\begin{proof}
Put $\rho_j=\rho_0/2+\rho_0/2^{j+1}$ and $\eta_j=(\rho_j-\rho_{j+1})/4=\rho_02^{-j-4}$; then $\eta_j\le1$ and $\eta_j<\rho_j/4$. We show inductively $e_j\le e_02^{-10j}$, together with the class conditions along every correction segment, $\abs{\lambda_j-\lb}\le1$ and $\abs{\Lambda(G_j,\nu)}\ge b_*$. By \eqref{eq:est3} the ratio of the bound for $e_{j+1}$ to $e_02^{-10(j+1)}$ is at most $C_B(16/\rho_0)^6\gamma^{-6}e_02^{10-4j}\le1$ by the first inequality in \eqref{eq:c}. The total $C^1$ correction on each later strip is at most $C_D(16/\rho_0)^4\gamma^{-3}e_0\sum_j2^{-6j}=A_D\gamma^{-3}e_0\le A_Dc_*\le\varsigma/2$; the total parameter change is at most $A_p\gamma^{-1}e_0\le A_pc_*\le\frac12$; and the total scalar change, by \eqref{eq:est4}, is at most $C_LA_D\gamma^{-4}e_0\le C_LA_Dc_*\le b_*$. Thus all conditions persist, and the estimates used for the next step are justified. The corrections converge absolutely in the $C^1$ Wiener norm on $\rho_0/2$ and in the parameter; continuity of $E$ on that space and $e_j\to0$ give $E=0$; summing proves \eqref{eq:newton4concl}.
\end{proof}

The constant $c_*$ in \eqref{eq:c} depends only on the class \eqref{eq:class}; it is independent of $m$, $\gamma$ and the iterate. The lemma is applied along a sequence for which $e_0/\gamma^6\to0$.

\section{Proof of the exact sequence theorem}\label{sec:d4proof}

\subsection{The nonvanishing scalar}
For the prepared $G_m$ of \eqref{eq:Gm}, the expansion \eqref{eq:Gexp} in a fixed Wiener norm, together with its derivative, gives
\[
\sin G_m(y)=\bigl(1-\tfrac{z_m}{12}\bigr)\sin y-\tfrac{z_m}{12}\sin3y+O(z_m^2).
\]
Since
\begin{gather*}
Q_\lambda(s)=-\frac{(\sin G_m)'(s+\nu_m/2)}{2\pi},\\
D_\nu\sin(ns)=-2\sin(n\nu/2)\cos\bigl(n(s+\nu/2)\bigr),
\end{gather*}
the first-order terms give the sine coefficients of $Y_\lambda$
\begin{equation}\label{eq:Y13}
Y_1=\frac{1-z_m/12}{4\pi\sin(\nu_m/2)},\qquad Y_3=-\frac{z_m}{16\pi\sin(3\nu_m/2)} .
\end{equation}
As $m\to\infty$, $\sin(3\nu_m/2)=-\tfrac{\sqrt3}{16}z_m+O(z_m^2)$, whence
\begin{equation}\label{eq:Y13lim}
Y_1\to-\frac1{2\pi\sqrt3},\qquad Y_3\to\frac1{\pi\sqrt3}.
\end{equation}
The entire remaining tail of $Q_\lambda$ is $O(z_m^2)$ in the initial Wiener norm. On any fixed smaller strip its cohomological inverse and the subsequent division by $\sin$ cost at most $C\gamma_m^{-1}$, so the whole contribution of the tail is $O(z_m^2/\gamma_m)=O(z_m)$; this uses $\gamma_m\ge c_*'z_m$ and estimates all Fourier modes, including all higher near-resonances. Further, $B=P/\sin=4\sqrt3+O(z_m)$ with inverse converging in the Wiener norm on fixed smaller strips, and $\mean(\sin s/\sin s)=\mean(\sin3s/\sin s)=1$. Combining,
\begin{equation}\label{eq:Lam13}
\Lambda(G_m,\nu_m)=\frac1{4\sqrt3}\Bigl(-\frac1{2\pi\sqrt3}+\frac1{\pi\sqrt3}\Bigr)+O(z_m)=\frac1{24\pi}+O(z_m),
\end{equation}
so eventually $\abs\Lambda\ge1/(48\pi)=2b_*$, which is the condition required in \eqref{eq:hyp4}. The frequency-three coefficient contributes to the leading value: it comes from the first-order term $-\frac{z_m}6\sin2y$ of $G_m$, divided by the small denominator $\sin(3\nu_m/2)$ of order $z_m$.

\subsection{Application of the Newton lemma}
Take $G_0=G_m$, $\nu=\nu_m$, $\gamma=\gamma_m$, $\lambda_0=\lambda_{z_m^*}$. By \eqref{eq:Gexp}, \eqref{eq:rho0}, \eqref{eq:Lam13} and $z_m\to0$ all smallness and scalar assumptions of \eqref{eq:hyp4} hold eventually; by \eqref{eq:d4res} and \cref{sec:d4conj}, $e_0=O(z_m^{13})$, while $\gamma_m$ is comparable to $z_m$. Hence $e_0/\gamma_m^6=O(z_m^7)\to0$ and \cref{lem:newton4} applies for all $m$ above a common $m_0$. Write the exact pair as $(G_m^\infty,\lambda_m)$. On the strip $\rho_0/2$,
\begin{equation}\label{eq:11-1}
G_m^\infty-G_m=O(z_m^{10})\ \text{in $C^1$ Wiener norm},\qquad\lambda_m-\lambda_{z_m^*}=O(z_m^{12}).
\end{equation}
Shrinking to a smaller fixed strip gives the same order for every fixed number of derivatives. The maps $k$ reconstructed via \eqref{eq:khG} differ by $O(z_m^{10})$ from $k_m$, their derivatives are bounded away from zero, and their inverses on the real circle and the compositions differ by $O(z_m^{10})$ in every fixed $C^j$ norm (use one additional fixed analytic strip for derivatives). Consequently the reconstructed meridian lift $h_m$ differs from $h_{z_m^*}$ by $O(z_m^{10})$ in real $C^j$; in particular
\begin{equation}\label{eq:hm}
h_m(t)=t-\alpha+\sqrt3\,z_m\bigl[\tfrac1{48}+\tfrac{\cos2t}{12}\bigr]+O(z_m^2).
\end{equation}
The $\cos2t$ coefficient is $\sqrt3z_m/12+O(z_m^2)$, so $\psi_m=t-h_m$ is nonconstant for all large $m$ and tends to $\alpha$ in every fixed $C^j$ norm. The chosen branch has $0<\psi_m<\pi/2$, so $\sec$ is injective on its range, and the radial function is nonconstant as well.

\subsection{Global geometry}
The analytic inverse function theorem on a fixed smaller strip applied to $k$ in \eqref{eq:khG} shows that $h_m$ and $\psi_m$ are real analytic, with $\psi_m$ even and $\pi$-periodic, with uniform holomorphic bounds on a fixed strip and convergence to $\alpha$. Thus the radial function $r_m(t)=\sec\psi_m(t)$ is even, $\pi$-periodic, analytic and close to $2$. By \cref{lem:poles} it is real analytic on all of $S^3$, and $r_m\to2$ in every fixed $C^j(S^3)$ norm. Let $K_m=\{su:u\in S^3,0\le s\le r_m(u)\}$. By \cref{lem:convex}, $K_m$ is strictly convex with analytic boundary and $B^4$ lies in its interior. Origin symmetry implies that a Euclidean ball with this radial function would be centered at the origin, and nonconstancy in \eqref{eq:hm} excludes it.

\subsection{Conclusion}
By \cref{lem:Esubst}, $E(G_m^\infty,\lambda_m)=0$ implies $W(h_m)=\lambda_m\cos$ on $\R$. \Cref{prop:d4eq}, applied to the strictly convex analytic body $K_m$, gives that all distance-one section volumes equal $\lambda_m$. Positivity of $\lambda_m$ and $\lambda_m\to\lb$ follow from \eqref{eq:11-1} and the definition of $\lambda_z$. This proves \cref{thm:d4seq}. \qed

\section{Positive measure of the volume set in dimension four}\label{sec:d4meas}

In this section the rotation number varies in a whole interval, so the single-frequency scheme of \cref{sec:d4newton} is carried out for families indexed by a compact set of rotation numbers, with weighted Lipschitz norms. We keep the notation of \cref{part:d4}: $\eps$ is the width of the rotation interval, $\theta$ the unit-circle rotation, $\nu=2\pi\theta$, $\rho$ an angular complex strip width, $z$ the parameter of the finite body, and $\lambda$ the actual section volume. The rotation $\theta$ of the endpoint map is different from the angle that \cite{HR} computes from $\lambda$.

\subsection{A compact frequency set of positive measure}
Let $\theta_0=-1/3$ and $I_\eps^+=[\theta_0+\eps,\theta_0+2\eps]$ or $I_\eps^-=[\theta_0-2\eps,\theta_0-\eps]$, with $\eps$ small enough that the chosen interval lies in $(-1,0)$. Put $\gamma=\eps/16$ and
\[
\Theta_\eps=I_\eps\cap\{\theta:\ \norm{n\theta}\ge\gamma/n^2\ \text{for all }n\ge1\}.
\]
This set is compact and contains only irrationals. For each $n$ the excluded subset of one unit interval consists of $n$ circle intervals of length $2\gamma/n^3$, so has measure at most $2\gamma/n^2$; countable subadditivity gives excluded measure at most $(\pi^2/3)\gamma<\eps/4$. Hence
\begin{equation}\label{eq:Theta}
\Leb(\Theta_\eps)\ge\tfrac{3\eps}4,\qquad\abs{1-e^{in\nu}}\ge\frac{4\gamma}{n^2}\quad(n\ne0,\ \nu=2\pi\theta\in2\pi\Theta_\eps).
\end{equation}
The measure of the set of actual volumes is obtained from this in \cref{sec:d4noncollapse}.

\subsection{A Lipschitz family of approximate conjugacies}
Use the finite order-$12$ initializer $h_z$ and $\lambda_z=\lb+\sqrt3\pi z$ of \cref{sec:d4init}. The geometric residual $R_z=W(h_z)-\lambda_z\cos$ is holomorphic on a fixed strip with $R_z=O(z^{13})$, $\partial_zR_z=O(z^{12})$, and \eqref{eq:kz} gives $k_z$ and $\bar\nu(z)$. For $\theta\in I_\eps$ let $z_0(\theta)$ be the real inverse of $\bar\nu(z)=2\pi\theta$. It is analytic on an interval containing both intervals $I_\eps^\pm$, $z_0=O(\eps)$, $\Lip(z_0)=O(1)$ and
\begin{equation}\label{eq:z0prime}
z_0'(\theta)=\frac{48\pi}{\sqrt3}+O(\eps).
\end{equation}
All constants below are uniform in $\eps$ after fixing the sign; one may take the maximum over the two signs. The initial circle error $e(\theta)=\mathcal S_{z_0}\circ k_{z_0}-k_{z_0}\circ T_{2\pi\theta}$ has supremum norm $O(\eps^3)$ and Lipschitz constant $O(\eps^2)$ on fixed smaller strips, as it is an analytic function of $z$ vanishing through order two and $z_0$ has bounded derivative.

We apply \cite[Theorem~4.1]{FHL} and its proof to $F_z(x)=\mathcal S_z(2\pi x)/(2\pi)$, $H_\theta(x)=k_{z_0(\theta)}(2\pi x)/(2\pi)$ and the parameter $z_0(\theta)$, with the fixed widths of \eqref{eq:fhlwidths}. Here the strip loss $\delta$ of \cite{FHL} is fixed and unrelated to $\eps$. Condition G1 asks for family derivatives through order three, supplied by the finite holomorphic family and the inverse-domain buffers. Condition G2 holds with $\tau=2$ and $\gamma=\eps/16$ by \eqref{eq:Theta}. Condition H1 has uniformly bounded first derivative, reciprocal and second derivative with fixed slack. Condition H2 has base mean parameter derivative $\sqrt3/(48\pi)>0$ and the same slack as in \cref{sec:d4conj}. Condition H3 holds with fixed bounds strictly above $\Lip(H-\id)$, $\Lip(\partial_xH)$ and $\Lip(z_0)$, all $O(1)$. The parameter stays in the middle half of the fixed interval and the image margins are unchanged. The weighted initial error $\max\{\sup\abs e,\gamma\delta^{\tau+1}\Lip(e)\}$ is $O(\eps^3)$. The constants of the proof of \cite[Theorem~4.1]{FHL} are bounded uniformly as $\gamma\to0$: after their explicit inverse powers, the quantities in \cite[(4.9)--(4.13), (4.20), (4.31), (4.34), (4.37)]{FHL} carry nonnegative powers of $\gamma$, and the fixed H3 slacks keep \cite[(4.39)]{FHL} bounded. Both iteration ratios of that proof are $O(\eps^3/\gamma^2)=O(\eps)$, so the sums occurring there, including the negative-index sums in \cite[(4.33) and (4.35)]{FHL}, are bounded by fixed constants for small $\eps$. The parameter budget is treated as in \cref{rem:fhl}.

Thus there are a Lipschitz family of exact normalized conjugacies $k_\theta$ and parameters $z^*(\theta)$ with, on fixed angular strips,
\begin{equation}\label{eq:d4fam}
\begin{gathered}
\sup\abs{z^*-z_0}=O(\eps^3),\qquad \Lip(z^*-z_0)=O(\eps^2),\\
\sup\norm{k-k_{z_0}}=O(\eps^2),\qquad\Lip(k-k_{z_0})=O(\eps).
\end{gathered}
\end{equation}
The same orders hold for any fixed number of spatial derivatives after a further fixed strip restriction. For the parameter Lipschitz bound one uses \cite[(4.40)]{FHL}, whose correction is of the order of the weighted error divided by $\gamma$, i.e.\ $O(\eps^2)$; for the conjugacy one uses \cite[(4.33), (4.35)]{FHL}, whose correction is of the order of the weighted error divided by $\gamma^2$, i.e.\ $O(\eps)$. Normalization, uniqueness of irrational conjugacies, reversibility and $\pi$-equivariance give the same odd, $\pi$-equivariant geometry as in \cref{prop:d4conj} for every $\theta$. Put
\[
G_\theta(y)=h_{z^*(\theta)}\bigl(k_\theta(y-\nu/2)\bigr),\qquad\lambda_{\mathrm{pre}}(\theta)=\lb+\sqrt3\pi z^*(\theta).
\]
Then $G_\theta$ is odd and $\pi$-equivariant, reconstructs $h_{z^*}$ exactly, and on a common complex strip
\begin{equation}\label{eq:Gfam}
G_\theta=\id-\tfrac{z_0}6\sin2y+\mathcal E_\theta,\qquad\sup\abs{\mathcal E}=O(\eps^2),\quad\Lip(\mathcal E)=O(\eps).
\end{equation}
All fixed-derivative versions hold on smaller strips. Suprema and Lipschitz seminorms transfer to Wiener norms by the same fixed Fourier coefficient buffer as in \cref{sec:d4conj}, applied also to difference quotients. No differentiability of the exact family in $\theta$ is needed.

\subsection{Weighted algebra and a shift-stable class}
For a family $f_\theta$ of periodic functions on a compact set $\Theta_\eps$ write
\[
\Lip_\rho(f)=\sup_{\theta\ne\theta'}\frac{\abs{f_\theta-f_{\theta'}}_\rho}{\abs{\theta-\theta'}},\qquad N_\rho(f)=\sup_\theta\abs{f_\theta}_\rho+\eps\Lip_\rho(f).
\]
Thus $\Lip_\rho$ is the Lipschitz seminorm of $\theta\mapsto f_\theta$ with values in the Wiener space, not the Lipschitz seminorm of the scalar function $\theta\mapsto\abs{f_\theta}_\rho$. In the same way $\Lip_\Theta(g;\abs\cdot_\rho+\abs{\cdot'}_\rho)$ is the supremum of $(\abs{g_\theta-g_{\theta'}}_\rho+\abs{g'_\theta-g'_{\theta'}}_\rho)/\abs{\theta-\theta'}$. The norm $N$ is submultiplicative, since the two product Lipschitz terms are covered by the product of the two $N$ norms, constant scalar families use the same definition, and Cauchy's estimate gives $N_{\rho-\eta}(f')\le\eta^{-1}N_\rho(f)$. For $g=G-\id$ use the stronger mixed control
\begin{equation}\label{eq:S}
S_\rho(g)=\sup_\theta\bigl(\abs g_\rho+\abs{g'}_\rho+\abs{g''}_\rho\bigr)+\eps\,\Lip_\Theta\bigl(g;\abs\cdot_\rho+\abs{\cdot'}_\rho\bigr),
\end{equation}
that is, supremum $C^2$ plus weighted Lipschitz $C^1$. For a fixed $l$ with $\abs l\le1$, $\eps\Lip[g'_\theta(\cdot+l\nu(\theta))]\le\eps\Lip(g')+2\pi\abs l\eps\sup\abs{g''}$, and the analogous bound for $g$ itself needs only $g'$; hence translated $g'$ factors and their differences are controlled by $S$, with constants independent of $\eps\le1$. The affine part of $G$ contributes only the bounded scalar slope $2\pi l$, separated from periodic functions. For a correction $u$,
\begin{equation}\label{eq:shiftstable}
N(u(\cdot+l\nu))\le N(u)+2\pi\abs l\,\eps\sup\abs{u'}\le C\,S(u).
\end{equation}
Take the core width $\rho_0=1/128$ and a small fixed $\varsigma_1>0$ so that $S(g)\le\varsigma_1$ implies the class \eqref{eq:class}. Require also $\eps\le\varsigma_1/(4\pi)$, $N(\lambda-\lb)\le1$ and $\abs{\nu+2\pi/3}\le\varsigma$. The small initial family satisfies these with arbitrarily large relative slack as $\eps\to0$.

For $k=\frac12(G(\cdot-\frac\nu2)+G(\cdot+\frac\nu2))$ and $q=\frac12(G(\cdot-\frac\nu2)-G(\cdot+\frac\nu2))$ define $P$ exactly as in \eqref{eq:P}, so that $P=\sin(s)\,B$. Pointwise in $\theta$, \cref{lem:sinquot} gives $\abs{B-4\sqrt3}<\frac1{10}$ and $B^{-1}$ bounded. In the weighted norm the same expressions have fixed bounds under $S(g)\le\varsigma_1$: the removable division of $k-\id$ by $\sin$ uses the weighted $C^1$ bound just established, and all other products and analytic functions use submultiplicativity. Thus $N(B),N(B^{-1})\le C$ uniformly in $\eps$, $\theta$ and the iterate. The reciprocal Lipschitz bound follows from $B_\theta^{-1}-B_{\theta'}^{-1}=B_\theta^{-1}(B_{\theta'}-B_\theta)B_{\theta'}^{-1}$. Similarly $N(B_{\mathrm{new}}-B_{\mathrm{old}})\le CS(G_{\mathrm{new}}-G_{\mathrm{old}})$, with constants depending only on the fixed strip and class.

\subsection{Weighted difference inverses}
Consider the zero-mean Fourier inverse of either $D_\nu f=f-f(\cdot+\nu)$ or the centered difference $f(\cdot+\frac\nu2)-f(\cdot-\frac\nu2)$ (absorb the half-shift into the first multiplier). In each case the multiplier $\mathfrak m_n(\theta)$ satisfies
\begin{equation}\label{eq:mult}
\abs{\mathfrak m_n}\le C\gamma^{-1}n^2,\qquad\frac{\abs{\mathfrak m_n(\theta)-\mathfrak m_n(\theta')}}{\abs{\theta-\theta'}}\le C\gamma^{-2}n^5\qquad(\theta,\theta'\in\Theta_\eps).
\end{equation}
For the first inverse subtract the two endpoint fractions: the numerator exponential difference is $\le2\pi n\abs{\theta-\theta'}$ and both endpoint denominators obey \eqref{eq:Theta}; the half-shift adds only a term $\le C\gamma^{-1}n^3$; the centered sine multiplier has the same endpoint-product estimate. Only the two endpoints $\theta,\theta'\in\Theta_\eps$ are used. Since $\eps/\gamma=16$, for $0<\eta\le1$ the Fourier sum and $\sup_tt^je^{-\eta t}\le C_j\eta^{-j}$ give
\begin{equation}\label{eq:wsum}
N_{\rho-\eta}(\text{inverse of }f)\le C\gamma^{-1}\eta^{-5}N_\rho(f).
\end{equation}
For an odd $\pi$-antiperiodic family, by the coefficient identity for $\sin(ns)/\sin s$,
\begin{equation}\label{eq:wsin}
N_{\rho-\eta}(f/\sin)\le\eta^{-1}N_\rho(f).
\end{equation}
All frequencies, including arbitrarily high near-resonances, are covered.

\subsection{Weighted scalar and one quasi-Newton step}
The equation and the derivative identity are those of \cref{sec:d4newton}, separately for every $\theta$, and preserve the odd, $\pi$-equivariant shape. Let $\mathcal C_\nu$ denote the first inverse composed with its half-shift. Then $Y_0=\mathcal C_\nu(G'E/(2\pi))$, $Y_\lambda=\mathcal C_\nu(-G'\cos G/(2\pi))$ and $\Lambda=\mean(Y_\lambda/P)$. Use fixed inner strips below $\rho_0/2$ to define and estimate scalar means; these gaps do not accumulate. Assume initially $\inf\abs\Lambda\ge2b_*$ with $b_*=1/(96\pi)$ and $N(\Lambda)\le L_\Lambda$ for a fixed $L_\Lambda$; we maintain $\inf\abs\Lambda\ge b_*$ and $N(\Lambda)\le L_\Lambda+b_*$, so $N(1/\Lambda)\le1/b_*+(L_\Lambda+b_*)/b_*^2$ is fixed. The correction is \eqref{eq:corr}. All zero means and reflection symmetries hold exactly as before, being algebraic identities at each $\theta$.

Let $e=N_\rho(E)$ with $\rho_0/2\le\rho\le\rho_0$. By \eqref{eq:wsum} and \eqref{eq:wsin} with fixed gaps, $N(\delta\lambda)\le C_p\gamma^{-1}e$. For the shape correction pay three successive losses $\eta$ (first inverse $\gamma^{-1}\eta^{-5}$, sine division $\eta^{-1}$, second inverse $\gamma^{-1}\eta^{-5}$), with one extra $\gamma^{-1}$ in the $\delta\lambda\,Y_\lambda$ term; hence $N_{\rho-3\eta}(\delta G)\le C\gamma^{-3}\eta^{-11}e$. Two further losses $\eta$ supply the supremum second derivative and the weighted Lipschitz first derivative:
\begin{equation}\label{eq:wcorr}
S_{\rho-5\eta}(\delta G)\le C_D\gamma^{-3}\eta^{-13}e.
\end{equation}
The principal derivative cancels $E$ exactly and leaves $E'v$, with weighted bound $\le C\gamma^{-3}\eta^{-12}e^2$. For the Taylor remainder, $E$ is an analytic function of three translated values of $G$ and of $\lambda$ with coefficient families of bounded weighted norms on the input class; by \eqref{eq:shiftstable}, translated correction norms are bounded by $S(\delta G)$ and the scalar correction by $N(\delta\lambda)$, so the quadratic remainder is $\le C[S(\delta G)+N(\delta\lambda)]^2$. A conservative bound for the new residual is
\begin{equation}\label{eq:wres}
N_{\rho-5\eta}(E_{\mathrm{new}})\le C_B\gamma^{-6}\eta^{-26}e^2.
\end{equation}
Finally, comparing two families $G$ at the same $\theta$, using fixed inner strips, \eqref{eq:wsum}, \eqref{eq:wsin} and the variation bound for $B$,
\begin{equation}\label{eq:wLam}
N(\Lambda_{\mathrm{new}}-\Lambda_{\mathrm{old}})\le C_L\gamma^{-1}S_{\rho_*}(\delta G),\qquad\rho_*=\rho_0/4 ,
\end{equation}
where $Y_\lambda/\sin$ has norm $O(\gamma^{-1})$ on fixed gaps. This gives the weighted scalar stability, including the reciprocal Lipschitz control.

\subsection{A strip induction for the weighted scheme}
Put $\rho_j=\rho_0/2+\rho_0/2^{j+1}$ and $\eta_j=(\rho_0/20)2^{-j}$, so that $\rho_j-\rho_{j+1}=5\eta_j$; enlarge $C_p,C_D,C_B,C_L$ to at least one, and set $A_D=C_D(20/\rho_0)^{13}/(1-2^{-51})$, $A_p=C_p/(1-2^{-64})$. Choose a fixed $c_*>0$ with
\begin{equation}\label{eq:c2}
c_*\le\frac{2^{-64}}{C_B(20/\rho_0)^{26}},\quad c_*\le\frac{\varsigma_1}{2A_D},\quad c_*\le\frac1{2A_p},\quad c_*\le\frac {b_*}{C_LA_D}.
\end{equation}

\begin{lemma}\label{lem:newton4w}
Assume $S(g_0)\le\varsigma_1/2$, $N(\lambda_0-\lb)\le\frac12$, $\inf\abs{\Lambda_0}\ge2b_*$, $N(\Lambda_0)\le L_\Lambda$ and $e_0\le c_*\gamma^6$. Then the corrections converge in supremum $C^2$ plus weighted Lipschitz $C^1$ on $\rho_0/2$, $\lambda$ converges in its weighted Lipschitz norm, $E=0$ in the limit for every $\theta$, and
\begin{equation}\label{eq:wlimit}
N(\lambda_\infty-\lambda_0)\le A_p\gamma^{-1}e_0,\qquad\Lip(\lambda_\infty-\lambda_0)\le A_p\gamma^{-1}e_0/\eps.
\end{equation}
\end{lemma}

\begin{proof}
By \eqref{eq:wres} one shows inductively $e_j\le e_02^{-64j}$: the ratio to the desired bound at $j+1$ is at most $C_B(20/\rho_0)^{26}\gamma^{-6}e_02^{64-38j}\le1$. The shape corrections sum to $\le A_D\gamma^{-3}e_0$, the parameter corrections to $\le A_p\gamma^{-1}e_0$, and the scalar drift is $\le C_LA_D\gamma^{-4}e_0$. Since $\gamma\le1$, \eqref{eq:c2} maintains all three margins and $N(\Lambda)\le L_\Lambda+b_*$ throughout, which proves the assumptions before each step and along every correction segment. The remaining statements follow as in \cref{lem:newton4}.
\end{proof}

\subsection{Uniform initialization with the entire weighted resonant tail}
Equation \eqref{eq:Gfam} is valid in the weighted norm: the remainder is $O(\eps^2)$, since its supremum is $O(\eps^2)$ and $\eps$ times its Lipschitz constant is $O(\eps^2)$; the same holds after fixed spatial derivatives. The scalar computation of \cref{sec:d4proof} is repeated in this norm. With $G=\id-(z_0/6)\sin2y+O_N(\eps^2)$ and $\nu=\bar\nu(z_0)$, the two leading sine coefficients of $\sin G$ sit at frequencies $1$ and $3$, the third being $-z_0/12+O_N(\eps^2)$. Solving the first difference gives for $Y_\lambda$ the coefficient $na_n/(4\pi\sin(n\nu/2))$, where $a_n$ is the $n$-th sine coefficient of $\sin G$. Since $\sin(3\nu/2)=-\tfrac{\sqrt3}{16}z_0+O(z_0^2)$, the third leading quotient tends to $1/(\sqrt3\pi)$ and the first to $-1/(2\sqrt3\pi)$. The smooth ratios $z_0/\sin(3\nu/2)$ have weighted errors $O(\eps)$ uniformly on each interval, including their sign, and their difference quotients are bounded by analytic cancellation at $z_0=0$. The entire differentiated remainder has weighted norm $O(\eps^2)$, so by \eqref{eq:wsum}, \eqref{eq:wsin} across fixed gaps its total scalar contribution is $O(\eps^2/\gamma)=O(\eps)$, which controls every higher resonance. Since moreover $B=4\sqrt3+O_N(\eps)$,
\begin{equation}\label{eq:wLam0}
N\bigl(\Lambda-\tfrac1{24\pi}\bigr)=O(\eps).
\end{equation}
Hence $\inf\abs\Lambda\ge2b_*$ and $N(\Lambda)\le L_\Lambda$ with a fixed $L_\Lambda$ for small $\eps$; the weighted bound \eqref{eq:wLam0} also controls the Lipschitz seminorm of $1/\Lambda$.

The prepared geometric residual is $E_{\mathrm{pre}}(\theta)=R_{z^*(\theta)}\circ G_\theta$. On fixed strips $R_z=O(z^{13})$, $\partial_zR_z=O(z^{12})$ and the spatial derivative is $O(z^{13})$; by \eqref{eq:d4fam}, $\Lip(z^*)=O(1)$ and $\Lip(G)=O(1)$. Therefore $\sup\abs{E_{\mathrm{pre}}}=O(\eps^{13})$, $\Lip(E_{\mathrm{pre}})=O(\eps^{12})$, and after the fixed supremum-to-Wiener buffers
\begin{equation}\label{eq:e0}
e_0=N(E_{\mathrm{pre}})=O(\eps^{13}).
\end{equation}
Since $\gamma=\eps/16$, $e_0/\gamma^6=O(\eps^7)\to0$. Also $S(g_0)=O(\eps)$ and $N(\lambda_{\mathrm{pre}}-\lb)=O(\eps)$, so all initial margins hold for small $\eps$, and \cref{lem:newton4w} applies uniformly to every $\theta\in\Theta_\eps$.

\subsection{Noncollapse of the actual volume map}\label{sec:d4noncollapse}
The prepared volume is $\lambda_{\mathrm{pre}}=\lb+\sqrt3\pi z^*$. By \eqref{eq:z0prime}, $\frac d{d\theta}[\lb+\sqrt3\pi z_0(\theta)]=48\pi^2+O(\eps)$ on the whole interval $I_\eps$, so its lower Lipschitz constant on $\Theta_\eps$ is also $48\pi^2+O(\eps)$. The Lipschitz contribution of $\sqrt3\pi(z^*-z_0)$ is $O(\eps^2)$ by \eqref{eq:d4fam}, and the weighted Newton estimate \eqref{eq:wlimit}, \eqref{eq:e0} gives $\Lip(\lambda_\infty-\lambda_{\mathrm{pre}})=O(\eps^{11})$. Hence for small $\eps$ and distinct $\theta,\theta'\in\Theta_\eps$,
\begin{equation}\label{eq:noncollapse}
\abs{\lambda_\infty(\theta)-\lambda_\infty(\theta')}\ge24\pi^2\abs{\theta-\theta'}.
\end{equation}
The map $\lambda_\infty$ is Lipschitz and strictly increasing on this compact set, since both perturbations have smaller Lipschitz constant than the increasing base map. Its range $V_\eps$ is compact, and applying the Lipschitz bound to its inverse gives
\begin{equation}\label{eq:Vmeasure}
\Leb(V_\eps)\ge24\pi^2\Leb(\Theta_\eps)\ge18\pi^2\eps>0.
\end{equation}
The leading expansion $\lambda_\infty(\theta)-\lb=48\pi^2(\theta-\theta_0)+O(\eps^2)$ also shows that $V_\eps$ lies on the chosen side of $\lb$ at distances comparable to $\eps$, namely $\pm(\lambda-\lb)\in[40\pi^2\eps,100\pi^2\eps]$ for small $\eps$.

\subsection{Proof of the positive-measure theorem in dimension four}
The fixed-strip shape correction is $O(\eps^{10})$, while the prepared meridian is the finite $h_{z^*}=t-\pi/3+\sqrt3z_0R_1+O(\eps^2)$. The $\cos2t$ coefficient is $\sqrt3z_0/12+O(\eps^2)$, uniformly nonzero because $\abs{z_0}$ is bounded below by a fixed multiple of $\eps$ on either interval; so no output is a ball. The reconstruction, \cref{lem:poles}, \cref{lem:convex} and the integrated section-volume identity of \cref{prop:d4eq} apply pointwise with uniform smallness: the bodies are analytic, globally strictly convex, origin-symmetric bodies of revolution, strictly contain $B^4$ and satisfy every distance-one section equality with volume $\lambda_\infty(\theta)$. Fixed complex widths give uniform convergence to the ball of radius two in every fixed $C^j$ norm of the radial function. Together with \eqref{eq:Vmeasure} this proves \cref{thm:d4meas}, with $V_\eps^\pm$ the set $V_\eps$ for the corresponding sign. \qed

\part{Dimensions at least five}\label{part:d5}

Throughout \cref{part:d5}, $d\ge5$ is a fixed integer, and
\begin{gather*}
\lambda=\tfrac{d-2}2,\qquad \beta=\lambda-1=\tfrac{d-4}2>0,\qquad R_0=\sqrt d,\\
\theta_0=\arccos\tfrac1{\sqrt d},\qquad \omega_0=2\theta_0,\qquad I_h=[\omega_0+h,\omega_0+2h].
\end{gather*}
In this part $\lambda$ is the Gegenbauer parameter; section volumes are denoted by $A$ or $V$. For a frequency $\omega=2\theta$ we write $R=\sec\theta$ for the radius of the round reference ball; the physical mean radius $R_*$ of a body is distinguished from $R$ in \cref{sec:center}. All constants may depend on $d$, and no uniformity in $d$ is claimed. We use the zonal polynomials $Y_n=C_n^\lambda/C_n^\lambda(1)$ in the variable $x=\cos t$, where $C^\lambda_n$ are the Gegenbauer polynomials, normalized by the generating function $(1-2xs+s^2)^{-\lambda}=\sum_nC^\lambda_n(x)s^n$ \cite[(18.12.4)]{DLMF}. Only even degrees $n$ occur, because all bodies are origin symmetric. We write $\Pi_n$ for the orthogonal projection onto the multiples of $Y_n$ in $L^2(S^{d-1})$ (physical projections), and $\Pi_0$ for the spherical mean.

\section{Section operators}\label{sec:sections}

\subsection{The section volume as an analytic functional}
Let $r$ be the radial function of a body of revolution close to the ball of radius $R>1$, written as a function of the axial cosine $s=\langle u,e_1\rangle$ of $u\in S^{d-1}$. The hyperplane with unit normal $\nu_\phi=(\cos\phi,\sin\phi,0,\dots,0)$ at distance one is parametrized by $x=\nu_\phi+q\,\varpi$ with $\varpi$ a unit vector orthogonal to $\nu_\phi$. Write $z\in[-1,1]$ for the cosine of the angle between $\varpi$ and the normalized projection of $-e_1$ to $\nu_\phi^\perp$. The point $x$ has $\abs x=\sqrt{1+q^2}$ and axial cosine $s=(\cos\phi-q\sin\phi\,z)/\sqrt{1+q^2}$. The boundary of the section in direction $(\phi,z)$ is the positive root $q=q(\phi,z)$ of
\begin{equation}\label{eq:q}
\sqrt{1+q^2}=r(s),\qquad s=\frac{\cos\phi-q\sin\phi\,z}{\sqrt{1+q^2}} .
\end{equation}

\begin{lemma}[Fixed-ray formulas]\label{lem:fixedray}
Fix $d\ge4$, a compact interval $\mathcal J\subset(1,\infty)$ and a Bernstein ellipse $\mathfrak E\supset[-1,1]$. There are $\delta_2>0$ and a complex neighborhood $U$ of $\R\times[-1,1]$, both depending only on $d$, $\mathcal J$ and $\mathfrak E$, with the following property. Let $R\in\mathcal J$ and let $r$ be holomorphic on $\mathfrak E$, real on $[-1,1]$, with $\sup_{\mathfrak E}\abs{r-R}\le\delta_2$. Then \eqref{eq:q} has, for $(\phi,z)\in U$, a unique solution $q(\phi,z)$ with $\abs{q-\sqrt{R^2-1}}<\frac12\sqrt{R^2-1}$; it is holomorphic in $(\phi,z)$, depends holomorphically on $r$, and for real $(\phi,z)$ it is the length of the section ray. The section volume $A_r(\phi)=\vol_{d-1}(K_r\cap\{x\cdot\nu_\phi=1\})$ and the radial derivative $T_rf(\phi)=\frac d{d\eps}A_{r+\eps f}(\phi)\bigr|_{\eps=0}$ are
\begin{align}
A_r(\phi)&=\frac{\abs{S^{d-3}}}{d-1}\int_{-1}^1q^{d-1}(1-z^2)^\beta\,dz,\notag\\
T_rf(\phi)&=\abs{S^{d-3}}\int_{-1}^1\frac{q^{d-2}f(s)}{q/\sqrt{1+q^2}-r'(s)\,\partial_qs}(1-z^2)^\beta\,dz,\label{eq:Tr}\\
\partial_qs&=\frac{-\sin\phi\,z-q\cos\phi}{(1+q^2)^{3/2}}.\notag
\end{align}
Both are holomorphic in $\phi$ on the projection of $U$, including $\phi\in\pi\Z$. For $0<\phi<\pi$ the equivalent coarea formula in the axial variable $s$ is
\begin{equation}\label{eq:coarea}
\begin{gathered}
T_rf(\phi)=\frac{\abs{S^{d-3}}}{(\sin\phi)^{2\lambda-1}}\int r(s)^2\,\Disc(\phi,s)^\beta f(s)\,ds,\\
\Disc=r(s)^2\sin^2\phi\,(1-s^2)-(1-r(s)s\cos\phi)^2 ,
\end{gathered}
\end{equation}
the integral being over the interval where $\Disc\ge0$.
\end{lemma}

\begin{proof}
For real $(\phi,z)$, Cauchy--Schwarz gives $\abs{\cos\phi-q\sin\phi z}\le\sqrt{1+q^2}$, so $s\in[-1,1]$; for $(\phi,z)$ in a small complex neighborhood of the real set, $s$ stays in $\mathfrak E$. At the ball, $q_0=\sqrt{R^2-1}$ solves \eqref{eq:q} for all $(\phi,z)$, and the derivative of $\sqrt{1+q^2}-r(s)$ in $q$ is $q_0/R>0$. The holomorphic implicit function theorem, with the uniform bounds on $\mathcal J$ and $\mathfrak E$, gives $\delta_2$, $U$ and the solution; for real data, convexity of $K_r$ (\cref{lem:convex}) identifies the root with the section ray. Polar coordinates in the $(d-1)$-dimensional section give $q^{d-1}/(d-1)$ for each direction; integrating over the latitude of the unit sphere $S^{d-2}$ of directions with respect to the axial direction gives the factor $\abs{S^{d-3}}(1-z^2)^\beta dz$ with $\beta=(d-4)/2\ge0$. Implicit differentiation of \eqref{eq:q} with $r\to r+\eps f$ gives the denominator in \eqref{eq:Tr}, which equals $q_0/R$ at the ball and stays nonzero. Since $(1-z^2)^\beta dz$ is a finite measure and the integrands are holomorphic, $A_r$ and $T_rf$ are holomorphic on the projection of $U$. For \eqref{eq:coarea} we write the variation of the volume as an integral over the boundary of the section: the normal speed of the boundary point of the section with axial cosine $s$ under $r\to r+\eps f$ is $f(s)$ times the ratio of the area elements, and integrating out the transverse latitude, with Jacobian $1/(r\sin\phi\sqrt{1-s^2})$, leaves exactly $r^2\Disc^\beta$.
\end{proof}

\begin{lemma}[Slice formula]\label{lem:slice}
In the situation of \cref{lem:fixedray}, write the meridian curve as $t\mapsto r(t)(\cos t,\sin t)$, put $X(t)=r(t)\cos t$ and $\Disc(\phi,t)=r(t)^2\sin^2\phi\sin^2t-(1-r(t)\cos t\cos\phi)^2$. For $0<\phi<\pi$,
\begin{equation}\label{eq:slice}
A_r(\phi)=\kappa_{d-2}\,(\sin\phi)^{-2\lambda-1}\int\Disc(\phi,t)^{\lambda}\,\abs{X'(t)}\,dt,
\end{equation}
the integral being over the $t$-interval between the two points where the line $\{x_1\cos\phi+x_2\sin\phi=1\}$ meets the meridian curve.
\end{lemma}

\begin{proof}
Let $x=(x_1,x_2,y)\in\R\times\R\times\R^{d-2}$. On the hyperplane $x_1\cos\phi+x_2\sin\phi=1$ we have $x_2=(1-x_1\cos\phi)/\sin\phi$, and the slice at fixed $x_1$ of $K_r$ is $\{y:\ x_2^2+\abs y^2\le\varrho(x_1)^2\}$, where $\varrho(x_1)=r(t)\sin t$ is the radius of $K_r$ at the axial coordinate $x_1=r(t)\cos t$. Hence the slice is a $(d-2)$-ball of radius $\sqrt{\Disc}/\sin\phi$. The arclength in the hyperplane along the $x_1$-direction is $\abs{dx_1}/\sin\phi$, and near the ball $t\mapsto X(t)$ is monotone. Integrating the slice volumes gives \eqref{eq:slice}, with $(d-2)/2=\lambda$.
\end{proof}

By \cref{lem:fixedray}, $A_r$ is holomorphic through $\phi\in\pi\Z$, so the apparent singularity of \eqref{eq:slice} at $\sin\phi=0$ is removable; quantitatively, if the numerator of \eqref{eq:slice} is bounded on a strip, the quotient is bounded on a narrower strip with a loss $C_k\sigma^{-k}$ for a width loss $\sigma$ and $k=2\lambda+1$, by the maximum principle on circles of radius $\sigma$ around the zeros of $\sin$.

\subsection{The round operator}
Let $L_R=T_{R}$ be the derivative at the ball of radius $R>1$.

\begin{lemma}[Round multipliers]\label{lem:round}
If $f=\sum_nf_nY_n$ is an even zonal function expanded in zonal polynomials, then
\begin{equation}\label{eq:Ln}
\begin{gathered}
L_Rf(\phi)=\sum_nL_n(R)f_nY_n(\cos\phi),\\
L_n(R)=\abs{S^{d-2}}\,R(R^2-1)^{\frac{d-3}2}\,Y_n(1/R).
\end{gathered}
\end{equation}
\end{lemma}

\begin{proof}
The boundary of $B_R\cap H_\phi$ is the $(d-2)$-sphere of radius $\varrho_0=\sqrt{R^2-1}$ about the foot point $\nu_\phi$. A radial change $\eps f(u)$ changes $q$ by $R\eps f/\varrho_0$ because $\abs x^2=1+q^2$. Hence $L_Rf(\phi)=\abs{S^{d-2}}\varrho_0^{d-2}(R/\varrho_0)\,\operatorname{avg}\{f(u):u\in S^{d-1},\langle u,\nu_\phi\rangle=1/R\}$. By the Funk--Hecke formula the average of a zonal harmonic $Y_n(\langle u,e_1\rangle)$ over this latitude sphere equals $Y_n(1/R)Y_n(\langle\nu_\phi,e_1\rangle)$.
\end{proof}

\begin{lemma}[Radial Taylor formula]\label{lem:taylor}
For $r=R+f$ with $f$ analytic and small and every $N\ge1$,
\begin{equation}\label{eq:taylor}
A_{R+f}=A_R+\sum_{j=1}^N\frac1{j!}\,\partial_R^{\,j-1}L_R\bigl(f^j\bigr)+O\bigl(\norm f^{N+1}\bigr),
\end{equation}
where $\partial_R$ acts only on the $R$-dependence of the operator $L_R$ and $f^j$ is held fixed; the remainder is bounded in every fixed analytic norm.
\end{lemma}

\begin{proof}
The set $K_r\cap H_\phi$ is the central projection from the origin of $\{u\in S^{d-1}:\langle u,\nu_\phi\rangle\,r(u)\ge1\}$ to $H_\phi$, so $A_r(\phi)=\int_{S^{d-1}}\Phi_\phi(r(u),u)\,du$ with $\Phi_\phi(\varrho,u)=\mathbf 1_{\varrho\langle u,\nu_\phi\rangle\ge1}\langle u,\nu_\phi\rangle^{-d}$, a distribution in $\varrho$. The $j$-th Taylor term of $\varrho\mapsto\Phi(\varrho+f(u),u)$ at $\varrho=R$ is $f(u)^j\partial_\varrho^j\Phi/j!$, and $\int f^j\partial_\varrho^j\Phi\,du=\partial_R^{j-1}\bigl[\int f^j\,\partial_\varrho\Phi\,du\bigr]$ with $f^j$ independent of $R$; the bracket is $L_R(f^j)$. The analytic dependence of $A_r$ on $r$ in \cref{lem:fixedray} gives the remainder bound.
\end{proof}

\section{Resonances, arithmetic and finite initializers}\label{sec:resonances}

\subsection{Low-degree identities}
Write $c_{nm}$ for the coefficients in $Y_2Y_n=\sum_mc_{nm}Y_m$. Put $L_n=L_n(R_0)$ and $L_n'=\partial_RL_n(R_0)$.

\begin{lemma}[Exact low coefficients]\label{lem:low}
\begin{enumerate}[label=\textup{(\alph*)}]
\item The first zonal polynomials are
\[
Y_2(x)=\frac{dx^2-1}{d-1},\qquad Y_4(x)=\frac{(d+2)(d+4)x^4-6(d+2)x^2+3}{(d-1)(d+1)},
\]
and
\begin{equation}\label{eq:Y2sq}
Y_2^2=\frac2{(d-1)(d+2)}+\frac{4(d-2)}{(d-1)(d+4)}Y_2+\frac{d^2(d+1)}{(d-1)(d+2)(d+4)}Y_4 .
\end{equation}
\item At $R_0=\sqrt d$ one has $Y_2(1/R_0)=0$, $L_0=\abs{S^{d-2}}\sqrt d\,(d-1)^{(d-3)/2}\ne0$, $L_2=0$ and $L_2'=-2\abs{S^{d-2}}(d-1)^{(d-5)/2}\ne0$.
\item The values at $1/\sqrt d$ of the next two zonal polynomials are
\[
Y_4(1/\sqrt d)=-\frac{2(d+4)}{d^2(d+1)}\ne0,\qquad Y_6(1/\sqrt d)=\frac{16(d-2)(d+6)}{d^3(d+1)(d+3)}\ne0.
\]
\item The diagonal coefficient of multiplication by $Y_2$ is
\begin{equation}\label{eq:cnn}
c_{nn}=\frac{\lambda\,n(n+2\lambda)}{(2\lambda+1)\bigl((n+\lambda)^2-1\bigr)}=\frac\lambda{2\lambda+1}-\frac{\lambda(\lambda^2-1)}{(2\lambda+1)\bigl((n+\lambda)^2-1\bigr)},
\end{equation}
strictly increasing in $n\ge0$ since $\lambda>1$.
\end{enumerate}
\end{lemma}

\begin{proof}
(a) is a direct computation with the explicit Gegenbauer polynomials. (b) follows from \eqref{eq:Ln}: $Y_2(1/R)=(d/R^2-1)/(d-1)$ has derivative $-2/(\sqrt d(d-1))$ at $R_0$. (c): by the generating function at $x=1/\sqrt d$ put $g_n=d^{n/2}C_n^\lambda(1/\sqrt d)$, so $\sum g_ns^n=(1-2s+ds^2)^{-\lambda}$, and read off $g_4,g_6$. (d): apply the three-term recurrence $(n+1)C_{n+1}=2(n+\lambda)xC_n-(n+2\lambda-1)C_{n-1}$ twice; the coefficient of $C_n$ in $x^2C_n$ is $\frac{(n+1)(n+2\lambda)}{4(n+\lambda)(n+\lambda+1)}+\frac{n(n+2\lambda-1)}{4(n+\lambda)(n+\lambda-1)}$, and substitution into $Y_2$ gives \eqref{eq:cnn}. The derivative in $n$ is $2\lambda(n+\lambda)(\lambda^2-1)/[(2\lambda+1)((n+\lambda)^2-1)^2]>0$.
\end{proof}

\begin{definition}\label{def:r1}
Put $r_1=-\tfrac12c_{22}=-\dfrac{2(d-2)}{(d-1)(d+4)}$, $\ w_1=r_1+Y_2$, and
\[
a_{4,d}=-\frac{d^{3/2}(d+1)(d^3-2d^2-9d+12)}{2(d-1)^2(d+2)(d+4)^2},\qquad h_2=a_{4,d}\,Y_4.
\]
\end{definition}
The value $r_1$ is the branch direction of the degree-two equation: by \eqref{eq:Y2sq} the quadratic Taylor term of \eqref{eq:taylor} along $w_1$ has degree-two part $\frac12L_2'(2r_1+c_{22})$. Similarly the degree-four equation gives $a_{4,d}=-c_{24}L_4'/(2L_4)$, with $c_{24}$ the $Y_4$ coefficient in \eqref{eq:Y2sq}, and logarithmic differentiation of $L_4$ at $R_0$ gives $L_4'/L_4=(d^3-2d^2-9d+12)/[\sqrt d(d-1)(d+4)]$. For instance $a_{4,5}=-\tfrac{5\sqrt5}{72}$ and $a_{4,6}=-\tfrac{1071\sqrt6}{10000}$.

\subsection{Resonant degrees}
A degree $n\ge4$ is \emph{resonant} if $C_n^\lambda(1/\sqrt d)=0$, equivalently $L_n(R_0)=0$.

\begin{lemma}[Structure of the round kernel]\label{lem:kernel}
\begin{enumerate}[label=\textup{(\alph*)}]
\item Every zero of $C_n^\lambda$ in $(-1,1)$ is simple; hence $L_n'\ne0$ at every resonant degree.
\item Two resonant even degrees do not differ by $2$; the degrees $4$ and $6$ are not resonant.
\item For every resonant $n$ the detuning $\delta_n=L_n'(r_1+c_{nn})$ is nonzero, since $c_{nn}>c_{22}$ gives $r_1+c_{nn}>c_{22}/2>0$. Distinct resonances are not coupled to first order in the direction $w_1$, nor are they coupled to degree two.
\end{enumerate}
\end{lemma}

\begin{proof}
(a) An interior zero with zero derivative of a solution of $(1-x^2)y''-(2\lambda+1)xy'+n(n+2\lambda)y=0$ forces $y\equiv0$ by uniqueness for this regular second-order equation. (b) If $C_n(x_0)=C_{n+2}(x_0)=0$ with $x_0=1/\sqrt d>0$ then the recurrence at the intervening degree gives $C_{n+1}(x_0)=0$, and repeating backward gives $C_0(x_0)=0$, a contradiction. Degrees $4$ and $6$ are nonresonant by \cref{lem:low}(c). (c) By \eqref{eq:Y2sq} and the recurrence, multiplication by $Y_2$ changes the degree by at most two, so by (b) it does not couple two distinct resonant degrees, and it does not couple degree two to a resonant degree since resonant degrees are at least eight.
\end{proof}

For odd $d$ the full zero set $\{n:C_n^\lambda(1/\sqrt d)=0\}$ is not classified, and it is not asserted to be finite. For even $d$ it is finite, by \cref{prop:evenres} below.

\subsection{The Darboux representation in even dimensions}
For even $d$ we have $\lambda\in\N$, and we put $k=\lambda-1\ge1$ and $k_\lambda=2^kk!$. Let $\mathsf D_j=j\cot t-\partial_t$.

\begin{lemma}[Darboux representation]\label{lem:darboux}
Let $d=2\lambda+2\ge6$ be even and $m\ge\lambda$ an integer.
\begin{enumerate}[label=\textup{(\alph*)}]
\item $C^\lambda_{m-\lambda}(\cos t)=\dfrac{\mathsf D_k\cdots\mathsf D_1\sin(mt)}{k_\lambda\sin^\lambda t}$.
\item $\mathsf D_k\cdots\mathsf D_1\sin(mt)=a_k(m,\cot t)\sin(mt)+b_k(m,\cot t)\cos(mt)$, where $a_0=1$, $b_0=0$ and, with $\mathsf q=\cot t$,
\begin{equation}\label{eq:abrec}
a_j=j\mathsf qa_{j-1}+(1+\mathsf q^2)\partial_{\mathsf q}a_{j-1}+mb_{j-1},\quad b_j=j\mathsf qb_{j-1}+(1+\mathsf q^2)\partial_{\mathsf q}b_{j-1}-ma_{j-1}.
\end{equation}
The polynomial $a_k$ is even in $m$ and has parity $k$ in $\mathsf q$; $b_k$ is odd in $m$ and has parity $k+1$ in $\mathsf q$.
\item With $S_k=k(k+1)/2$: if $k=2\ell-1$, then $b_k=(-1)^\ell m^k+O(m^{k-2})$ and $a_k=(-1)^{\ell-1}S_k\mathsf q\,m^{k-1}+O(m^{k-3})$; if $k=2\ell$, then $a_k=(-1)^\ell m^k+O(m^{k-2})$ and $b_k=(-1)^\ell S_k\mathsf q\,m^{k-1}+O(m^{k-3})$, with coefficients polynomial in $\mathsf q$.
\item For $t\in(0,\pi)$ the numbers $a_k(m,\cot t)$ and $b_k(m,\cot t)$ do not vanish simultaneously.
\end{enumerate}
\end{lemma}

\begin{proof}
(a) The differentiation formula \cite[(18.9.19)]{DLMF} reads $\frac d{dx}C^\mu_N=2\mu C^{\mu+1}_{N-1}$. Hence $C^{\lambda}_{m-\lambda}=(k_\lambda)^{-1}(d/dx)^kC^1_{m-1}$, and $C^1_{m-1}(\cos t)=\sin(mt)/\sin t$. With $d/dx=-(\sin t)^{-1}\partial_t$, induction on $j$ gives $[-(\sin t)^{-1}\partial_t]^j[\sin(mt)/\sin t]=\mathsf D_j\cdots\mathsf D_1\sin(mt)/\sin^{j+1}t$, since $-(\sin t)^{-1}\partial_t[g\sin^{-(j+1)}t]=\mathsf D_{j+1}g/\sin^{j+2}t$. (b) Apply $\mathsf D_j$ to $a\sin(mt)+b\cos(mt)$ with $a,b$ functions of $\mathsf q$ and use $\partial_t\mathsf q=-(1+\mathsf q^2)$. The parities follow by induction from \eqref{eq:abrec}. (c) Induction on $k$ using \eqref{eq:abrec}: the top term of $b_{2\ell}$ or $a_{2\ell+1}$ receives $(k+1)\mathsf q$ times the top term of the same function and $\pm m$ times the top term of the other, and $S_k+k+1=S_{k+1}$. (d) Put $\mathsf L_j=-\partial_t^2+j(j-1)\csc^2t$. One checks $\mathsf D_j^\dagger\mathsf D_j=\mathsf L_j-j^2$ and $\mathsf D_j\mathsf D_j^\dagger=\mathsf L_{j+1}-j^2$ with $\mathsf D_j^\dagger=j\cot t+\partial_t$, hence $\mathsf D_j\mathsf L_j=\mathsf L_{j+1}\mathsf D_j$. Thus $\mathsf D_j\cdots\mathsf D_1$ maps the solutions $\sin(mt),\cos(mt)$ of $\mathsf L_1y=m^2y$ to solutions of $\mathsf L_{k+1}y=m^2y$. On solutions of $\mathsf L_jy=m^2y$ the map $(y,y')\mapsto(\mathsf D_jy,(\mathsf D_jy)')$ has the matrix $\begin{psmallmatrix}j\cot t&-1\\ m^2-j^2\csc^2t&j\cot t\end{psmallmatrix}$ of determinant $m^2-j^2$, so the Wronskian of the two transformed solutions is $-m\prod_{j=1}^k(m^2-j^2)\ne0$ for $m\ge k+1$. The transformed cosine is $a_k\cos(mt)-b_k\sin(mt)$. If $a_k$ and $b_k$ vanished at $t$, both transformed solutions would vanish there, contradicting the nonvanishing Wronskian.
\end{proof}

\begin{proposition}[Finite resonance in even dimensions]\label{prop:evenres}
Let $d=2\lambda+2\ge6$ be even. Then $C_n^\lambda(1/\sqrt d)$ vanishes for only finitely many $n$, and the set of zero degrees is effectively computable.
\end{proposition}

\begin{proof}
Put $\sigma=\sqrt{d-1}$, so $\cot\theta_0=1/\sigma$. By \cref{lem:darboux}(b) the parities of $a_k,b_k$ give $\sigma^ka_k(m,1/\sigma)=\tilde a(m)$ and $\sigma^kb_k(m,1/\sigma)=\sigma\tilde b(m)$ with $\tilde a,\tilde b\in\Z[m]$. Hence, for $m=n+\lambda$, the zero equation is $\tilde a(m)\sin m\theta_0+\sigma\tilde b(m)\cos m\theta_0=0$, and by \cref{lem:darboux}(d), $N_d(m):=\tilde a(m)^2+(d-1)\tilde b(m)^2>0$ for $m\ge\lambda$. With $s=i\sigma$ we have $e^{i\theta_0}=(1+s)/\sqrt d$, and the zero equation becomes $(\tilde a+\tilde bs)(1+s)^m=(\tilde a-\tilde bs)(1-s)^m$. If an odd prime $p$ divides $d$, Hensel lifting gives $r_p$ modulo $p^m$ with $r_p^2\equiv1-d$ and $r_p\equiv-1$ modulo $p$ (the derivative $2r_p$ is a unit); applying the ring homomorphism $\Z[s]\to\Z/p^m\Z$, $s\mapsto r_p$, the left side vanishes and $(1-r_p)^m$ on the right is a unit, hence $p^m$ divides $\tilde a-\tilde br_p$ and therefore $N_d(m)$. The remaining even dimensions are $d=2^j$, $j\ge3$: with $u=(1+s)/2$ one has $u^2-u+d/4=0$ and $u\bar u=d/4\ge2$; expanding $(u+\bar u)^{2m-1}=1$ shows that $u^m$ and $\bar u^m$ generate the unit ideal of $\Z[u]$, so $u^m$ divides $\tilde a-\tilde bs$ in $\Z[u]$, and taking norms $(d/4)^m$ divides $N_d(m)$. Since $N_d$ is a fixed positive polynomial, these exponential divisibilities fail beyond an effectively computable bound.
\end{proof}

\subsection{Circle arithmetic}\label{sec:circlearith}
Let $z_0=e^{2i\theta_0}=(2-d+2i\sqrt{d-1})/d$. Then $z_0+z_0^{-1}=4/d-2$. For $d\ge5$ this is rational and not an integer; since the trace of a root of unity is an algebraic integer, $z_0$ is not a root of unity and $\theta_0/\pi$ is irrational. The number $z_0$ is a root of $dz^2+(2d-4)z+d$, whose discriminant $16(1-d)$ is negative; so this polynomial is irreducible over $\Q$, and dividing it by the greatest common divisor $\gcd(d,4)$ of its coefficients gives the minimal polynomial of $z_0$ with relatively prime integer coefficients, all bounded in absolute value by $2d$.

\begin{lemma}[Logarithmic forms {\cite[p.~20]{BW}}]\label{lem:BW}
Let $\alpha_1,\dots,\alpha_n$ be algebraic numbers different from $0$ and $1$, let $D$ be the degree over $\Q$ of the field $\Q(\alpha_1,\dots,\alpha_n)$, and let $\Log$ denote the principal branch of the logarithm. Let $b_1,\dots,b_n$ be integers and $B\ge\max_j\abs{b_j}$, $B\ge e$. For each $j$ let $A_j\ge e$ be an upper bound for the absolute values of the coefficients of the minimal polynomial of $\alpha_j$ over $\Z$ with relatively prime coefficients. If $\Xi=b_1\Log\alpha_1+\dots+b_n\Log\alpha_n\ne0$, then
\[
\log\abs\Xi>-(16nD)^{2(n+2)}\log A_1\cdots\log A_n\,\log B .
\]
\end{lemma}

\begin{lemma}[Rotation arithmetic at $\omega_0$]\label{lem:circle}
There are $c_d>0$ and $\tau_0=64^8\log(2d)$ such that
\begin{equation}\label{eq:sinbw}
\abs{\sin k\theta_0}\ge c_d\,k^{-\tau_0}\qquad(k\ge1).
\end{equation}
Fix $\tau>\tau_0+1$ and $0<\gamma\le c_d/2$, and set
\[
\Dh=\{\omega\in I_h:\ \abs{\sin(k\omega/2)}\ge\gamma k^{-\tau}\ \text{for all integers }k\ge1\}.
\]
Then $\Dh$ is compact and $\Leb(I_h\setminus\Dh)=o_d(h)$.
\end{lemma}

\begin{proof}
Apply \cref{lem:BW} with $n=2$, $D=2$, $\alpha_1=z_0$, $\alpha_2=-1$, $A_1=2d$, $A_2=e$, $B=\max(e,2k)$ to $\Xi=k\Log z_0-2j\Log(-1)=2i(k\theta_0-j\pi)$, with $j$ the integer nearest to $k\theta_0/\pi$; here $\Log z_0=2i\theta_0$ because $0<2\theta_0<\pi$, and $\Log(-1)=i\pi$. The hypotheses hold: $\Xi\ne0$ as $\theta_0/\pi$ is irrational, $\abs{2j}\le k+1\le2k$, both numbers are different from $0$ and $1$, $\Q(z_0,-1)$ has degree two, and the coefficient bounds were noted above. Hence $\abs\Xi\ge(\max(e,2k))^{-64^8\log(2d)}$, and $\abs{\sin k\theta_0}\ge2\dist(k\theta_0,\pi\Z)/\pi=\abs\Xi/\pi$.

For $\omega\in I_h$ and $k\le K:=\lfloor(c_d/(2h))^{1/(\tau_0+1)}\rfloor$ we have $\abs{\sin(k\omega/2)-\sin k\theta_0}\le kh\le\frac{c_d}2k^{-\tau_0}$, hence $\abs{\sin(k\omega/2)}\ge\frac{c_d}2k^{-\tau_0}\ge\gamma k^{-\tau}$. For $k>K$ the set $\{\omega\in I_h:\abs{\sin(k\omega/2)}<\gamma k^{-\tau}\}$ consists of at most $kh/(2\pi)+2$ intervals of length at most $C\gamma k^{-\tau-1}$. Summing over $k>K$ gives at most $C\gamma[hK^{1-\tau}+K^{-\tau}]=o_d(h)$, because $\tau>1$ and $\tau>\tau_0+1$. Compactness is clear.
\end{proof}

\begin{lemma}[Even base arithmetic]\label{lem:evenlow}
Let $d\ge6$ be even. There are $c,C>0$ such that $\abs{C^\lambda_n(1/\sqrt d)}\ge c\exp[-C\log^2(n+2)]$ for every $n\ge0$ with $C^\lambda_n(1/\sqrt d)\ne0$.
\end{lemma}

\begin{proof}
Let $m=n+\lambda$ and use the notation of the proof of \cref{prop:evenres}. Writing $a_k=a_k(m,1/\sigma)$, $b_k=b_k(m,1/\sigma)$ and $x=m\theta_0$,
\[
\sigma^k(a_k\sin x+b_k\cos x)=\frac{(\tilde a+\tilde bs)e^{-ix}}{2i}\bigl(e^{2ix}-\xi_m\bigr),\qquad\xi_m=\frac{\tilde a-\tilde bs}{\tilde a+\tilde bs},
\]
and $\abs{\tilde a+\tilde bs}=N_d(m)^{1/2}\ge1$. Since $s$ is purely imaginary, $\abs{\xi_m}=1$. By \cref{lem:darboux}(a) it suffices to bound $\abs{z_0^m-\xi_m}$ from below when it is nonzero. If $\tilde b(m)=0$ then $\xi_m=1$, and if $\tilde a(m)=0$ then $\xi_m=-1$; in these cases $\abs{z_0^m\mp1}$ is bounded below by the two-logarithm argument of \cref{lem:circle}, which gives a polynomial bound. Otherwise $\xi_m\ne\pm1$, so $\xi_m\notin\Q$, and $\xi_m$ is a root of $N_d(m)z^2-2(\tilde a^2-(d-1)\tilde b^2)z+N_d(m)$; its minimal polynomial with relatively prime integer coefficients divides this quadratic, so its coefficients are bounded by $2N_d(m)\le C_dm^{2k}$. Choose $\ell\in\Z$ with $\abs\ell\le m+1$ such that $\Xi=m\Log z_0-\Log\xi_m-2\ell\Log(-1)$ lies in $i(-\pi,\pi]$. Then $\abs{z_0^m-\xi_m}=\abs{e^\Xi-1}\ge\frac2\pi\abs\Xi$, and $\Xi\ne0$ whenever the Gegenbauer value is nonzero. \Cref{lem:BW} with $n=3$, $D=2$ (all three numbers lie in $\Q(z_0)$), $A_1=2d$, $A_2=\max(e,C_dm^{2k})$, $A_3=e$ and $B=2m+2$ gives $\abs\Xi\ge\exp[-C\log^2(m+2)]$. The factors $\sigma^{-k}$, $\frac12N_d(m)^{1/2}\ge\frac12$ and $(k_\lambda\sin^\lambda\theta_0)^{-1}$ only change the constants.
\end{proof}

\subsection{Integrality of the base values for odd \texorpdfstring{$d$}{d}}
Let $d$ be odd, $m_*=(d-3)/2$ so that $\lambda=m_*+\frac12$, and define rational numbers $b_n$ by $(1-2s+ds^2)^{-\lambda}=\sum_nb_ns^n$, so that $b_n=d^{n/2}C_n^\lambda(1/\sqrt d)$. The recurrence is $b_0=1$, $b_1=2\lambda$, $(n+1)b_{n+1}=2(n+\lambda)b_n-d(n+2\lambda-1)b_{n-1}$.

\begin{lemma}\label{lem:integral}
With $c=(2m_*)!/m_*!$, $c\,4^nb_n\in\Z$ for all $n$. Consequently $b_n\ne0$ implies $\abs{b_n}\ge1/(c\,4^n)$, and the normalized round value $(1+n)^{\lambda+1}C_n^\lambda(\cos\theta_0)/C_n^\lambda(1)$ has absolute value at least $e^{-C_dn}$ whenever it is nonzero.
\end{lemma}

\begin{proof}
For every integer $j\ge0$, $c\,4^j(m_*+\frac12)_j/j!=\binom{2j+2m_*}{j+m_*}(j+m_*)!/j!\in\Z$. Expand $(1-(2s-ds^2))^{-\lambda}=\sum_j\frac{(\lambda)_j}{j!}(2s-ds^2)^j$; the coefficient of $s^n$ is a finite sum over $j\le n$, and multiplying each term by $c\,4^n$ clears its denominator. Finally $C_n^\lambda(1)=\binom{n+d-3}{n}$ is a positive polynomial in $n$ of degree $d-3$.
\end{proof}

This lemma does not determine which $b_n$ vanish, nor whether their zero set is finite.

\subsection{Finite formal solutions}
For the next proposition we write the physical radial function as a function of the axial angle $t$, and we use the scalar quantities
\begin{gather*}
R_*=\Pi_0r,\qquad \eps=\text{the coefficient of $Y_2$ in $r$},\\
a_2=\mean_tY_2(\cos t)=\frac{d-2}{2(d-1)},
\end{gather*}
where $\mean_t$ is the angular mean in $t$ (not the spherical mean).

\begin{proposition}[Finite initializers]\label{prop:init}
Fix $\rho_*\in(0,1]$. For every $N\ge2$ there is $h_N>0$ such that for $0<h<h_N$ there is a family $v^{(N)}_\omega$, $\omega\in I_h$, polynomial in $\omega-\omega_0$ with coefficients that are real odd $\pi$-periodic trigonometric polynomials, such that the body $K(v^{(N)}_\omega,\omega)$ of \cref{def:centered} with radial function $r=r_\omega$ has the following properties, uniformly in $\omega\in I_h$:
\begin{enumerate}[label=\textup{(\roman*)}]
\item $(A_r-\Pi_0A_r)/\eps=O(h^N)$ and $h\,\partial_\omega[(A_r-\Pi_0A_r)/\eps]=O(h^N)$ in every norm $\norm\cdot_{\rho_*,q}$ of the output angle;
\item $c_{\mathrm t}h\le\eps\le C_{\mathrm t}h$, $R_*-R_0-r_1\eps=O(h^2)$, and $r=R_*+\eps Y_2+\eps^2h_2+\eps u_*$ with $\Pi_0u_*=\Pi_2u_*=0$ and $u_*=O(h^2)$, together with the same bounds for $h\,\partial_\omega$;
\item the frequency slope is
\[
\omega_1=\frac{2(r_1+a_2)}{\sqrt{d(d-1)}}=\frac{\sqrt d\,(d-2)}{(d-1)^{3/2}(d+4)}>0,
\]
and the mean section volume $V(\omega)=\Pi_0A_r$ satisfies
\begin{equation}\label{eq:Vslope}
V'(\omega)=\frac{L_0\,r_1}{\omega_1}+O(h)=-2\abs{S^{d-2}}(d-1)^{\frac{d-2}2}+O(h).
\end{equation}
\end{enumerate}
\end{proposition}

\begin{proof}
\emph{Step 1: the radial expansion.} We seek $r_\eps=R_0+\eps w_1+\sum_{j\ge2}\eps^jv_j$ with $v_1:=w_1$ and the gauge $\Pi_2v_j=0$ for $j\ge2$, so that $\eps$ is the physical degree-two coefficient. Write
\[
v_j=c_j+\sum_{n\in\mathcal Z}\mathsf a_{j,n}Y_n+\sum_{n\in\mathcal N}\mathsf b_{j,n}Y_n,
\]
where $\mathcal Z$ is the set of resonant even degrees and $\mathcal N$ the set of nonresonant even degrees $n\ge4$; we put $c_1=r_1$ and $\mathsf a_{1,n}=\mathsf b_{1,n}=0$. By \cref{lem:taylor}, with $L=L_{R_0}$ and $L'=\partial_RL_R|_{R_0}$, the coefficient of $\eps^j$ in $A_{r_\eps}$ is
\[
Lv_j+L'(w_1v_{j-1})+\mathcal F_j\quad(j\ge3),\qquad Lv_2+\tfrac12L'(w_1^2)\quad(j=2),
\]
where $\mathcal F_j$ is a finite polynomial expression in $v_1,\dots,v_{j-2}$: in the $i$-th Taylor term the coefficient of $\eps^j$ is a sum of products $v_{l_1}\cdots v_{l_i}$ with $l_1+\dots+l_i=j$, and only the terms with $i=1$ or with $i=2$ and $\{l_1,l_2\}=\{1,j-1\}$ involve $v_j$ or $v_{j-1}$. We require that all components of degree $n\ge2$ of these coefficients vanish for $j\le N$; the degree-zero component is the volume and is free.

At order $j=2$ the degree-two component is $\frac12L_2'(2r_1+c_{22})$, which vanishes by the choice of $r_1$; the forcing $\frac12L'(w_1^2)$ has degrees $0,2,4$ only, so the degree-four equation $L_4\mathsf b_{2,4}+\frac12L_4'c_{24}=0$ gives $\mathsf b_{2,4}=a_{4,d}$, all other $\mathsf b_{2,n}$ vanish, and the resonant equations at order two read $0=0$ since resonant degrees are at least eight.

At order $j\ge3$ the equations are solved in the following order.
\begin{enumerate}[label=(\arabic*)]
\item \emph{Mean of $v_{j-1}$.} The degree-two component is $L_2'\bigl(c_{j-1}+c_{42}\mathsf b_{j-1,4}\bigr)+\Pi_2\mathcal F_j$, because $L_2=0$, $\Pi_2v_j=0$ and $Y_2Y_n$ has a $Y_2$ component only for $n\in\{0,2,4\}$. Since $L_2'\ne0$ this determines $c_{j-1}$.
\item \emph{Resonant part of $v_{j-1}$.} For $n\in\mathcal Z$ the degree-$n$ component is $\delta_n\mathsf a_{j-1,n}+L_n'(c_{n-2,n}\mathsf b_{j-1,n-2}+c_{n+2,n}\mathsf b_{j-1,n+2})+\Pi_n\mathcal F_j$, because $L_n=0$, the mean contributes only to degrees $0,2$, and no other resonant degree is within distance two (\cref{lem:kernel}). The first nonzero coefficient of $\mathsf a_{j-1,n}$ is the detuning $\delta_n\ne0$, so this determines $\mathsf a_{j-1,n}$.
\item \emph{Nonresonant part of $v_j$.} For $n\in\mathcal N$ the degree-$n$ component is $L_n\mathsf b_{j,n}+\Pi_n[L'(w_1v_{j-1})+\mathcal F_j]$, where $v_{j-1}$ is now completely known; since $L_n\ne0$ this determines $\mathsf b_{j,n}$.
\end{enumerate}
Inductively every $v_l$ with $l\le j-1$ is a finite combination of zonal polynomials, and $L_R$ is diagonal on them, so every forcing is a finite combination and only finitely many coefficients determined at stage $j$ are nonzero, even if $\mathcal Z$ is infinite. The procedure determines all coefficients uniquely except the mean $c_N$ and the resonant part $\mathsf a_{N,\cdot}$ of $v_N$, which enter first at order $N+1$; we set them equal to zero. Then $r^{(N)}_\eps=R_0+\sum_{j=1}^N\eps^jv_j$ satisfies $A_{r^{(N)}_\eps}-\Pi_0A_{r^{(N)}_\eps}=O(\eps^{N+1})$ in every analytic norm, by \cref{lem:taylor} with remainder of order $N+1$. Its physical degree-two coefficient is $\eps$, its mean is $R_0+r_1\eps+O(\eps^2)$, and $v_2=c_2+h_2$ since $\mathsf a_{2,n}=0$ for $n\in\mathcal Z$ (the forcing at order three in a degree $\ge8$ vanishes, its degrees being at most six); hence the decomposition in (ii) holds with $\eps u_*=\sum_{j\ge3}\eps^j(v_j-c_j)$.

\emph{Step 2: the centered representation.} We seek $p_\eps=u+\sum_{j=1}^N\eps^jp_j$, $p_j$ odd $\pi$-periodic trigonometric polynomials, and $\theta_\eps=\theta_0+\sum_{j=1}^N\theta_j\eps^j$ such that the radial function of \cref{def:centered} built from $(p_\eps,\theta_\eps)$ agrees with $r^{(N)}_\eps$ up to $O(\eps^{N+1})$, i.e.\ $\sec\alpha(y)-r^{(N)}_\eps(\chi(y))=O(\eps^{N+1})$. At $\eps=0$ we have $\chi=\id$ and $\alpha=\theta_0$. Expanding in $\eps$, the terms of order $j$ that contain $p_j$ or $\theta_j$ are $R_0\tan\theta_0\bigl[\theta_j+\frac12(p_j(y+\theta_0)-p_j(y-\theta_0))\bigr]$, all other terms being determined by $p_1,\dots,p_{j-1}$, $\theta_1,\dots,\theta_{j-1}$ and $r^{(N)}$. If $p_j=\sum_{\ell}p_{j,\ell}\sin(\ell u)$, then $\frac12(p_j(y+\theta_0)-p_j(y-\theta_0))=\sum_\ell p_{j,\ell}\sin(\ell\theta_0)\cos(\ell y)$. Hence the angular mean of the order-$j$ equation determines $\theta_j$, and its $\cos(\ell y)$ coefficient determines $p_{j,\ell}$ after division by $R_0\tan\theta_0\sin(\ell\theta_0)\ne0$ ($\ell\ge2$ even), which is nonzero because $\theta_0/\pi$ is irrational. Each forcing is a trigonometric polynomial, so each $p_j$ is one. At order one the forcing is $w_1(\cos y)$, so $\theta_1=\mean_y w_1/(R_0\tan\theta_0)=(r_1+a_2)/\sqrt{d(d-1)}$, i.e.\ $\omega_1=2\theta_1$. Since $p_\eps,\theta_\eps$ are polynomials in $\eps$ and the centered construction is analytic in $(p,\theta)$, the difference of the two radial functions is an analytic function of $\eps$ whose first $N$ Taylor coefficients vanish, hence $O(\eps^{N+1})$ on a fixed strip.

\emph{Step 3: the frequency as parameter.} Since $\theta_1=\omega_1/2>0$, the polynomial $\eps\mapsto2\theta_\eps-\omega_0$ has a truncated inverse $\mathfrak e(\omega)=\sum_{i=1}^Ne_i(\omega-\omega_0)^i$ with $e_1=1/\omega_1$ and $2\theta_{\mathfrak e(\omega)}=\omega+O((\omega-\omega_0)^{N+1})$. Let $v^{(N)}_\omega$ be the truncation at degree $N$ in $\omega-\omega_0$ of $p_{\mathfrak e(\omega)}-u$. Then $p=u+v^{(N)}_\omega$ at $\theta=\omega/2$ differs from $(p_{\mathfrak e(\omega)},\theta_{\mathfrak e(\omega)})$ by $O(h^{N+1})$, so its radial function differs from $r^{(N)}_{\mathfrak e(\omega)}$ by $O(h^{N+1})$, and $A_r-\Pi_0A_r=O(h^{N+1})$ by \cref{lem:fixedray}. Dividing by $\eps\asymp h$ gives (i) for the values. All quantities are polynomial, hence holomorphic, in complex $\omega$ in a disk of radius $3h$ about $\omega_0$; Cauchy's estimate on the disk of radius $h$ about each point of $I_h$ gives the bounds for $h\partial_\omega$. The physical degree-two coefficient of $r$ is $\mathfrak e(\omega)+O(h^{N+1})\asymp h$, which gives (ii) after replacing $\eps$ by the exact coefficient. Finally $\Pi_0A_{r^{(N)}_\eps}=A_{R_0}+\eps L_0r_1+O(\eps^2)$, since the degree-zero part of the first-order term is $L_0c_1$, and $d\eps/d\omega=1/\omega_1+O(h)$; with $L_0=\abs{S^{d-2}}\sqrt d(d-1)^{(d-3)/2}$ and $r_1/\omega_1=-2\sqrt{d-1}/\sqrt d$ this gives \eqref{eq:Vslope}. The only choices made are the gauge $\Pi_2v_j=0$, the terminal values $c_N=0$, $\mathsf a_{N,\cdot}=0$, and the two truncations at degree $N$; the threshold $h_N$ depends on $N$ through the size of the finitely many coefficients, while the width $\rho_*$ is fixed in advance.
\end{proof}

\section{Normalized linearizations}\label{sec:normalization}

The linearization of the section map at a profile near the ball is a nonlocal operator, and in the native coordinates of \cref{def:centered} it has small divisors coming from the rotation $\omega$. In this section we construct, for both parities of $d$, exact global normalizations which exhibit it as a diagonal operator plus a localized $O(h)$ perturbation.

\subsection{Matrix classes}
For $\rho>0$ and a matrix $B=(B_{mn})$ put
\[
N_\rho(B)=\max\Bigl\{\sup_m\sum_n\abs{B_{mn}}e^{\rho\abs{m-n}},\ \sup_n\sum_m\abs{B_{mn}}e^{\rho\abs{m-n}}\Bigr\}.
\]
$N_\rho$ is submultiplicative. Let $\ell^2_\rho=\{u:\sum e^{2\rho n}\abs{u_n}^2<\infty\}$. By Schur's test a matrix with $N_\rho(B)<\infty$ acts boundedly on $\ell^2_{r}$ for $0\le r\le\rho$, with norm at most $N_\rho(B)$. Let $\Lambda_0=\operatorname{diag}(1+n)$.

\begin{definition}[Normalized families]\label{def:matrixclass}
Fix $M\ge1$ and $0<\rho_\infty<\rho_{\max}$. A \emph{normalized family} on $I_h$ is a family of real matrices $K(\omega)=D(\omega)+A(\omega)$ acting on the modes $n\ge M$ of a fixed parity (step two), where $D=\operatorname{diag}(d_n(\omega))$, $\theta=\omega/2$, is of one of the following types.
\begin{enumerate}[label=\textup{(\Roman*)}]
\item[\textup{(O)}] $d_n=(1+n)^{\lambda+1}C_n^\lambda(\cos\theta)/C_n^\lambda(1)$ (odd $d$; $n$ is the degree of an orthonormal Gegenbauer polynomial).
\item[\textup{(B)}] $d_m=-m\cot(m\theta)+q_m$, where $(q_m)$ is a fixed real sequence, independent of $\omega$, with $q_m=q_\infty+O(m^{-2})$ (even $d$ with $\lambda$ even; $m$ labels the mode $\sin(m\,\cdot)$ of a primitive).
\item[\textup{(E)}] $d_m=-\cot(m\theta)/m+q_m$, where $(q_m)$ is a fixed real sequence, independent of $\omega$, with $q_m=-b_E+O(m^{-2})$ and $b_E>0$ (even $d$ with $\lambda$ odd).
\end{enumerate}
The perturbation $A$ is real and, for every $\rho\in[\rho_\infty,\rho_{\max}]$,
\begin{equation}\label{eq:Aclass}
\begin{gathered}
N_\rho(A(\omega))\le Ch,\qquad\partial_\omega A=A_0+\Lambda_0A_1,\\
N_\rho(A_0)\le C,\qquad N_\rho(A_1)\le Ch,
\end{gathered}
\end{equation}
where $\omega\mapsto A(\omega)$ is $C^1$ on $I_h$ as a map into the bounded operators $\ell^2_\rho\to\ell^2_{\rho-\sigma}$ for every $\sigma>0$. In types (B) and (E), $A_1=0$.
\end{definition}

In type (O) the frequency derivative is of the weaker form $A_0+\Lambda_0A_1$, which costs one power of the mode; it is the form provided by the odd normalization below.

\subsection{Odd dimensions}
Let $d\ge5$ be odd. Fix $\theta$ in a compact subinterval of $(0,\pi/2)$, $\omega=2\theta$ and $R=\sec\theta$. Let $p=u+v$ be as in \cref{def:centered}, with $\abs v_{\rho,1}$ small on a fixed strip, and let $G$ be the odd analytic function on a Bernstein ellipse with $\cos p(u)=G(\cos u)$; it exists because $\cos p(u)$ is even and $\pi$-antiperiodic, and $G(\pm1)=\pm1$. The quotients $\sin p(u)/\sin u$ and $\sin\chi(y)/\sin y$ are analytic, nonzero and close to $1$, including at the meridians $u,y\in\pi\Z$, and fractional powers of them are taken on the branches near $1$.

Let $\tilde T$ be the section derivative $T_r$ of \cref{lem:fixedray} acting on the input $\tilde f(y)=f(\cos\chi(y))$, with output at the normal angle $\phi=p(u)$. Put
\begin{equation}\label{eq:Lnorm}
\begin{gathered}
a_0=R^{2\lambda}(\sin\theta)^\beta,\qquad a_{\mathrm{in}}(y)=\hat r^{2\lambda}(\sin\alpha)^\beta\chi'\Bigl[\frac{\sin\chi}{\sin y}\Bigr]^\lambda,\\
a_{\mathrm{out}}(u)=\frac{[\sin p/\sin u]^\lambda}{(p')^\beta},\qquad
\mathcal L=a_0\,a_{\mathrm{out}}\,\tilde T\,a_{\mathrm{in}}^{-1}.
\end{gathered}
\end{equation}
These are regular invertible analytic multipliers, equal to $a_0$, $a_0$ and $1$ at the ball. Below $X=\cos u$, $Y=\cos y$, $\xi_\mp=\cos(y\mp\theta)$, and $G[a,b]=(G(a)-G(b))/(a-b)=\int_0^1G'(b+\tau(a-b))\,d\tau$ is the divided difference.

\begin{proposition}[Exact global normalization, odd $d$]\label{prop:oddnorm}
\begin{enumerate}[label=\textup{(\alph*)}]
\item At the output $\phi=p(u)$ and the axial cosine $s=\cos\chi(y)$ the discriminant of \eqref{eq:coarea} is $\Disc=\hat r^2[G(\xi_-)-G(X)][G(X)-G(\xi_+)]$.
\item Let $\Disc_0=R^2(1-X^2-Y^2)+2RXY-1=R^2(\xi_--X)(X-\xi_+)$ be the round discriminant, and
\[
\Psi=\frac{G[\xi_-,X]\,G[X,\xi_+]}{G[\xi_-,\xi_+]\,G'(X)} .
\]
Then, for $F$ a function of $Y$,
\begin{equation}\label{eq:Lkernel}
\mathcal LF(X)=\abs{S^{d-3}}R^2(1-X^2)^{-(\lambda-\frac12)}\int_{\Disc_0\ge0}\Disc_0^{\beta}\,\Psi^\beta\,F(Y)\,dY .
\end{equation}
For $\Psi\equiv1$ this is the round derivative $T_{d,R}$.
\item There is an analytic function $\mathfrak E=\mathfrak E_{d,\theta}(X,Y)$, depending analytically on $G$ and vanishing for $G=\id$, such that $\Psi^\beta-1=(X-\xi_-)(X-\xi_+)\,\mathfrak E$. Consequently
\begin{equation}\label{eq:dd2}
\mathcal LF=T_{d,R}F-\frac{\abs{S^{d-3}}}{(1-X^2)^{\lambda-1/2}}\int_{\Disc_0\ge0}\Disc_0^{\lambda}\,\mathfrak E(X,Y)F(Y)\,dY ,
\end{equation}
and for a separated amplitude $\mathfrak E=b_1(X)c_1(Y)$ the correction equals $-\dfrac{\abs{S^{d-3}}}{R^2\abs{S^{d-1}}}M_{b_1}M_{1-X^2}T_{d+2,R}M_{c_1}$, where $M_g$ denotes multiplication by $g$.
\end{enumerate}
\end{proposition}

\begin{proof}
(a) With $\phi=p(u)$, $s=\cos\chi$ and $\hat r=\sec\alpha$, multiplying $\Disc=\hat r^2\sin^2\phi\sin^2\chi-(1-\hat r\cos\chi\cos\phi)^2$ by $\cos^2\alpha$ and writing $\sin\phi\sin\chi$ and $\cos\phi\cos\chi$ through $\cos(\phi\mp\chi)$ gives $\Disc=\sec^2\alpha\,[\cos(\phi-\chi)-\cos\alpha][\cos\alpha-\cos(\phi+\chi)]$. The product-to-sum formulas show that this equals $\sec^2\alpha\,[\cos(\chi-\alpha)-\cos\phi][\cos\phi-\cos(\chi+\alpha)]$. Finally $\cos(\chi-\alpha)(y)=\cos p(y-\theta)=G(\xi_-)$, $\cos(\chi+\alpha)=G(\xi_+)$ and $\cos\phi=G(X)$.

(b) In the coarea formula \eqref{eq:coarea} substitute $s=\cos\chi(y)$, $ds=-\sin\chi\,\chi'\,dy$, $r(s)=\hat r(y)$, and use (a), writing $\Pi=[G(\xi_-)-G(X)][G(X)-G(\xi_+)]$. Since $G(\xi_-)-G(\xi_+)=2\sin\chi\sin\alpha$,
\[
(\sin\alpha)^{-\beta}\Pi^\beta=2^\beta(\sin\chi)^\beta\Bigl[\frac{\Pi}{G(\xi_-)-G(\xi_+)}\Bigr]^\beta .
\]
With the weights \eqref{eq:Lnorm}, all powers of $\hat r$ cancel because $2-2\lambda+2\beta=0$, the factors $\chi'$ cancel, and the powers of $\sin p$ and $\sin\chi$ cancel because $1-\lambda+\beta=0$. What remains is
\begin{gather*}
\mathcal LF(u)=a_0\abs{S^{d-3}}2^\beta\int\mathcal K^\beta\Bigl(\frac{\sin y}{\sin u}\Bigr)^\lambda F\,dy,\\
\mathcal K=\frac{\Pi}{[G(\xi_-)-G(\xi_+)]\sin p(u)\,p'(u)} .
\end{gather*}
Since $\sin p\,p'=\sin u\,G'(X)$, the ratio of $\mathcal K$ to its value $\mathcal K_0$ for $G=\id$ is exactly $\Psi$. For $G=\id$ (the ball) the weights are constant and $\mathcal L=T_{d,R}$; substituting $dY=-\sin y\,dy$ and $\Pi=\Disc_0/R^2$ shows that the kernel with $\mathcal K_0$ is the round coarea kernel $\abs{S^{d-3}}R^2(1-X^2)^{-(\lambda-1/2)}\Disc_0^\beta$. This gives \eqref{eq:Lkernel}. The identity $\Disc_0=R^2(\xi_--X)(X-\xi_+)$ follows from $\xi_-+\xi_+=2Y\cos\theta$, $\xi_-\xi_+=Y^2-\sin^2\theta$, $R\cos\theta=1$ and $R^2\sin^2\theta=R^2-1$. The integration domains agree because $G$ is increasing near the identity.

(c) $\Psi$ is analytic near $1$ and symmetric in $(\xi_-,\xi_+)$; hence it is an analytic function $\Phi(X,\varsigma_1,\varsigma_2)$ of $X$ and of $\varsigma_1=\xi_-+\xi_+$, $\varsigma_2=\xi_-\xi_+$, including at $\xi_-=\xi_+$. Since $G[a,a]=G'(a)$, $\Psi=1$ when $X=\xi_-$ or $X=\xi_+$, i.e.\ $\Phi-1$ vanishes on the hypersurface $\varsigma_2=\varsigma_1X-X^2$. Division in the variable $\varsigma_2$ gives $\Phi-1=(\varsigma_2-\varsigma_1X+X^2)\,\mathfrak E_1$ with $\mathfrak E_1=\int_0^1\partial_{\varsigma_2}\Phi(X,\varsigma_1,\varsigma_1X-X^2+\tau(\varsigma_2-\varsigma_1X+X^2))\,d\tau$ analytic, and $\varsigma_2-\varsigma_1X+X^2=(X-\xi_-)(X-\xi_+)$. Put $\mathfrak E=\mathfrak E_1\int_0^1\beta[1+\tau(\Psi-1)]^{\beta-1}d\tau$; as $\varsigma_1=2Y\cos\theta$ and $\varsigma_2=Y^2-\sin^2\theta$, it is an analytic function of $(X,Y)$, including at the meridians where $\xi_-=\xi_+$. For $G=\id$ we have $\Psi\equiv1$ and $\mathfrak E=0$. Since $\Psi^\beta-1=-\Disc_0\mathfrak E/R^2$, \eqref{eq:Lkernel} becomes \eqref{eq:dd2}. For separated amplitudes use the coarea expression of $T_{d+2,R}$, whose denominator power is $\lambda+\frac12$ and whose numerator is $\abs{S^{d-1}}R^2\Disc_0^\lambda$.

The identities are first proved on a nondegenerate output interval. The left side of \eqref{eq:dd2} is globally analytic by \eqref{eq:Tr} and the regular weights, the right side by the fixed-ray formula for $T_{d+2,R}$ and the convergent expansion of $\mathfrak E$ below, and analytic continuation on the connected common domain gives the identity through every meridian and turning output. It is an identity of operators.
\end{proof}

\begin{lemma}[Physical section map]\label{lem:physmap}
Let $\mathcal X(Y)=\hat r(y)\cos\chi(y)$ and $S_G(X)=(1-G(X)^2)/(1-X^2)$, which is analytic, positive and close to $1$. Then
\begin{equation}\label{eq:Qp}
\begin{gathered}
\cV_p(u):=A_r(p(u))=\frac{\kappa_{d-2}}{\abs{S^{d-1}}R^d}\,S_G(X)^{-\lambda-\frac12}\,T_{d+2,R}\bigl[\Omega(X,\cdot)\bigr](X),\\
\Omega(X,Y)=\hat r(Y)^{2\lambda}\mathcal X'(Y)\bigl[G[\xi_-,X]G[X,\xi_+]\bigr]^\lambda .
\end{gathered}
\end{equation}
For complex $u$ on the working strip the fixed-ray argument $Y=\cos\theta\cos u-\sin\theta\sin u\,z$, $z\in[-1,1]$, of $T_{d+2,R}$ stays in the same Bernstein ellipse, so \eqref{eq:Qp} has a common analytic domain.
\end{lemma}

\begin{proof}
By \cref{lem:slice} in the variable $y$, $A_r(\phi)=\kappa_{d-2}(\sin\phi)^{-2\lambda-1}\int\Disc^\lambda\,d\mathcal X$. Insert $\Disc=\hat r^2R^{-2}\Disc_0\,G[\xi_-,X]G[X,\xi_+]$ from \cref{prop:oddnorm}(a), $(\sin p)^{2}=S_G(X)(1-X^2)$, and the coarea expression $T_{d+2,R}\Omega(X)=\abs{S^{d-1}}R^2(1-X^2)^{-\lambda-1/2}\int\Disc_0^\lambda\Omega\,dY$; the power of $R$ is $2\lambda+2=d$. At the ball $\Omega=R^{2\lambda+1}$ and $T_{d+2,R}[1]=\abs{S^d}R(R^2-1)^{(d-1)/2}$, and \eqref{eq:Qp} reduces to $\kappa_{d-1}(R^2-1)^{(d-1)/2}$, as it should. The point $Y$ is a real linear combination $a\cos u+b\sin u$ with $a^2+b^2=\cos^2\theta+z^2\sin^2\theta\le1$, i.e.\ a real contraction of a real translate of $\cos u$; Bernstein ellipses are convex and invariant under these operations.
\end{proof}

\begin{proposition}[Matrix class (O)]\label{prop:Oclass}
Let $L^2_\lambda=L^2((-1,1),(1-X^2)^{\lambda-1/2}dX)$ with its orthonormal Gegenbauer basis. Fix $s_d=\lceil\lambda\rceil+5$ and $r_d=\lceil\lambda\rceil+8$. Suppose the coefficients of $\mathfrak E=\sum_{k,l}e_{kl}T_k(X)T_l(Y)$ satisfy $\sum_{k,l}\abs{e_{kl}}(1+k+l)^{s_d}e^{\rho(k+l)}\le C_dh$; by two contour shifts and $r_d$ integrations by parts in each angular variable this follows from bounds on the mixed angular derivatives of order $r_d$ of $\mathfrak E(\cos u,\cos y)$ on the closed polystrip of width $\rho$, because $r_d>s_d+2$. Then
\begin{equation}\label{eq:matrixO}
\begin{gathered}
\frac{\Lambda_0^{\lambda+1}\mathcal L}{\kappa_d(R)}=D+A_r,\qquad \kappa_d(R)=\abs{S^{d-2}}R(R^2-1)^{\frac{d-3}2},\\
d_n=\frac{(1+n)^{\lambda+1}C_n^\lambda(\cos\theta)}{C_n^\lambda(1)},\qquad N_\rho(A_r)\le C_dh,
\end{gathered}
\end{equation}
and if the coefficient fields have first $\omega$-jets of size $O_d(1)$, then $\partial_\omega A_r=A_0+\Lambda_0A_1$ with $N_\rho(A_0)\le C_d$, $N_\rho(A_1)\le C_dh$. Reality and the even sector are preserved. At fixed $\omega$, $D$ and the round factor of order $d+2$ are fixed, so $A_r$ is Lipschitz in the coefficient fields at the same radius.
\end{proposition}

\begin{proof}
The inclusion $\mathsf I:L^2_\lambda\to L^2_{\lambda+1}$ is a contraction with adjoint $\mathsf I^*g=(1-X^2)g$, and the connection formula $C_n^\lambda=\frac\lambda{n+\lambda}(C_n^{\lambda+1}-C_{n-2}^{\lambda+1})$ \cite[(18.9.7)]{DLMF} makes $\mathsf I$ two-diagonal in the orthonormal bases; hence $M_{1-X^2}T_{d+2,R}=\mathsf I^*\operatorname{diag}(\ell_{d+2,n})\mathsf I$ with $\ell_{d+2,n}=\abs{S^d}R(R^2-1)^{(d-1)/2}C_n^{\lambda+1}(1/R)/C_n^{\lambda+1}(1)$. For fixed $d$ and $\theta$ in a compact subinterval of $(0,\pi/2)$, the fixed-parameter Darboux expansion of Jacobi polynomials \cite[(18.15.4\_5)]{DLMF}, specialized to Gegenbauer polynomials as in \cite[\S18.15(ii)]{DLMF}, gives $\abs{C_N^\mu(\cos\theta)}\le C_\mu N^{\mu-1}$; with $C_N^\mu(1)=\binom{N+2\mu-1}{N}$ this gives $\abs{\ell_{d+2,n}}\le C_d(1+n)^{-\lambda-1}$. Differentiating the exact formula first, by $\frac{d}{dx}C^\mu_N=2\mu C^{\mu+1}_{N-1}$ \cite[(18.9.19)]{DLMF}, and then applying the same bound gives $\abs{\partial_\omega\ell_{d+2,n}}\le C_d(1+n)^{-\lambda}$. Thus $\Lambda_0^{\lambda+1}$ times the correction in \eqref{eq:dd2} is a band matrix composed with multiplications, and its frequency derivative costs one power of the mode. Multiplication by $T_k$ has bandwidth $k$, operator norm $\le1$ on $L^2_\lambda$, and hence $N_\rho(M_{T_k})\le(2k+1)e^{\rho k}$; conjugation by $\Lambda_0^q$ costs at most $(1+k)^{\abs q}$ more. Summing over $(k,l)$ gives \eqref{eq:matrixO}. The round part $T_{d,R}$ is diagonal by \cref{lem:round}, with $L_n(R)/\kappa_d(R)=C_n^\lambda(\cos\theta)/C_n^\lambda(1)$. The frequency derivative falls on the coefficient fields (giving $A_0$), on $\ell_{d+2,n}$ (giving $\Lambda_0A_1$ with $A_1=O(h)$) and on the scalar $\kappa_d(R)$ (giving $A_0$).
\end{proof}

\subsection{Even dimensions: primitive, regularizer and endpoint formula}
Let $d=2\lambda+2\ge6$ be even, $k=\lambda-1$ and $k_\lambda=2^kk!$, as in \cref{lem:darboux}. For a smooth function $H$ on $(-1,1)$ the identity
\begin{equation}\label{eq:rodrigues}
\prod_{j=1}^k(\partial_t^2+j^2)\bigl[\sin t\,H(\cos t)\bigr]=\partial_c^k\bigl[(1-c^2)^{k+\frac12}H^{(k)}(c)\bigr]\Big|_{c=\cos t}
\end{equation}
holds. Indeed, for $k=0$ both sides equal $\sin t\,H(\cos t)$; if it holds for $k$, put $\mathsf u=(1-c^2)^{k+1/2}H^{(k)}$, so that $(1-c^2)^{k+3/2}H^{(k+1)}=(1-c^2)\mathsf u'+(2k+1)c\,\mathsf u$, and the Leibniz rule gives $\partial_c^{k+1}[(1-c^2)\mathsf u'+(2k+1)c\mathsf u]=[(1-c^2)\partial_c^2-c\,\partial_c+(k+1)^2]\partial_c^k\mathsf u$, which is $(\partial_t^2+(k+1)^2)\partial_c^k\mathsf u$ since $\partial_t^2=(1-c^2)\partial_c^2-c\,\partial_c$.

\emph{Input primitive.} For an even radial variation $f$ (a function of the axial cosine $s$) and the profile $r$, let $H$ be a $k$-fold primitive, $H^{(k)}=r^{2\lambda}f$, and put $p_f(t)=k_\lambda\sin t\,H(\cos t)$. The polynomial ambiguity of $H$ (degree $<k$) changes $p_f$ by a combination of $\sin(jt)$, $j\le k$, of the parity of $\lambda$; we fix it by requiring that these $\lfloor(\lambda-1)/2\rfloor$ sine coefficients of $p_f$ vanish. The \emph{native primitive} is $v=p_f\circ\chi$. At the ball, by \cref{lem:darboux}(a), $f=R^{-2\lambda}C^\lambda_{m-\lambda}$ corresponds to $p_f=\sin(mt)$.

\emph{Output regularizer.} For an even function $g$ of the normal angle $\phi$ let $S_\lambda g=k_\lambda\sin\phi\,H_g(\cos\phi)$ with $H_g^{(k)}=g$, and $P_\lambda=\partial_\phi\prod_{j=1}^k(\partial_\phi^2+j^2)$. By \eqref{eq:rodrigues},
\begin{equation}\label{eq:PS}
P_\lambda S_\lambda g(\phi)=-k_\lambda\sin\phi\;\partial_c^{\lambda}\bigl[(1-c^2)^{\lambda-\frac12}g(c)\bigr]_{c=\cos\phi},
\end{equation}
which does not depend on the polynomial ambiguity of $H_g$. Since $P_\lambda\sin(mt)=(-1)^km\prod_{j=1}^k(m^2-j^2)\cos(mt)$, the regularized output has no cosine modes $\cos(jt)$ with $j\le k$ of the parity of $\lambda$.

\emph{Native operator.} Put $\Phi=\chi+\alpha$, so that $\Phi(y)=p(y+\theta)$ and $(\chi-\alpha)(y+\omega)=\Phi(y)$. Let $T$ be the translation $(Tv)(y)=v(y+\omega)$, $\Delta=I-T$, $P_0$ the angular mean (zero on antiperiodic functions), $\Delta^\dagger$ the zero-mean inverse of $\Delta$ on periodic and the inverse on antiperiodic functions (where it exists), and $\cA=(I+T)\Delta^\dagger$. Then $\cA$ vanishes on constants and $\cB:=\cA\partial_y$ acts on $\sin(my)$ and $\cos(my)$ as multiplication by $-m\cot(m\theta)$. The native derivative is
\[
\mathcal Lv(y)=\Phi'(y)\,\bigl(P_\lambda S_\lambda T_rf\bigr)(\Phi(y)),\qquad v=p_f\circ\chi .
\]
With $\mathfrak q(c,s)=1-s^2-r(s)^{-2}+2(s/r(s))c-c^2$ we have $\Disc=r^2\mathfrak q$ in \eqref{eq:coarea}, and we put $\mathcal Q=\partial_c^\lambda\mathfrak q^k$, a polynomial in $c$ of degree at most $\lambda-2$.

\begin{lemma}[Endpoint formula]\label{lem:endpoint}
Let $I(c)=\int f\,r^2\Disc^k\,ds$ be the coarea numerator, the integral being over $[\mathsf a(c),\mathsf b(c)]=\{\Disc\ge0\}$. Then $P_\lambda S_\lambda(T_rf)(\phi)=-k_\lambda\abs{S^{2\lambda-1}}\sin\phi\,I^{(\lambda)}(\cos\phi)$ and
\begin{equation}\label{eq:endpoint}
I^{(\lambda)}=\bigl[k!\,f\,r^2(\partial_c\Disc)^k\,\mathsf s'(c)\bigr]_{\mathsf a}^{\mathsf b}+\int_{\mathsf a}^{\mathsf b} H^{(k)}(s)\,\mathcal Q(c,s)\,ds ,
\end{equation}
where $\mathsf s'(c)$ denotes the derivative of the respective endpoint. Integrating the last term $k$ times by parts gives finitely many local endpoint derivatives of $H$ of order $\le k-1$ and finitely many terms $J_i(y)\int_y^{y+\omega}g_i(\tilde y)v(\tilde y)\,d\tilde y$ with $J_i=\Phi'\sin\Phi\,(\cos\Phi)^i$, $0\le i\le\lambda-2$, and real analytic native coefficients $g_i$. Consequently
\begin{equation}\label{eq:Lendpoint}
\mathcal Lv=E_+v-T(E_-v)+\sum_iJ_i\int_y^{y+\omega}g_iv,
\end{equation}
where $E_\pm$ are differential operators of order $k$ in the native variable, with coefficients analytic away from the meridians.
\end{lemma}

\begin{proof}
The first identity is \eqref{eq:PS} combined with \eqref{eq:coarea}, since $(\sin\phi)^{2\lambda-1}=(1-c^2)^{\lambda-1/2}$. Since $\Disc^k$ vanishes to order $k$ at both endpoints, $I^{(j)}=\int f r^2\partial_c^j\Disc^k\,ds$ for $j\le k$; one more derivative produces the boundary term, where $\partial_c^k\Disc^k=k!(\partial_c\Disc)^k$, and the interior term $\int fr^2\partial_c^\lambda\Disc^k\,ds=\int H^{(k)}\mathcal Q\,ds$ since $fr^{2+2k}=H^{(k)}$. Integrating by parts $k$ times,
\[
\int H^{(k)}\mathcal Q\,ds=\sum_{i=0}^{k-1}(-1)^i\bigl[H^{(k-1-i)}\partial_s^i\mathcal Q\bigr]_{\mathsf a}^{\mathsf b}+(-1)^k\int H\,\partial_s^k\mathcal Q\,ds .
\]
In native coordinates $H(\cos\chi)=v/(k_\lambda\sin\chi)$ and $H\,ds=-v\,\chi'\,dy/k_\lambda$, so no division by $\sin\chi$ remains in the residual integral. The upper endpoint $\mathsf b$ (larger $s$) is the profile point with $\chi+\alpha=\phi$, i.e.\ $y=x$ when $\phi=\Phi(x)$; the lower endpoint is the point with $\chi-\alpha=\phi$, i.e.\ $y=x+\omega$. Writing $\partial_s^k\mathcal Q=\sum_{i\le\lambda-2}c^i\mathfrak g_i(s)$ and multiplying by the output factor $-k_\lambda\abs{S^{2\lambda-1}}\Phi'\sin\Phi$ gives \eqref{eq:Lendpoint}.
\end{proof}

\begin{lemma}[Removability at the meridians]\label{lem:meridian}
The coefficients of the differential operators $E_\pm$ are real analytic; in particular they have no poles at the meridians.
\end{lemma}

\begin{proof}
It suffices to work near $t=0$. The profile is even and analytic, hence analytic in $s=\cos t$ near $s=1$; put $\mathsf y=1-s$. Let $\mathsf b(c)$ be the endpoint crossing the meridian and $\delta(c)=1-\mathsf b(c)$. The root is simple: $\partial_s\Disc=-2(r_{\mathrm{tip}}^2-1)\ne0$ at the meridian, so $\Disc=(s-\mathsf b(c))\,\mathsf d_1(s,c)$ with $\mathsf d_1$ jointly analytic and nonzero. The endpoint angle $t_{\mathsf b}$ is analytic in $c$ with nonzero derivative, since $c=\cos(t_{\mathsf b}+\alpha(t_{\mathsf b}))$ has nonzero $t$-derivative at the crossing, and $\delta=t_{\mathsf b}^2$ times an analytic unit.

Test the local endpoint expression on $p(t)=t^j$, $0\le j\le k$. For odd $j$, $H=p/(k_\lambda\sin t)$ is analytic in $s$ and all moving-endpoint integrals are analytic. For even $j$ put $\beta_j=(j-1)/2\ge-\frac12$: the local expansion of $H$ has powers $\mathsf y^{\beta_j+i}$, $i\ge0$; after $k$ axial derivatives they become $\mathsf y^{\beta_j-k+i}$. Expand the analytic unit $\mathsf d_1^k$ in nonnegative powers $\mathsf y^{\iota}$ and $(\delta-\mathsf y)^k=\sum_l\binom kl\delta^{k-l}(-\mathsf y)^l$. A primitive term evaluated at the moving endpoint has the power $\delta^{k-l}\delta^{\beta_j-k+i+\iota+l+1}=\delta^{\beta_j+i+\iota+1}$, with the nonzero half-integer integration denominator $\beta_j-k+i+\iota+l+1$; so no logarithms occur, and the endpoint exponent starts at $\beta_j+1=(j+1)/2>0$. Since $\delta=t_{\mathsf b}^2$ times a unit, these are integer powers $t_{\mathsf b}^{\,j+1+2i+2\iota}$ of positive degree. Hence the whole finite-part primitive is analytic before the $\lambda$ derivatives in $c$ are taken. The residual integral is analytic for every test since $H\,ds=-p\,dt/k_\lambda$. An operator of order $k$ with meromorphic coefficients whose values on $1,t,\dots,t^k$ are analytic has analytic coefficients (recover the coefficients triangularly from $E(1)$, $E(t)$, \dots). The finite-part primitive is used only for these test functions; on physical inputs all integrals are ordinary convergent integrals.
\end{proof}

\begin{proposition}[Cohomological pull-through]\label{prop:pullthrough}
Write $E_+-TE_-=\Delta\mathsf M+(I+T)\mathsf N$ with $\mathsf M=(E_++E_-)/2$ and $\mathsf N=(E_+-E_-)/2$. For an analytic coefficient $J$ put $M_J=(J+T^{-1}J)/2$ and $Y_J=(J-T^{-1}J)/2$, as multiplication operators. Then, on the relevant parity sector,
\begin{equation}\label{eq:pull}
\begin{gathered}
\Delta^\dagger J\Delta=(I-P_0)M_J+\cA Y_J,\qquad \Delta^\dagger J_i=\tfrac12(\cA+I)J_i,\\
\int_y^{y+\omega}u=-\Delta\partial_y^\dagger(u-P_0u)+\omega P_0u,
\end{gathered}
\end{equation}
where $\partial_y^\dagger$ is the zero-mean primitive. Consequently the normalized derivative is
\begin{equation}\label{eq:KVBW}
K_{\mathrm{raw}}:=\Delta^\dagger\mathcal L=V+\cA C=V+\cB W,\qquad W=\partial_y^\dagger(I-P_0)C,
\end{equation}
where $V$, $C$, $W$ are finite sums of differential operators, multiplications, zero-mean primitives, translations and finite-rank mean operators with real analytic coefficients depending on the native profile; they contain no cotangent multiplier.
\end{proposition}

\begin{proof}
Direct multiplication gives $\Delta M_J+(I+T)Y_J=J\Delta$ and $P_0J\Delta=2P_0Y_J$. Since $\Delta\cA=I+T-2P_0$, applying $\Delta$ to the right side of the first identity gives $(I-P_0)J\Delta$, and both sides have zero mean. The interval integral splits as stated because $\Delta\partial_y^\dagger(u-P_0u)$ is minus the integral of $u-P_0u$ over $[y,y+\omega]$. Each $J_i=-(\cos^{i+1}\Phi)'/(i+1)$ is a derivative, so $P_0J_i=0$ and $\Delta^\dagger J_i=\frac12(\cA+I)J_i$. Combining these with the residual terms $J_i\int_y^{y+\omega}g_iv$, the term $\omega P_0(g_iv)\cdot\frac12(\cA+I)J_i$ equals $\cA[\frac\omega2P_0(g_iv)J_i]+\frac\omega2P_0(g_iv)J_i$ because $P_0(g_iv)$ is a scalar; it contributes the same rank-one operator to $V$ and to $C$. With $\Delta^\dagger\Delta=I-P_0$ and $\Delta^\dagger(I+T)=\cA$ for the local terms we obtain $K_{\mathrm{raw}}=V+\cA C$. Finally $\cA$ kills constants and commutes with $\partial_y$, so $\cA C=\cA\partial_y\partial_y^\dagger(I-P_0)C=\cB W$.
\end{proof}

\subsection{Even dimensions: orders and ball symbols}
For an operator of the type in \cref{prop:pullthrough} the \emph{order} is the highest order of differentiation that occurs, primitives counting $-1$.

\begin{lemma}[Orders and ball symbols]\label{lem:orders}
\begin{enumerate}[label=\textup{(\alph*)}]
\item For $\sigma=\pm$ the coefficient of $\partial_y^k$ in $E_\sigma$ is $\mathfrak c_\sigma=-k_\lambda\abs{S^{2\lambda-1}}(-\sigma)^k\sin(\alpha)^k(\chi')^{1-k}$, evaluated at the native point of the endpoint. Consequently, if $\lambda$ is even ($k$ odd), $C$ has order $k$ with nonvanishing leading coefficient, $W$ has order $\lambda-2$ and $V$ has order at most $k-1$; if $\lambda$ is odd ($k$ even), $V$ has order $k$, $C$ has order at most $k-1$ and $W$ has order at most $\lambda-3$.
\item At the ball of radius $R=\sec\theta$ the operators $E_\pm$, $\mathsf M$, $\mathsf N$, $V$, $C$, $W$ have constant coefficients, the residual terms vanish, and after division by the nonzero scalar $\mathfrak n(\theta)$ below, $V$ and $W$ act on $\sin(my)$, $m\ge\lambda$ of the parity of $\lambda$, as multiplication by
\begin{equation}\label{eq:ballsymbols}
V_0(m)=a_k(m,\cot\theta),\qquad W_0(m)=-b_k(m,\cot\theta)/m .
\end{equation}
\item Consequently $W_0$ is a polynomial in $m$ of exact degree $\lambda-2$ ($\lambda$ even), resp.\ $\lambda-3$ ($\lambda$ odd), and $q^\circ_m(\theta):=V_0(m)/W_0(m)$ satisfies
\begin{equation}\label{eq:Q0}
q^\circ_m(\theta)=\begin{cases}S_k\cot\theta+O(m^{-2}),&\lambda\text{ even},\\-\dfrac{\tan\theta}{S_k}\,m^2+O(1),&\lambda\text{ odd},\end{cases}
\end{equation}
uniformly for $\theta$ in compact subsets of $(0,\pi/2)$, together with the $\theta$-derivative.
\end{enumerate}
\end{lemma}

\begin{proof}
(a) At the endpoint $\sigma$ we have $\phi=\chi+\sigma\alpha$, $s/r=\cos\chi\cos\alpha$, hence $\partial_c\Disc=r^2\partial_c\mathfrak q=2r^2(\cos\chi\cos\alpha-\cos(\chi+\sigma\alpha))=2\sigma r^2\sin\chi\sin\alpha$. The endpoint derivative is $\mathsf s'(c)=\sin\chi\,\chi'/(\sin\Phi\,\Phi')$; for the lower endpoint this uses $(\chi-\alpha)'(y+\omega)=\Phi'(y)$. The highest term of $H^{(k)}$ is $(-1)^k\partial_y^kv/(k_\lambda(\chi')^k\sin^{k+1}\chi)$, because $d/ds=-(\chi'\sin\chi)^{-1}\partial_y$. Hence the boundary term of \eqref{eq:endpoint} has highest part
\begin{multline*}
k!\,\frac{(-1)^k\partial_y^kv}{k_\lambda(\chi')^k\sin^{k+1}\chi}\,r^{2-2\lambda}\,(2\sigma r^2\sin\chi\sin\alpha)^k\,\frac{\sin\chi\,\chi'}{\sin\Phi\,\Phi'}\\
=\frac{(-\sigma)^k\sin^k\alpha\,(\chi')^{1-k}}{\sin\Phi\,\Phi'}\,\partial_y^kv,
\end{multline*}
using $r^{2-2\lambda+2k}=1$ and $k!\,2^k=k_\lambda$. Multiplying by the output factor $-k_\lambda\abs{S^{2\lambda-1}}\sin\Phi\,\Phi'$ gives $\mathfrak c_\sigma$. All other endpoint terms involve $H^{(j)}$ with $j\le k-1$, and the residual terms have order $-1$. As functions of the native variable, $\mathfrak c_+=(-1)^k\mathfrak c_-$; therefore the top coefficient of $\mathsf M$ vanishes for odd $k$ and that of $\mathsf N$ for even $k$. By \cref{prop:pullthrough}, $V$ contains $\mathsf M$ and terms of order $-1$, $C$ contains $\mathsf N$ and terms of order $-1$, and $W=\partial_y^\dagger(I-P_0)C$. This gives the orders. That $W$ has exactly the stated order near the ball follows from (b) and (c).

(b) At the ball $\chi=\id$, $\alpha=\theta$, $\Phi(y)=y+\theta$ and $r\equiv R$. Then $\mathfrak q=\sin^2\theta(1-c^2)-(s-c\cos\theta)^2$ has total degree two in $(c,s)$, so $\mathcal Q$ has degree at most $k-1$ in $s$ and $\partial_s^k\mathcal Q=0$: the residual terms vanish. By \cref{lem:round}, \cref{lem:darboux}(a), $P_\lambda\sin(mt)=(-1)^km\prod_j(m^2-j^2)\cos(mt)$ and $C^\lambda_{m-\lambda}(1)=m\prod_{j=1}^k(m^2-j^2)/(2\lambda-1)!$, the regularized round operator maps the primitive $p=\sin(m\,\cdot)$ to
\begin{gather*}
\kappa(\theta)\,[\mathsf D_k\cdots\mathsf D_1\sin(m\,\cdot)](\theta)\,\cos(m\phi),\\
\kappa(\theta)=\frac{(-1)^k(2\lambda-1)!\,\abs{S^{d-2}}R^{1-2\lambda}(R^2-1)^{\frac{d-3}2}}{k_\lambda\sin^\lambda\theta}\ne0 .
\end{gather*}
Since $\mathsf D_j$ acts on the variable $\theta$ and $\sin(m\theta)\cos(m\phi)=\frac12[\sin m(\phi+\theta)-\sin m(\phi-\theta)]$, this equals $\frac{\kappa(\theta)}2(\mathsf D_k\cdots\mathsf D_1)_\theta[p(\phi+\theta)-p(\phi-\theta)]=\frac{\kappa(\theta)}2\sum_{j=0}^k\mathfrak d_j(\theta)\bigl[(-1)^jp^{(j)}(\phi+\theta)-p^{(j)}(\phi-\theta)\bigr]$, where $\mathsf D_k\cdots\mathsf D_1=\sum_j\mathfrak d_j(\theta)(-\partial_\theta)^j$ with $\mathfrak d_k=1$. On the other hand, by \cref{lem:endpoint} the same operator is $\sum_j[\alpha_j(\phi)p^{(j)}(\phi-\theta)+\beta_j(\phi)p^{(j)}(\phi+\theta)]$ with coefficients analytic in $\phi$. The difference of the two expressions vanishes on $p=\sin(m\,\cdot)$ for all $m$ in the progression $m_0+2\N$; writing $\partial^j\sin(mt)=\operatorname{Im}((im)^je^{imt})$, it is a combination $\sum_\zeta P_\zeta(m)\zeta^m$ with polynomial coefficients $P_\zeta$ and $\zeta\in\{e^{\pm i(\phi-\theta)},e^{\pm i(\phi+\theta)}\}$. For $\phi$ outside a discrete set the four numbers $\zeta^2$ are distinct, and the sequences $\ell\mapsto\ell^i\mu^\ell$ with distinct $\mu$ are linearly independent; hence all $P_\zeta$ vanish, i.e.\ $\alpha_j(\phi)=-\frac{\kappa}2\mathfrak d_j(\theta)$ and $\beta_j(\phi)=\frac\kappa2(-1)^j\mathfrak d_j(\theta)$ for all $\phi$. (For $j=k$ this agrees with (a).) Thus $E_\pm$ have constant coefficients at the ball, and so do $\mathsf M$, $\mathsf N$; then $V=(I-P_0)\mathsf M$, $C=\mathsf N$ and $W$ are Fourier multipliers. Put $\mathfrak n(\theta)=-\kappa(\theta)/2$. On $\sin(my)$, $\Delta\sin(my)=-2\sin(m\theta)\cos(m(y+\theta))$, so the full normalized operator acts as multiplication by $\mathfrak n(\theta)[a_k(m)+b_k(m)\cot(m\theta)]$, using \cref{lem:darboux}(b). Writing $\mathsf M\sin(my)=\mu_1(m)\sin(my)+\mu_2(m)\cos(my)$ and $\mathsf N\sin(my)=\nu_1(m)\sin(my)+\nu_2(m)\cos(my)$ with polynomials $\mu_i,\nu_i$, and using that $\cB$ is multiplication by $-m\cot(m\theta)$ on sines and cosines, we obtain
\begin{multline*}
[\mu_1(m)-\cot(m\theta)\nu_2(m)]\sin(my)+[\mu_2(m)+\cot(m\theta)\nu_1(m)]\cos(my)\\
=\mathfrak n(\theta)[a_k(m)+b_k(m)\cot(m\theta)]\sin(my).
\end{multline*}
If $\theta/\pi$ is irrational, $\cot(m\theta)$ takes values arbitrarily close to every real number along $m\in m_0+2\N$, which is impossible for a nonzero rational function of $m$; hence $\mu_2=\nu_1=0$, $\mu_1=\mathfrak n a_k$ and $\nu_2=-\mathfrak nb_k$. This gives \eqref{eq:ballsymbols} for $\theta/\pi$ irrational, and for all $\theta$ by continuity, since all symbols are polynomial in $m$ with coefficients analytic in $\theta$.

(c) follows from \cref{lem:darboux}(c): for $k=2\ell-1$, $W_0=-(-1)^\ell m^{k-1}+O(m^{k-3})$ and $V_0=(-1)^{\ell-1}S_k\cot\theta\,m^{k-1}+O(m^{k-3})$; for $k=2\ell$, $W_0=-(-1)^\ell S_k\cot\theta\,m^{k-2}+O(m^{k-4})$ and $V_0=(-1)^\ell m^k+O(m^{k-2})$. The quotients are rational in $m$ with coefficients polynomial in $\cot\theta$, which gives the uniformity in $\theta$.
\end{proof}

From now on we replace $\mathcal L$ by $\mathfrak n(\theta)^{-1}\mathcal L$ and $K_{\mathrm{raw}}$ by $\mathfrak n(\theta)^{-1}K_{\mathrm{raw}}$; this does not change the equations, and the ball symbols are \eqref{eq:ballsymbols}.

\subsection{Even dimensions: localization and the matrix classes (B) and (E)}
We work in the sine basis $\sin(my)$ of the native primitive, $m\ge1$ of the parity of $\lambda$, and in the Wiener norms $\norm a_{\rho,q}$ of coefficient functions.

\begin{lemma}[Same-radius localization]\label{lem:local}
Let $a$ be an analytic $\pi$-periodic or $\pi$-antiperiodic coefficient with $\norm a_{\rho,q}<\infty$, let $j\ge0$, and let $\mathfrak w_{\mathrm{in}}(n)$, $\mathfrak w_{\mathrm{out}}(m)$ be positive weights with $\mathfrak w_{\mathrm{out}}(m)\,n^j\,\mathfrak w_{\mathrm{in}}(n)\le C_0(1+\abs{m-n})^q$ for all $m,n\ge1$. Let $B$ be the matrix of $v\mapsto\mathfrak w_{\mathrm{out}}\cdot\Pi_{\sin}[a\,\partial_y^j(\mathfrak w_{\mathrm{in}}v)]$ in the sine basis, where $\Pi_{\sin}$ is the projection onto the sine modes and the weights act diagonally on the modes. Then $N_\rho(B)\le4C_0\norm a_{\rho,q}$. The same bound holds with $\partial_y^j$ replaced by $\partial_y^{\dagger}$ (with $j=-1$), by a translation $T^{\pm1}$, or with $a$ replaced by a finite-rank operator $v\mapsto P_0(gv)\,a$.
\end{lemma}

\begin{proof}
We have $a(y)\sin(ny)=\sum_k\hat a_k(e^{i(k+n)y}-e^{i(k-n)y})/(2i)$, whose $\sin(my)$ coefficient involves only $\hat a_k$ with $\abs k\in\{\abs{m-n},m+n\}$: the first case gives Toeplitz entries, the second Hankel entries. In both cases $e^{\rho\abs{m-n}}\le e^{\rho\abs k}$ and, since $\abs{m-n}\le\abs k$, the weight factor is at most $C_0(1+\abs k)^q$. For fixed $m$ and each $k$ at most two $n$ occur, so each row sum is at most $4C_0\sum_k(1+\abs k)^qe^{\rho\abs k}\abs{\hat a_k}$, and the same holds for column sums. The derivative contributes the factor $n^j$ already included; primitives contribute $n^{-1}$, translations unimodular factors, and the finite-rank case is a product of two bounded analytic vectors.
\end{proof}

The estimate uses the coefficient function at the same radius $\rho$, with $q$ additional powers of the mode, i.e.\ finitely many derivatives of the coefficient; no width is lost. It is distinct from the composition estimates of \cref{sec:newton}, which lose width.

\begin{proposition}[Matrix classes (B) and (E)]\label{prop:evenclass}
Choose a finite cutoff $m_\sharp$ above every zero of the polynomial $W_0(m)$ (uniformly for $\theta$ near $\theta_0$) and above all $m=n+\lambda$ with $n$ resonant, and let $F$ be a fixed diagonal operator supported on the modes $m<m_\sharp$ such that $W_0+F$ is invertible on the whole sine sector. Put $\widetilde W=W+F$. Suppose the native coefficient fields of $V,C$ are within $O(h)$ of their values at the ball of radius $\sec\theta$ in the norm $\norm\cdot_{\rho,q_W}$ for a fixed $q_W$ and all $\rho\in[\rho_\infty,\rho_{\max}]$, with first $\omega$-jets of size $O(1)$. Then:
\begin{enumerate}[label=\textup{(\roman*)}]
\item $\widetilde W(W_0+F)^{-1}=I+O(h)$ in $N_\rho$; hence $\widetilde W$ has a Neumann inverse of order minus the order of $W_0$, without rotation small divisors.
\item $K_{\mathrm{raw}}\widetilde W^{-1}=\cB+Q$ with $Q=V\widetilde W^{-1}-\cB F\widetilde W^{-1}$; the last term has output in the finite set of modes $m<m_\sharp$, on which $\cB$ is analytic for $\omega$ near $\omega_0$.
\item If $\lambda$ is even, the compression of $\cB+Q$ to the modes $m\ge M\ge m_\sharp$ is $\operatorname{diag}(-m\cot(m\theta)+q_m)+A$ with the frozen sequence $q_m=q^\circ_m(\theta_0)$ of \eqref{eq:Q0} and a real matrix $A$ satisfying \eqref{eq:Aclass} with $A_1=0$. This is type \textup{(B)}.
\item If $\lambda$ is odd, with $\Lambda=\operatorname{diag}(m)$, the compression of $\Lambda^{-1}(\cB+Q)\Lambda^{-1}$ to $m\ge M$ is $\operatorname{diag}(-\cot(m\theta)/m+q_m)+A$ with the frozen sequence $q_m=q^\circ_m(\theta_0)/m^2$ and $A$ as in \textup{(iii)}. This is type \textup{(E)} with $b_E=\tan\theta_0/S_k$.
\end{enumerate}
The matrices are real and in general nonsymmetric.
\end{proposition}

\begin{proof}
(i) By \cref{lem:orders}, $W-W_0$ is a finite sum of terms $(a-a^0)\partial_y^j$ with $j\le\operatorname{ord}W_0$, primitives, translations and finite-rank terms, all with coefficient differences $O(h)$; composed with $(W_0+F)^{-1}$, whose symbol is $O(m^{-\operatorname{ord}W_0})$, \cref{lem:local} with $\mathfrak w_{\mathrm{in}}=\abs{W_0+F}^{-1}$ gives $N_\rho=O(h)$; the mismatch between input and output powers is bounded by a power of $1+\abs k$ and absorbed by $q_W$. (ii) is \eqref{eq:KVBW}. (iii) By \cref{lem:orders} and \eqref{eq:ballsymbols}, at the ball $Q$ is diagonal on $m\ge m_\sharp$ with entries $q^\circ_m(\theta)$. Away from the ball, $Q-Q_{\mathrm{ball}}=(V-V_0)\widetilde W^{-1}+V_0(\widetilde W^{-1}-(W_0+F)^{-1})-\cB F(\widetilde W^{-1}-(W_0+F)^{-1})$ has $N_\rho=O(h)$ by (i) and \cref{lem:local}, because $\operatorname{ord}V\le\operatorname{ord}W$ for even $\lambda$; the last term has finite output support and does not meet the high compression. Put $A=[Q-Q_{\mathrm{ball}}]_{\mathrm{high}}+\operatorname{diag}(q^\circ_m(\theta)-q^\circ_m(\theta_0))$. By \eqref{eq:Q0} and its $\theta$-derivative, $\abs{q^\circ_m(\theta)-q^\circ_m(\theta_0)}\le C\abs{\theta-\theta_0}\le Ch$ and $\abs{\partial_\theta q^\circ_m(\theta)}\le C$ uniformly in $m\ge m_\sharp$; the coefficient jets give $N_\rho(\partial_\omega A)=O(1)$. (iv) is identical after the congruence by $\Lambda^{-1}$, which turns $-m\cot(m\theta)$ into $-\cot(m\theta)/m$ and $q^\circ_m(\theta)$ into $q^\circ_m(\theta)/m^2=-\tan\theta/S_k+O(m^{-2})$; here $\operatorname{ord}V=\operatorname{ord}W+2$, and the two factors $\Lambda^{-1}$ compensate the extra two powers in \cref{lem:local}.
\end{proof}

\begin{remark}[Global coefficient families]\label{rem:extension}
The coefficient families $V$, $C$, $\widetilde W$ and $Q$ are given by explicit finite formulas in terms of $p$, $\chi$, $\alpha$, $\Phi$, translations, primitives and means, for \emph{all} $\omega\in I_h$ and all small $v$; they are $C^1$ in $\omega$. The rotation small divisors enter only through the diagonal symbol $-m\cot(m\theta)$ of $\cB$, which has poles at isolated frequencies. The physical equation is solved only on the retained set of \cref{sec:schur}, where all divisors are bounded below.
\end{remark}

\section{The high-mode inverse}\label{sec:schur}

In this section we prove an inverse theorem for the normalized families of \cref{def:matrixclass}. The matrices are real and not assumed to be symmetric; no spectral theorem is used.

\subsection{Hypotheses on the diagonal}
Put $L_h=\log(1/h)$, $\delta=1/\tau$ with $\tau$ from \cref{lem:circle}, and $N_h=\lfloor L_h/(4C_1)\rfloor$ with $C_1$ from (S3) below. Let $d_n(\omega)$ be the diagonal of a normalized family on $I_h$; in types (B) and (E) it is meromorphic with poles at the zeros of $\sin(n\omega/2)$. Fix $\mu>0$, with $\mu\le b_E/8$ in type (E). We consider the following hypotheses, for all $n\ge M$.

\begin{enumerate}[label=\textup{(S\arabic*)}]
\item\emph{Transversality.} (a) At nonpole $\omega\in I_h$ with $\abs{d_n(\omega)}\le2\mu$: $c_Sn^2\le\abs{\partial_\omega d_n(\omega)}\le C_Sn^2$. (b) At nonpole $\omega\in I_h$ with $\abs{d_n(\omega)}\ge\mu$: $\abs{\partial_\omega d_n(\omega)}\le C_Sn^2d_n(\omega)^2$. (c) The open set $\{\omega\in I_h:\ \omega\text{ nonpole},\ \abs{d_n(\omega)}<2\mu\}$ has at most $C_S(hn+1)$ connected components, the \emph{active cells}, and on each of them $\partial_\omega d_n$ has a constant sign.
\item\emph{Separation.} For $\omega\in\Dh$ and distinct $n,l\ge M$ with $\abs{d_n(\omega)},\abs{d_l(\omega)}<2\mu$: $\abs{n-l}\ge c_S\min(n,l)^\delta$.
\item\emph{Base values.} (a) If $d_n(\omega_0)\ne0$ then $\abs{d_n(\omega_0)}\ge c_Se^{-C_1n}$. (b) If $n\le N_h$, $\omega\in I_h$, and $\abs{d_n(\omega)}<2\mu$ or $\abs{d_n(\omega_0)}<\mu$, then the closed segment between $\omega_0$ and $\omega$ contains no pole of $d_n$, and $\abs{\partial_\omega d_n}\le C_Sn^2$ on it.
\end{enumerate}

The sign in (S1)(c) may change from one active cell to the next; only its constancy on each cell is used.

\begin{lemma}[Verification for type (O)]\label{lem:verifyO}
For odd $d$, the diagonal \textup{(O)} satisfies \textup{(S1)--(S3)} once $M$ is large, and moreover $\abs{\partial_\omega d_n}\le Cn^2$ and $\abs{\partial_\omega^2d_n}\le Cn^3$ on all of $I_h$.
\end{lemma}

\begin{proof}
By the fixed-parameter interior Darboux expansion \cite[(18.15.4\_5) and \S18.15(ii)]{DLMF}, applied to $C^\lambda_n$ and, after differentiating the exact formula by \cite[(18.9.19)]{DLMF}, to $C^{\lambda+1}_{n-1}$, uniformly for $\theta$ near $\theta_0$,
\begin{equation}\label{eq:darboux}
\begin{gathered}
d_n=\mathsf c(\theta)\,n\cos\varphi_n+O(1),\qquad \partial_\omega d_n=-\tfrac12\mathsf c(\theta)\,n^2\sin\varphi_n+O(n),\\
\varphi_n=(n+\lambda)\theta-\tfrac{\lambda\pi}2,
\end{gathered}
\end{equation}
with $\mathsf c(\theta)=2\Gamma(2\lambda)/(\Gamma(\lambda)(2\sin\theta)^\lambda)>0$. Differentiating once more with \cite[(18.9.19)]{DLMF} gives $\abs{\partial^2_\omega d_n}\le Cn^3$. Choose $C'$ with $\mathsf c(\theta)C'>2\mu+C$, where $C$ is the constant in \eqref{eq:darboux}. Outside the phase windows $\{\abs{\cos\varphi_n}\le C'/n\}$ we have $\abs{d_n}>2\mu$; on each window $\abs{\sin\varphi_n}\ge\frac12$ for $n\ge M$, so $\abs{\partial_\omega d_n}\ge\frac18\mathsf c\,n^2$ there and $\partial_\omega d_n$ has the fixed sign of $-\sin\varphi_n$, which alternates from one window to the next. Each window is an interval, the windows are disjoint, and the phase changes by $(n+\lambda)h/2$ across $I_h$, so at most $(n+\lambda)h/(2\pi)+2$ windows meet $I_h$; on each window $d_n$ is strictly monotone, so $\{\abs{d_n}<2\mu\}$ meets it in one interval. This gives (S1)(a),(c); (S1)(b) follows from $\abs{\partial_\omega d_n}\le Cn^2\le C\mu^{-2}n^2d_n^2$. For (S2), two active indices have phases within $O(1/n+1/l)$ of $\frac\pi2+\pi\Z$, so $\dist((n-l)\theta,\pi\Z)\le C(1/n+1/l)$, while on $\Dh$, $\abs{\sin((n-l)\theta)}\ge\gamma\abs{n-l}^{-\tau}$. (S3)(a) is \cref{lem:integral}, and (S3)(b) follows from the global bound on $\partial_\omega d_n$, since there are no poles.
\end{proof}

\begin{lemma}[Verification for types (B) and (E)]\label{lem:verifyBE}
For even $d$, the diagonals \textup{(B)} and \textup{(E)} of \cref{prop:evenclass} satisfy \textup{(S1)--(S3)} once $M$ exceeds $m_\sharp$ and is large. In both types $\partial_\omega d_m>0$ at every nonpole frequency.
\end{lemma}

\begin{proof}
\emph{Type (B).} Here $\partial_\omega d_m=\frac{m^2}2\csc^2(m\omega/2)>0$. If $\abs{d_m}\le2\mu$ then $\abs{m\cot(m\omega/2)}\le2\mu+\sup\abs{q_m}$, so $\abs{\cot(m\omega/2)}\le C/m$ and $\frac{m^2}2\le\partial_\omega d_m\le m^2$ for $m\ge M$. If $\abs{d_m}\ge\mu$, then $\csc^2=1+\cot^2=1+(q_m-d_m)^2/m^2\le1+2(q_m^2+d_m^2)/m^2$, whence $\partial_\omega d_m/d_m^2\le\frac{m^2}2\bigl[\mu^{-2}(1+2\sup q^2/m^2)+2/m^2\bigr]\le Cm^2$. On each cotangent cell $\{2\pi j<m\omega<2\pi(j+1)\}$ the function $d_m$ increases from $-\infty$ to $+\infty$, so $\{\abs{d_m}<2\mu\}$ meets it in one interval, and at most $mh/(2\pi)+2$ cells meet $I_h$. An active index has $m\omega$ within $C/m$ of an odd multiple of $\pi$; for two active indices $(m-n)\omega$ is within $C(1/m+1/n)$ of $2\pi\Z$, while on $\Dh$, $\abs{\sin((m-n)\omega/2)}\ge\gamma\abs{m-n}^{-\tau}$. This gives (S2).

\emph{Type (E).} Here $\partial_\omega d_m=\frac12\csc^2(m\omega/2)>0$. If $\abs{d_m}\le2\mu\le b_E/4$, then $q_m=-b_E+O(m^{-2})$ forces $\frac12b_Em\le\abs{\cot(m\omega/2)}\le2b_Em$ for $m\ge M$, hence $cm^2\le\partial_\omega d_m\le Cm^2$. If $\abs{d_m}\ge\mu$, then $\cot^2=m^2(q_m-d_m)^2\le2m^2(q_m^2+d_m^2)$ gives $\partial_\omega d_m/d_m^2\le Cm^2$. Cells are as in type (B). An active index has $m\omega$ within $C/m$ of $2\pi\Z$, which gives (S2) as before.

\emph{(S3) for both types.} By \cref{lem:orders} and \cref{lem:darboux},
\begin{align*}
d_m(\omega_0)&=-\frac{m\,[a_k\sin m\theta_0+b_k\cos m\theta_0]}{b_k\sin m\theta_0}\\
&=-\frac{m\,k_\lambda\sin^\lambda\theta_0\,C^\lambda_{m-\lambda}(\cos\theta_0)}{b_k(m)\sin m\theta_0}\qquad\text{(type (B))},
\end{align*}
and the same expression divided by $m^2$ in type (E), with $a_k,b_k$ evaluated at $(m,\cot\theta_0)$. Since $\abs{b_k(m)}\le Cm^{k}$ and $\abs{\sin m\theta_0}\le1$, \cref{lem:evenlow} gives $\abs{d_m(\omega_0)}\ge c\exp[-C\log^2(m+2)]\ge c_Se^{-C_1m}$ whenever $d_m(\omega_0)\ne0$; for $m\ge M>m_\sharp$ it is nonzero. For (S3)(b) let $m\le N_h$. In type (B), the condition $\abs{d_m}<2\mu$ at $\omega$, or $\abs{d_m}<\mu$ at $\omega_0$, gives $\abs{\sin(m\varpi/2)}\ge c_0>0$ at that endpoint $\varpi$; the phase $m\omega/2$ changes by at most $mh\le N_hh\to0$ along the segment, so $\abs{\sin}\ge c_0/2$ there and $\partial_\omega d_m\le2m^2/c_0^2$. In type (E) the same condition gives $c/m\le\abs{\sin(m\varpi/2)}$ and $\abs{\cot(m\varpi/2)}\le Cm$ at that endpoint, i.e.\ a phase distance comparable to $1/m$ from the pole; since $mh\le N_hh=o(1/N_h)$, the segment keeps a distance at least half of that, and $\partial_\omega d_m\le Cm^2$ on it.
\end{proof}

\subsection{Smooth diagonal elimination}
Let $\psi_\mu$ be a fixed smooth even function with $0\le\psi_\mu\le1$, $\psi_\mu=0$ on $\abs t\le\mu$ and $\psi_\mu=1$ on $\abs t\ge2\mu$. Put $g_\mu(t)=\psi_\mu(t)/t$ and $\nu_\mu(t)=1-\psi_\mu(t)$, and
\[
\cG=\operatorname{diag}g_\mu(d_n),\qquad\cP=\operatorname{diag}\nu_\mu(d_n),
\]
with $\cG_n=\cP_n=0$ at poles. Near a pole $\psi_\mu(d_n)=1$, so $\cG_n=1/d_n$ is a quotient of trigonometric functions with nonzero denominator and $\cP_n\equiv0$; thus $\omega\mapsto\cG_n(\omega),\cP_n(\omega)$ are $C^1$ on all of $I_h$. Moreover $\norm{\cG}\le\mu^{-1}$ and, by (S1)(a),(b),
\begin{equation}\label{eq:Gderiv}
\abs{\partial_\omega\cG_n}+\abs{\partial_\omega\cP_n}\le Cn^2 :
\end{equation}
on $\mu\le\abs{d_n}\le2\mu$ we use $\abs{\partial_\omega d_n}\le C_Sn^2$, and on $\abs{d_n}\ge2\mu$ we use $\abs{\partial_\omega d_n}/d_n^2\le C_Sn^2$.

For small $h$, $N_\rho(\cG A)<\frac12$. Set
\begin{equation}\label{eq:JC}
\begin{gathered}
\cJ=(I+\cG A)^{-1},\qquad\cC=A\cJ=(I+A\cG)^{-1}A,\\
(D+A)\cJ=D+\cP\cC,
\end{gathered}
\end{equation}
where the identity follows from $\cJ=I-\cG\cC$: the left side is $D-D\cG\cC+\cC$ and $I-D\cG=\cP$ (on finite analytic vectors, then on the analytic diagonal domain). Then $N_\rho(\cJ)\le2$ and $N_\rho(\cC)\le2Ch$. No circle condition is needed to define \eqref{eq:JC}. Define
\begin{equation}\label{eq:eff}
c_n=\cC_{nn},\qquad e_n=d_n+\nu_\mu(d_n)\,c_n .
\end{equation}
Outside the active cells $e_n=d_n$, and $\abs{e_n}\ge2\mu$ there. Differentiating \eqref{eq:JC},
\begin{equation}\label{eq:Comega}
\begin{gathered}
\partial_\omega\cC=(I-\cC\cG)\,\partial_\omega A\,(I-\cG\cC)-\cC\,(\partial_\omega\cG)\,\cC,\\
\abs{(\partial_\omega\cC)_{nn}}\le C(1+hn+h^2n^2).
\end{gathered}
\end{equation}
For the diagonal bound, the first term has entries $O(1+hn)$ since $\partial_\omega A=A_0+\Lambda_0A_1$ and an inserted $\Lambda_0$ is localized near $n$; the last term is bounded by $\sum_l(Ch)^2e^{-2\rho_\infty\abs{n-l}}C(1+l)^2\le Ch^2(1+n)^2$ by \eqref{eq:Gderiv}.

\begin{lemma}[Cellwise transversality of the effective diagonal]\label{lem:effmono}
There are $M$ and $h_0(M)$ such that for $n\ge M$, $h<h_0$ and every active cell $\mathcal O$ of $d_n$: $e_n$ is $C^1$ on $\mathcal O$, $\partial_\omega e_n$ has the sign of $\partial_\omega d_n$ on $\mathcal O$, and $\abs{\partial_\omega e_n}\ge\frac{c_S}2n^2$. Consequently, for $0<\varepsilon'\le\mu$ and $\varsigma>0$,
\begin{equation}\label{eq:smallgap}
\begin{gathered}
\Leb\{\omega\in I_h:\abs{e_n(\omega)}<\varepsilon'\}\le\frac{4C_S}{c_S}\,\varepsilon'n^{-2}(hn+1),\\
\sum_{n>L}\Leb\{\abs{e_n}<\eta n^{-\varsigma}\}\le C\eta\bigl[L^{-\varsigma-1}+hL^{-\varsigma}\bigr].
\end{gathered}
\end{equation}
\end{lemma}

\begin{proof}
On $\mathcal O$, $\partial_\omega e_n=\partial_\omega d_n\,(1+\nu_\mu'(d_n)c_n)+\nu_\mu(d_n)\partial_\omega c_n$. Here $\abs{\nu_\mu'(d_n)c_n}\le C\mu^{-1}h\le\frac14$, and by \eqref{eq:Comega} $\abs{\partial_\omega c_n}\le C(1+hn+h^2n^2)\le\frac{c_S}4n^2$ once $n\ge M$ and $h\le h_0$. Since $\abs{\partial_\omega d_n}\ge c_Sn^2$ on $\mathcal O$, we get $\abs{\partial_\omega e_n-\partial_\omega d_n}\le\frac12\abs{\partial_\omega d_n}$, which gives the sign and the lower bound. Thus $e_n$ is strictly monotone on each active cell, and $\{\abs{e_n}<\varepsilon'\}\cap\mathcal O$ is an interval of length at most $4\varepsilon'/(c_Sn^2)$; outside the cells $\abs{e_n}\ge2\mu>\varepsilon'$. The number of cells is at most $C_S(hn+1)$ by (S1)(c). Summing with $\varepsilon'=\eta n^{-\varsigma}$ gives the second bound.
\end{proof}

\begin{lemma}[Same-frequency stability]\label{lem:samefreq}
For two coefficient matrices $A_{\mathrm{old}},A_{\mathrm{new}}$ at the same $\omega$ (so $D,\cG,\cP$ coincide), $\cC_{\mathrm{new}}-\cC_{\mathrm{old}}=(I+A_{\mathrm{new}}\cG)^{-1}(A_{\mathrm{new}}-A_{\mathrm{old}})(I+\cG A_{\mathrm{old}})^{-1}$, hence
\begin{equation}\label{eq:stab}
\begin{gathered}
N_\rho(\cC_{\mathrm{new}}-\cC_{\mathrm{old}})\le CN_\rho(A_{\mathrm{new}}-A_{\mathrm{old}}),\\
\sup_n\abs{e_{n,\mathrm{new}}-e_{n,\mathrm{old}}}\le CN_\rho(A_{\mathrm{new}}-A_{\mathrm{old}}),
\end{gathered}
\end{equation}
uniformly in $n$.
\end{lemma}

\subsection{The active equation and three classes of indices}
At a nonpole frequency let $\mathcal B_\omega=\{n:\nu_\mu(d_n)\ne0\}=\{n:\abs{d_n}<2\mu\}$ be the active set and $\mathcal I_\omega$ its complement. By \eqref{eq:JC} the equation $Ku=f$ is equivalent to $(D+\cP\cC)w=f$ with $u=\cJ w$, and
\begin{equation}\label{eq:active}
d_nw_n=f_n\ (n\in\mathcal I_\omega),\qquad(D_{\mathcal B}+\cP_{\mathcal B}\cC_{\mathcal B\mathcal B})w_{\mathcal B}=f_{\mathcal B}-\cP_{\mathcal B}\cC_{\mathcal B\mathcal I}D_{\mathcal I}^{-1}f_{\mathcal I}.
\end{equation}
Every inactive $\abs{d_n}\ge2\mu$, so \eqref{eq:active}, followed by inactive reconstruction and $u=\cJ w$, is equivalent to the original equation. The active matrix has diagonal $e_n$ and off-diagonal part $E_{\mathrm{off}}$, $(E_{\mathrm{off}})_{nl}=\nu_\mu(d_n)\cC_{nl}$ for $n\ne l$, with $N_\rho(E_{\mathrm{off}})\le Ch$. With $\varpi_n=\abs{e_n}^{-1/2}$, $\varpi=\operatorname{diag}(\varpi_n)$ and $S=\operatorname{diag}(\sgn e_n)$,
\begin{equation}\label{eq:middle}
D_{\mathcal B}+\cP_{\mathcal B}\cC_{\mathcal B\mathcal B}=\varpi^{-1}\bigl[S+\varpi E_{\mathrm{off}}\varpi\bigr]\varpi^{-1},
\end{equation}
valid for real nonsymmetric $E_{\mathrm{off}}$ since $S^{-1}=S$; if the middle perturbation has norm $<1$ the Neumann series applies.

Put $\eta_h=h\exp[-L_h^\delta/(\log L_h)^4]$ and fix $\varsigma>0$. A frequency $\omega$ is \emph{retained} if $\omega\in\Dh$ and $\abs{e_n(\omega)}\ge\eta_hn^{-\varsigma}$ for all $n\ge M$. On $\Dh$ there are no poles. By \eqref{eq:smallgap} the exclusion costs $O(\eta_h(1+h))=o(h)$, and the retained set is compact, because each set $\{\abs{e_n}<\eta_hn^{-\varsigma}\}$ is relatively open: it lies in the pole-free open set $\{\abs{d_n}<\mu\}$, on which $e_n$ is continuous. For a retained $\omega$ and an active index $n$:
\begin{equation}\label{eq:classes}
\begin{aligned}
&\mathsf Z:\ n\le N_h,\ d_n(\omega_0)=0,&&\abs{e_n}\ge\tfrac{c_S}2\,hn^2;\\
&\mathsf P:\ n\le N_h,\ d_n(\omega_0)\ne0,&&\abs{e_n}\ge h^{1/2};\\
&\mathsf H:\ n>N_h,&&\abs{e_n}\ge\eta_hn^{-\varsigma}.
\end{aligned}
\end{equation}
Indeed, if $d_n(\omega_0)=0$ and $n\le N_h$, then by (S3)(b) $\abs{d_n}\le2C_Sn^2h<\mu$ on $[\omega_0,\omega_0+2h]$ for small $h$, since $hN_h^2\to0$; so this segment lies in one active cell, where by (S1) the sign of $\partial_\omega d_n$ is fixed and $\abs{\partial_\omega d_n}\ge c_Sn^2$. Hence $\abs{d_n(\omega)}\ge c_Shn^2$ on $I_h$, and $\abs{e_n}\ge c_Shn^2-2Ch\ge\frac{c_S}2hn^2$ after increasing $M$. If $d_n(\omega_0)\ne0$, $n\le N_h$ and $n$ is active, (S3) gives $\abs{d_n(\omega)}\ge c_Se^{-C_1N_h}-2C_ShN_h^2\ge\frac{c_S}2h^{1/4}$, so $\abs{e_n}\ge h^{1/2}$ for small $h$. All bounds concern absolute values.

\subsection{A smooth extension of the cotangent diagonal}
For types (B) and (E) the diagonal has poles between retained frequencies, so its parameter jets on the retained set are compared through a smooth extension. Fix an even $\zeta\in C^\infty(\R)$ with $0\le\zeta\le1$, $\zeta=0$ on $[-\frac12,\frac12]$ and $\zeta=1$ on $\abs s\ge1$, and put $F_\varepsilon(s)=\zeta(s/\varepsilon)/s$.

\begin{lemma}[Smooth reciprocal cutoff]\label{lem:cutoff}
For $m\ge1$ let $\varepsilon_m=\gamma m^{-\tau}$ and define, in type \textup{(B)},
\[
d_m^\sharp(\omega)=-m\cos(m\omega/2)\,F_{\varepsilon_m}\bigl(\sin(m\omega/2)\bigr)+q_m,
\]
and in type \textup{(E)} the same expression with $-m$ replaced by $-1/m$. Then $d^\sharp_m\in C^\infty(\R)$; $d^\sharp_m=d_m$ and $\partial_\omega d^\sharp_m=\partial_\omega d_m$ on the closed set $\{\abs{\sin(m\omega/2)}\ge\varepsilon_m\}\supset\Dh$; and
\[
\abs{d^\sharp_m-q_m}\le C\gamma^{-1}m^{\tau+1},\quad\abs{\partial_\omega d^\sharp_m}\le C\gamma^{-2}m^{2\tau+2},\quad\abs{\partial^2_\omega d^\sharp_m}\le C\gamma^{-3}m^{3\tau+3},
\]
with $C$ depending only on $\zeta$. Consequently $D^\sharp=\operatorname{diag}(d^\sharp_m)$ is a $C^2$ family of operators $\ell^2_\rho\to\ell^2_{\rho-\sigma}$ with $\norm{\partial^j_\omega(D^\sharp-q)}_{\rho\to\rho-\sigma}\le C_j\gamma^{-j-1}(1+\sigma^{-(j+1)(\tau+1)})$, $j=0,1,2$.
\end{lemma}

\begin{proof}
Since $F_\varepsilon=0$ near $0$, $d^\sharp_m$ is smooth. Where $\abs{\sin(m\omega/2)}>\varepsilon_m$ we have $\zeta=1$ in a neighborhood, so $d^\sharp_m=d_m$ there; at points with $\abs{\sin(m\omega/2)}=\varepsilon_m$ all derivatives of $\zeta(s/\varepsilon_m)$ vanish, so value and first derivative still agree. On $\Dh$, $\abs{\sin(m\omega/2)}\ge\gamma m^{-\tau}$. The bounds follow from $\abs{F_\varepsilon}\le2/\varepsilon$, $\abs{F'_\varepsilon}\le C/\varepsilon^2$, $\abs{F''_\varepsilon}\le C/\varepsilon^3$ and the chain rule, each $\omega$-derivative producing a factor $m/2$. Finally $\sup_mm^ae^{-\sigma m}\le C_a\sigma^{-a}$.
\end{proof}

In type (O) we put $D^\sharp=D$; by \cref{lem:verifyO}, $\norm{\partial_\omega^jD}_{\rho\to\rho-\sigma}\le C(1+\sigma^{-j-1})$ for $j=1,2$. In all types let $K^\sharp=D^\sharp+A$. Then $K^\sharp=K$ and $\partial_\omega K^\sharp=\partial_\omega K$ on $\Dh$.

\subsection{The inverse theorem}

\begin{theorem}[High-mode inverse]\label{thm:highinverse}
Let $K=D+A$ be a normalized family satisfying \textup{(S1)--(S3)}. There exist $M_*$ and, for each $M\ge M_*$, $h_*(M)>0$ such that for $h<h_*$ and every retained $\omega$:
\begin{enumerate}[label=\textup{(\roman*)}]
\item \textup{(Inverse.)} $K(\omega)$ is invertible on analytic sequences. For $\sigma>0$ and $\rho,\rho-\sigma\in[\rho_\infty,\rho_{\max}]$,
\begin{equation}\label{eq:Rcost}
\log\norm{K(\omega)^{-1}}_{\ell^2_\rho\to\ell^2_{\rho-\sigma}}\le C\sigma^{-2\tau}\bigl[L_h+\abs{\log\sigma}+1\bigr]^{2\tau}=:\log B(\sigma,h),
\end{equation}
and for every fixed $A_*>0$, if $\sigma\ge c(A_*\log L_h+1)^{-2}$, then $B(\sigma,h)\le Ch^{-2}\sigma^{-\varsigma}$, with $C$ depending on $A_*$.
\item \textup{(Fixed-width response.)} Let $0<\mu_0\le\rho_\infty/2$ and $s\ge0$. Put $\zeta_0=\min\{\mu_0/(24C_1),1/8\}$ and $\kappa_*=1-2\zeta_0<1$. Then
\[
\norm{K(\omega)^{-1}f}_{H^s}\le\bigl[\epsilon_Mh^{-1}+C_sh^{-\kappa_*}\bigr]\norm f_{\ell^2_{\mu_0}},
\]
where $\epsilon_M\to0$ as $M\to\infty$, for fixed $\mu_0$ and $s$.
\item \textup{(Jets.)} On the retained set put $(K^{-1})^{[1]}=-K^{-1}(\partial_\omega K)K^{-1}$. With three equal portions $\sigma/3$ of the loss,
\begin{equation}\label{eq:Rjet}
\begin{gathered}
\norm{(K^{-1})^{[1]}}_{\rho\to\rho-\sigma}\le\mathfrak g_1(\sigma)\,B(\sigma/3,h)^2,\\
\mathfrak g_1(\sigma)=\begin{cases}C(1+\sigma^{-2})&\text{type (O)},\\ C\gamma^{-2}(1+\sigma^{-2\tau-2})&\text{types (B), (E)}.\end{cases}
\end{gathered}
\end{equation}
For retained $a,b$, $K(b)^{-1}-K(a)^{-1}=-K(b)^{-1}[K^\sharp(b)-K^\sharp(a)]K(a)^{-1}$, and with five portions $\sigma/5$,
\begin{multline}\label{eq:Rwhitney}
\norm{K(b)^{-1}-K(a)^{-1}-(b-a)(K^{-1})^{[1]}(a)}_{\rho\to\rho-\sigma}\\
\le\abs{b-a}\,\bigl[\mathfrak m_A(\abs{b-a})+\abs{b-a}\bigr]\,\mathfrak g_2(\sigma)B(\sigma/5,h)^3,
\end{multline}
where $\mathfrak m_A$ is the modulus of continuity of $\partial_\omega A$ as a map into the bounded operators $\ell^2_\rho\to\ell^2_{\rho-\sigma/5}$, and $\mathfrak g_2(\sigma)=C\bigl(1+\mathfrak g_1(\sigma/5)^2+\gamma^{-3}\sigma^{-3\tau-3}\bigr)$ in types (B), (E), resp.\ $C(1+\mathfrak g_1(\sigma/5)^2+\sigma^{-3})$ in type (O). Thus $K^{-1}$ with the jet $(K^{-1})^{[1]}$ is a Whitney $C^1$ family on the retained set with values in the bounded operators $\ell^2_\rho\to\ell^2_{\rho-\sigma}$, for every $\sigma>0$.
\end{enumerate}
\end{theorem}

\subsection{Proof of the high-mode inverse theorem}
\emph{Middle factor at a fixed interior radius.} Take a radius $r\in(0,\rho_\infty)$ and work in $\ell^2_r$, so that the coefficient matrices have the margin $\rho_\infty-r$. The weights satisfy $\varpi_n\le Ch^{-1/2}/n$ on $\mathsf Z$, $\varpi_n\le h^{-1/4}$ on $\mathsf P$, and $\varpi_n\le\eta_h^{-1/2}n^{\varsigma/2}$ on $\mathsf H$. The blocks of $\varpi E_{\mathrm{off}}\varpi$ below $N_h$ are bounded by
\[
\mathsf{ZZ}:\ C/M^2,\qquad\mathsf{ZP}:\ Ch^{1/4}/M,\qquad\mathsf{PP}:\ Ch^{1/2}.
\]
For a pair $(n,l)$ of active indices with $\max(n,l)>N_h$: if $\min(n,l)\le\frac12\max(n,l)$ then $\abs{n-l}\ge N_h/2$; otherwise $\min(n,l)>N_h/2$ and (S2) gives $\abs{n-l}\ge c_S(N_h/2)^\delta$. In both cases $\abs{n-l}\ge cN_h^\delta$, and polynomial factors in $n,l$ are absorbed by a fixed portion of $e^{-(\rho_\infty-r)\abs{n-l}}$, since near-shell pairs have $n\le2l\le4n$ and $\abs{n-l}\ge c\min(n,l)^\delta$. Hence, up to fixed polynomial factors in $N_h$ and $(\rho_\infty-r)^{-1}$,
\begin{equation}\label{eq:HHblocks}
\begin{gathered}
\mathsf{ZH}:\ \sqrt{h/\eta_h}\,e^{-c(\rho_\infty-r)N_h^\delta},\qquad\mathsf{PH}:\ h^{3/4}\eta_h^{-1/2}e^{-c(\rho_\infty-r)N_h^\delta},\\
\mathsf{HH}:\ (h/\eta_h)\,e^{-c(\rho_\infty-r)N_h^\delta}.
\end{gathered}
\end{equation}
The logarithm of the largest of these is at most $L_h^\delta/(\log L_h)^4-c(\rho_\infty-r)(L_h/(4C_1))^\delta+O(\log L_h)\to-\infty$. Choose $M$ once so that $C/M^2<\frac18$. Then the perturbation in \eqref{eq:middle} has norm $<\frac12$ on $\ell^2_r$, for every $r\in(0,\rho_\infty-\sigma_0]$ with a fixed $\sigma_0>0$, and the middle matrix has a Neumann inverse of norm $\le2$.

\emph{Interior solution and fixed-width response.} This gives, through \eqref{eq:active} and $u=\cJ w$, a unique solution in $\ell^2_r$, with $w_{\mathcal B}=\varpi(S+\varpi E_{\mathrm{off}}\varpi)^{-1}\varpi(\cdots)$. Let $f\in\ell^2_{\mu_0}$. Use the middle inverse on $\ell^2_{2\mu_0/3}$; the right factor $\varpi$ maps $\ell^2_{\mu_0}$ to $\ell^2_{2\mu_0/3}$ with norm $\sup_n\varpi_ne^{-\mu_0n/3}$, and the left factor maps $\ell^2_{2\mu_0/3}$ to $\ell^2_{\mu_0/3}$, which embeds into $H^s$ with norm $C_s\mu_0^{-s}$. With the buffer $e^{-\mu_0n/3}$ the weights cost at most
\[
\mathsf Z:\ \epsilon'_Mh^{-1/2},\qquad\mathsf P:\ h^{-1/4},\qquad\mathsf H:\ \eta_h^{-1/2}\sup_{n>N_h}n^{\varsigma/2}e^{-\mu_0n/3}\le Ch^{-\frac12+\zeta_0},
\]
where $\epsilon'_M=C\sup_{n\ge M}n^{-1}e^{-\mu_0n/3}\to0$. For the last bound, $e^{-\mu_0N_h/3}\le Ch^{\mu_0/(12C_1)}$, while $\eta_h^{-1/2}N_h^{\varsigma/2}\le h^{-1/2-\mu_0/(24C_1)}$ for small $h$. Each of the two external factors is therefore bounded by $\epsilon'_Mh^{-1/2}+h^{-1/4}+Ch^{-1/2+\zeta_0}$. Their product is at most
\[
(\epsilon'^2_M+C\epsilon'_M)h^{-1}+C\bigl(h^{-1/2}+h^{-3/4+\zeta_0}+h^{-1+2\zeta_0}\bigr)\le\epsilon_Mh^{-1}+C_sh^{-\kappa_*},
\]
since every term containing a $\mathsf Z$ weight carries a factor $\epsilon'_M$, and $\max(\frac12,\frac34-\zeta_0,1-2\zeta_0)=1-2\zeta_0$ for $\zeta_0\le\frac14$. The inactive reconstruction and $\cJ$ are bounded. This proves (ii).

\emph{Analytic width loss.} Let $\sigma>0$ and $N_\sigma=\lceil C\sigma^{-2/\delta}[L_h+\abs{\log\sigma}+1]^{2/\delta}\rceil$. The tail split is made in the already formed system $D+\cP\cC$ of \eqref{eq:JC}, with the same $\cC$ and the same retained effective diagonals $e_n$. Let $w=\cJ^{-1}u$ be the interior solution, keep its components with $n\le N_\sigma$, and solve the remaining active equation for $n>N_\sigma$ with these components moved to the right side; inactive coordinates keep their diagonal reconstruction. For active $n,l>N_\sigma$ the separation is at least $cN_\sigma^\delta$ as above, so the tail middle factor has norm $<\frac12$ at the target radius $\rho-\sigma$, the coefficient margin being $\sigma$: the factor $(h/\eta_h)e^{-c\sigma N_\sigma^\delta}$ and the polynomial weights are absorbed by the choice of $N_\sigma$. Embedding the known components $n\le N_\sigma$ into $\ell^2_{\rho-\sigma}$ costs $e^{CN_\sigma}$, and the interior bound costs a power of $h^{-1}$. Uniqueness on the interior radius identifies the result with $u$. This gives \eqref{eq:Rcost}. For $\sigma\ge c(A_*\log L_h+1)^{-2}$ the middle factor is small directly at the radius $\rho-\sigma$, without the tail split: the exponent in \eqref{eq:HHblocks} with $\rho_\infty-r$ replaced by $\sigma$ is at most $L_h^\delta/(\log L_h)^4-c'L_h^\delta/(\log L_h)^2+O(\log L_h)\to-\infty$. The two external weights then cost at most $h^{-1/2}$ on $\mathsf Z\cup\mathsf P$ and $\eta_h^{-1/2}\sup_nn^{\varsigma/2}e^{-\sigma n/3}\le Ch^{-1}\sigma^{-\varsigma/2}$ on $\mathsf H$, since $\eta_h\ge h^2$ for small $h$; this gives $B(\sigma,h)\le Ch^{-2}\sigma^{-\varsigma}$.

\emph{Jets and Whitney property.} On the retained set, $K$ is differentiable with $\partial_\omega K=\partial_\omega D+A_0+\Lambda_0A_1$. In type (O), $\abs{\partial_\omega d_n}\le C(1+n)^2$ gives $\norm{\partial_\omega K}_{\rho\to\rho-\sigma}\le C(1+\sigma^{-2})$. In types (B) and (E), on $\Dh$, $\abs{\partial_\omega d_m}\le\frac{m^2}2\csc^2(m\omega/2)\le\frac12\gamma^{-2}m^{2\tau+2}$, giving $\mathfrak g_1$. The identity $(K^{-1})^{[1]}=-K^{-1}(\partial_\omega K)K^{-1}$ with three portions of the loss gives \eqref{eq:Rjet}. For retained $a,b$ the resolvent identity holds on nested scales, and $K(b)-K(a)=K^\sharp(b)-K^\sharp(a)$ because $D^\sharp=D$ on $\Dh$. Writing $\Delta_{ab}=b-a$,
\begin{multline*}
K(b)^{-1}-K(a)^{-1}-\Delta_{ab}(K^{-1})^{[1]}(a)\\
=-K(b)^{-1}\mathsf r_{ab}K(a)^{-1}-\Delta_{ab}\bigl[K(b)^{-1}-K(a)^{-1}\bigr](\partial_\omega K)(a)K(a)^{-1},
\end{multline*}
with $\mathsf r_{ab}=K^\sharp(b)-K^\sharp(a)-\Delta_{ab}\partial_\omega K^\sharp(a)$. By \cref{lem:cutoff} the diagonal part of $\mathsf r_{ab}$ is bounded by $\frac12\Delta_{ab}^2\sup\norm{\partial^2_\omega D^\sharp}$, and the part from $A$ by $\abs{\Delta_{ab}}$ times the modulus of continuity of $\partial_\omega A$. The first term is bounded by $B^2\norm{\mathsf r_{ab}}$, the second by $\Delta_{ab}^2B^3\mathfrak g_1^2$, using three, resp.\ five, portions of the loss. This gives \eqref{eq:Rwhitney}; the same argument applied to $(K^{-1})^{[1]}$ gives continuity of the jet. No indicator function of the active set is differentiated. \qed

\section{The finite physical center}\label{sec:center}

The normalized high matrix of \cref{sec:normalization} acts on the modes $n\ge M$. The remaining finite system couples the scalar variables of the physical problem with the low modes and, for even $d$, with the primitive and output constraints. In this section we keep the notation $r_1,w_1,h_2,\omega_1,a_2$ of \cref{sec:resonances}.

\subsection{The tangent equation and the frequency functional}

\begin{lemma}[Centered tangent equation]\label{lem:tangent}
Let $p=u+v$ be as in \cref{def:centered} at $\theta=\omega/2$, and let $(\mathsf k,\delta\omega)$ be a variation of $(p,\omega)$ with $\mathsf k$ odd and $\pi$-periodic. Put $p'_\pm=p'(y\pm\theta)$, $\mathsf w=\mathsf k/p'$, and let $\eta=\delta r\circ\chi$ be the induced physical radial variation, at fixed physical angle. Then
\begin{equation}\label{eq:tangent}
F_\eta:=\frac{2\chi'\,\eta}{\sec\alpha\tan\alpha\;p'_-p'_+}=\mathsf w(y+\theta)-\mathsf w(y-\theta)+\delta\omega .
\end{equation}
In particular the \emph{frequency functional} $\cR_p(\eta):=\mean_yF_\eta$ equals $\delta\omega$. Conversely, let $\omega\in\Dh$ and let $\eta$ be an even, $\pi$-periodic, analytic physical radial variation with $\cR_p(\eta)=0$; write $F_\eta=\sum_{\ell\ge2\text{ even}}f_\ell\cos(\ell y)$. Then $\mathsf k=p'\mathsf w$ with
\begin{equation}\label{eq:tangentinv}
\mathsf w=\sum_\ell\frac{f_\ell}{2\sin(\ell\theta)}\sin(\ell y)
\end{equation}
is the unique odd $\pi$-periodic variation at fixed $\omega$ inducing $\eta$, and $\norm{\mathsf w}_{\rho-\sigma}\le C\gamma^{-1}\sigma^{-\tau}\norm{F_\eta}_\rho$. Its first $\omega$-jet on $\Dh$ costs $C\gamma^{-2}\sigma^{-2\tau-1}$, and replacing $1/(2\sin(\ell\theta))$ by the smooth reciprocal cutoff of \cref{lem:cutoff} gives a $C^2$ extension in $\omega$ with the same value and first jet on $\Dh$. At the ball of radius $R$, $\cR(\eta)=2\mean_y\eta/(R\sqrt{R^2-1})$.
\end{lemma}

\begin{proof}
With $\mathsf k_\pm=\mathsf k(y\pm\theta)$ we have $\delta\chi=\frac12(\mathsf k_++\mathsf k_-)+\alpha'\delta\theta$ and $\delta\alpha=\frac12(\mathsf k_+-\mathsf k_-)+\chi'\delta\theta$. Differentiating $r(\chi(y))=\sec\alpha(y)$ and using $r'(\chi)\chi'=\sec\alpha\tan\alpha\,\alpha'$ gives $\eta=\sec\alpha\tan\alpha\,[\delta\alpha-(\alpha'/\chi')\delta\chi]$. Since $\chi'=\frac12(p'_++p'_-)$ and $\alpha'=\frac12(p'_+-p'_-)$,
\[
\delta\alpha-\frac{\alpha'}{\chi'}\delta\chi=\frac{\mathsf k_+p'_--\mathsf k_-p'_+}{2\chi'}+\frac{\chi'^2-\alpha'^2}{\chi'}\delta\theta=\frac{p'_+p'_-}{2\chi'}\bigl[\mathsf w(y+\theta)-\mathsf w(y-\theta)+2\delta\theta\bigr],
\]
which is \eqref{eq:tangent}. The mean of $\mathsf w(y+\theta)-\mathsf w(y-\theta)$ vanishes. For the converse, $\mathsf w(y+\theta)-\mathsf w(y-\theta)=\sum_\ell2\mathsf w_\ell\sin(\ell\theta)\cos(\ell y)$ for $\mathsf w=\sum\mathsf w_\ell\sin(\ell y)$; on $\Dh$ all $\sin(\ell\theta)$ are nonzero with $\abs{\sin\ell\theta}\ge\gamma\ell^{-\tau}$, which gives existence, uniqueness and the bound via $\sup_\ell\ell^\tau e^{-\sigma\ell}\le C\sigma^{-\tau}$. The $\omega$-derivative of $1/\sin(\ell\omega/2)$ is bounded by $\ell^{2\tau+1}/(2\gamma^2)$ on $\Dh$. At the ball $\chi'=p'_\pm=1$ and $\sec\theta\tan\theta=R\sqrt{R^2-1}$.
\end{proof}

At the base radius $R_0=\sqrt d$ we put $q_0=2/\sqrt{d(d-1)}$, so that $\cR=q_0\mean_y$ at the ball, and $\omega_1=q_0(r_1+a_2)$ by \cref{prop:init}.

\subsection{The physical normalized derivative}
Let the profile lie in the \emph{near-formal tube}
\begin{equation}\label{eq:tube}
\begin{gathered}
c_{\mathrm t}h\le\eps\le C_{\mathrm t}h,\qquad R_*-R_0-r_1\eps=O(h^2),\\
r=R_*+\eps Y_2+\eps^2h_2+\eps u_*,\qquad \Pi_0u_*=\Pi_2u_*=0,\qquad u_*=O(h^2),
\end{gathered}
\end{equation}
in fixed strong analytic norms (specified in \cref{sec:newton}) together with the scaled first $\omega$-jets; here $R_*=\Pi_0r$ and $\eps$ is the coefficient of $Y_2$ in $r$. Write $\Pi^\perp=I-\Pi_0$ for the projection onto output functions of zero spherical mean, and put
\[
\cF(r)=\frac{A_r-\Pi_0A_r}{\eps},\qquad\mathfrak j=Y_2+2\eps h_2+u_* .
\]
A variation of $r$ preserving the form \eqref{eq:tube} is $\delta r=\delta R_*+\mathfrak j\,\delta\eps+\eps z$ with $z=\delta u_*$, $\Pi_0z=\Pi_2z=0$. Then
\begin{equation}\label{eq:border}
\delta\cF=\Pi^\perp T_rz+a_R\,\delta R_*+a_\eps\,\delta\eps,\qquad a_R=\frac{\Pi^\perp T_r1}\eps,\qquad a_\eps=\frac{\Pi^\perp T_r\mathfrak j-\cF}\eps ,
\end{equation}
the term $-\cF/\eps$ in the amplitude column coming from the derivative of the factor $1/\eps$. At a fixed frequency the variation must satisfy the frequency compatibility $\cR_p(\delta r)=0$ of \cref{lem:tangent}, i.e.
\begin{equation}\label{eq:freqrow}
\eps\,\cR_p(z)+\cR_p(1)\,\delta R_*+\cR_p(\mathfrak j)\,\delta\eps=0 ;
\end{equation}
the factor $\eps$ in front of $\cR_p(z)$ is the physical amplitude. The \emph{physical normalized derivative} with target $g$, $\Pi_0g=0$, is the system \eqref{eq:border} $=g$, \eqref{eq:freqrow}, for the unknowns $(z,\delta R_*,\delta\eps)$.

\begin{lemma}[Leading scalar structure]\label{lem:scalar}
In the tube \eqref{eq:tube}, $\cF=O(h^2)$, and uniformly
\begin{gather*}
a_R=L_2'Y_2+O(h),\qquad a_\eps=-r_1L_2'Y_2+O(h),\\
\cR_p(1)=q_0+O(h),\qquad\cR_p(\mathfrak j)=q_0a_2+O(h),
\end{gather*}
in fixed analytic norms. In particular the components of $a_R$ and $a_\eps$ of degree $\ge4$ are $O(h)$, and
\begin{equation}\label{eq:scalarmatrix}
\det\begin{pmatrix}L_2'&-r_1L_2'\\ q_0&q_0a_2\end{pmatrix}=L_2'\,q_0(r_1+a_2)=L_2'\,\omega_1\ne0 .
\end{equation}
\end{lemma}

\begin{proof}
By \cref{lem:taylor}, with $f=r-R_*=\eps Y_2+\eps^2h_2+\eps u_*$ and $T_{R_*+f}g=L_{R_*}g+L'_{R_*}(fg)+O(\norm f^2\norm g)$,
\begin{align*}
\Pi^\perp A_r&=\eps L_2(R_*)Y_2+\eps^2L_4a_{4,d}Y_4+\tfrac{\eps^2}2\Pi^\perp L'(Y_2^2)+O(h^3)\\
&=\eps^2\bigl[L_2'(r_1+\tfrac{c_{22}}2)Y_2+(L_4a_{4,d}+\tfrac12L_4'c_{24})Y_4\bigr]+O(h^3),
\end{align*}
using $L_2(R_*)=L_2'(R_*-R_0)+O(h^2)$ and \eqref{eq:Y2sq}; both brackets vanish by \cref{def:r1}. Hence $\cF=O(h^2)$. Next $\Pi^\perp T_r1=\eps L_2'Y_2+O(h^2)$, and
\begin{align*}
\Pi^\perp T_r\mathfrak j&=L_2(R_*)Y_2+\eps\,\Pi^\perp L'(Y_2^2)+2\eps L_4h_2+O(h^2)\\
&=\eps\bigl[L_2'(r_1+c_{22})Y_2+(L_4'c_{24}+2L_4a_{4,d})Y_4\bigr]+O(h^2),
\end{align*}
whose $Y_4$ coefficient vanishes and whose $Y_2$ coefficient is $-r_1L_2'\eps$, since $c_{22}=-2r_1$. With $\cF/\eps=O(h)$ this gives $a_R$ and $a_\eps$. Finally $\cR_p\to q_0\mean_y$ and $\mean_y\mathfrak j\to\mean_yY_2(\cos y)=a_2$.
\end{proof}

The scalar column $a_\eps$ has no degree-four component at leading order because of the equation defining $h_2$; this is used to show that the high parts of the scalar columns are $O(h)$.

\subsection{Normalized coordinates, constraints and dimension count}
To combine \eqref{eq:border}--\eqref{eq:freqrow} with the high matrix we write the system in the normalized native coordinates of \cref{sec:normalization}, and we replace the projected output equation by the equivalent unprojected equation
\begin{equation}\label{eq:unproj}
T_rz+a_R\,\delta R_*+a_\eps\,\delta\eps-\varkappa=g,
\end{equation}
with an additional scalar unknown $\varkappa$; for $\Pi_0g=0$ its spherical mean gives $\varkappa=\Pi_0T_rz$, and its projection is \eqref{eq:border}. The constraints $\Pi_0z=\Pi_2z=0$ are kept as two linear rows; $\Pi_0$ and $\Pi_2$ are evaluated natively by the spherical measure $(\sin\chi)^{d-2}\chi'\,dy$.

\emph{Odd $d$.} The unknowns are the normalized input $F=a_{\mathrm{in}}\tilde z$ (all degrees) and the three scalars $(\delta R_*,\delta\eps,\varkappa)$; the equations are the normalized output of \eqref{eq:unproj} (all degrees), the frequency row \eqref{eq:freqrow} and the two projection rows. Splitting the harmonic coordinates at the degree $M$, the \emph{finite center} consists of the degrees $n<M$ and the three scalars, and the rows of output degree $n<M$ together with the three rows; it is square.

\emph{Even $d$.} The unknown is the native primitive $v$ of $z$. By \eqref{eq:physrec} below, a primitive determines the physical variation, and two primitives give the same variation exactly when they differ by a combination of the $\lfloor(\lambda-1)/2\rfloor$ modes $\sin(j\chi)$, $j\le k$ of the parity of $\lambda$; we add the corresponding primitive rows, which require these sine coefficients of the physical primitive to vanish. They are the native integrals $\int v(y)\sin(j\chi(y))\chi'(y)\,dy$. On the output side, the regularized output has no cosine modes $\cos(l\phi)$, $l\le k$ of the parity of $\lambda$ (and $l\ne0$ in the periodic case); in native coordinates these are the functionals $\ell_l(\mathsf g)=\int[\Delta\mathsf g](y)\cos(l\Phi(y))\,dy$, for which $\ell_l(\sin(m\,\cdot))=-2\pi\sin(l\theta)\delta_{lm}$ at the ball; only the finitely many indices $1\le l\le k$ occur, and for each of them $\sin(l\theta_0)\ne0$ because $\theta_0/\pi$ is irrational (\cref{sec:circlearith}); since $\abs{\theta-\theta_0}=O(h)$ on $I_h$, continuity gives $\abs{\sin(l\theta)}\ge\frac12\min_{1\le l\le k}\abs{\sin(l\theta_0)}>0$ for all these $l$, uniformly for small $h$. Every transformed derivative of a physical input lies in the common kernel of these functionals, since the identity of \cref{prop:pullthrough} was derived from the physical derivative. We append $\lfloor(\lambda-1)/2\rfloor$ columns along a fixed low Fourier basis dual to the $\ell_l$; applying the $\ell_l$ to a solved equation with physical source forces the appended multipliers to vanish. The numbers of primitive rows and of appended columns are both $\lfloor(\lambda-1)/2\rfloor$, so the finite center, formed as in the odd case after the input change $\widetilde W$ (and the congruence by $\Lambda^{-1}$ if $\lambda$ is odd), is again square, and the high matrix is unchanged. The physical radial variation is recovered from the primitive by
\begin{equation}\label{eq:physrec}
z\circ\chi=\cos(\alpha)^{2\lambda}\Bigl[\frac{-1}{\chi'\sin\chi}\partial_y\Bigr]^k\Bigl[\frac v{k_\lambda\sin\chi}\Bigr],
\end{equation}
whose divisions by $\sin\chi$ are removable on the parity class of $v$ (finitely many removable divisions and derivatives).

\subsection{The finite matrix}
Fix $M$ above all lower degrees used above, above $m_\sharp$ in the even case, and large enough for \cref{thm:highinverse}. Split the finite variables into the extra resonances $\mathsf Z_0$ below $M$ (degrees $n<M$ with $C^\lambda_n(1/\sqrt d)=0$, $n\ge8$) and the primary block $\mathsf P_0$ (all other finite variables). On harmonic coordinates use the same exponential weight $e^{rn}$ as for the high system, at a fixed interior radius $r$; scalar coordinates have fixed weights. The finite matrix has the block form
\begin{equation}\label{eq:Fh}
\begin{gathered}
F_h=\begin{pmatrix}F_{\mathsf{PP}}&F_{\mathsf{PZ}}\\F_{\mathsf{ZP}}&F_{\mathsf{ZZ}}\end{pmatrix},\qquad F_{\mathsf{PP}}=F^0_{\mathsf{PP}}+O_{d,M}(h),\\
F_{\mathsf{PZ}},\,F_{\mathsf{ZP}}=O_{d,M}(h),\qquad F_{\mathsf{ZZ}}=\eps\operatorname{diag}(\bar\delta_n)+O_{d,M}(h^2),
\end{gathered}
\end{equation}
where $\bar\delta_n$ is the detuning $\delta_n$ of \cref{lem:kernel} multiplied by the fixed normalization factors at the ball.

\begin{lemma}\label{lem:Fhstructure}
The matrix $F^0_{\mathsf{PP}}$ is invertible, \eqref{eq:Fh} holds, and $\sup_n\abs{\bar\delta_n}^{-1}\le C_d$, the supremum being over all resonant degrees and $C_d$ independent of $M$.
\end{lemma}

\begin{proof}
At the ball limit $h\to0$ along \eqref{eq:tube}, in physical coordinates and the order: projection rows, output degrees $n\ge4$ nonresonant, output degree two, frequency row, output degree zero, the primary block is block triangular: the projection rows give $1$ on the input degrees $0$ and $2$; a nonresonant input degree $n$ has output only in degree $n$, with the nonzero entry $L_n$; the degree-two row and the frequency row restricted to $(\delta R_*,\delta\eps)$ form the matrix \eqref{eq:scalarmatrix}, while the input harmonics enter them only with $L_2=0$ or with the factor $\eps$; the degree-zero output row has the entry $-1$ in the column $\varkappa$, since $a_R,a_\eps$ have zero mean. In the even case the primitive rows and appended columns have the nonzero diagonal pairings computed above. The finite changes from physical to normalized native coordinates are at the ball diagonal with nonzero entries on each retained mode (weights $a_{\mathrm{in}},a_{\mathrm{out}}$, the factors $\Lambda_0^{\lambda+1}/\kappa_d$, the primitive $S_\lambda$, the multipliers of $P_\lambda$ for $m\ge\lambda$, the denominators $\sin(m\theta_0)$ and $W_0+F$), and no change of coordinates divides by a resonant multiplier. Hence $F^0_{\mathsf{PP}}$ is invertible. Both cross blocks vanish at the ball: a resonant input has zero round output and zero projection rows, and the frequency row carries $\eps$; conversely the scalar columns have only degree-two base output, the column $\varkappa$ is constant, and the other low harmonics are diagonal. The first-order part of $F_{\mathsf{ZZ}}$ in the direction $w_1$ is diagonal with entries $\eps\delta_n$ (\cref{lem:kernel}); the tube \eqref{eq:tube} gives the remainder $O(h^2)$. Passing from the ideal to the actually compressed native block changes $F_h$ by uniformly conditioned left and right factors and by $O(h^2)$: a transformed low basis vector has a low part which is an $O(h)$ perturbation of its base image and a high tail of size $O(h)$, and dropping the tail changes each finite row by $O(h^2)$, since all high-to-finite rows vanish at the ball and are $O(h)$. Finally, in the odd case, $\bar\delta_n=(1+n)^{\lambda+1}L_n'(r_1+c_{nn})/\kappa_d(R_0)$ up to the fixed weights, and at a zero of $d_n(\omega_0)$ the identity $\partial_\theta Y_n(\cos\theta)=-\sin\theta\,Y_n'(\cos\theta)$ and (S1)(a) give $\abs{\bar\delta_n}\ge cn^2(r_1+c_{nn})\ge c'n^2$ for $n\ge M_0$; the finitely many lower resonances have $\bar\delta_n\ne0$. In the even case there are finitely many resonances.
\end{proof}

\begin{proposition}[Smallest singular value]\label{prop:sigmamin}
For each fixed $M$ and small $h$, $\norm{F_h^{-1}}\le C_{d,M}h^{-1}$, equivalently $\sigma_{\min}(F_h)\ge c_{d,M}h$. Moreover $\limsup_{h\to0}h\norm{F_h^{-1}}\le C_d$ with $C_d$ independent of $M$.
\end{proposition}

\begin{proof}
Let $S_\eps=\operatorname{diag}(I_{\mathsf P},\eps^{-1}I_{\mathsf Z})$. Then
\[
S_\eps F_h=\begin{pmatrix}F^0_{\mathsf{PP}}+O(h)&O(h)\\ \eps^{-1}F_{\mathsf{ZP}}&\operatorname{diag}(\bar\delta_n)+O(h)\end{pmatrix},\qquad\eps^{-1}F_{\mathsf{ZP}}=O_{d,M}(1).
\]
The Schur complement of the upper left block is $\Sigma=\operatorname{diag}(\bar\delta_n)+O(h)-\eps^{-1}F_{\mathsf{ZP}}F_{\mathsf{PP}}^{-1}F_{\mathsf{PZ}}=\operatorname{diag}(\bar\delta_n)+O_{d,M}(h)$, which is invertible with $\norm{\Sigma^{-1}}\le2C_d$ for small $h$ by \cref{lem:Fhstructure}. The block formula
\[
(S_\eps F_h)^{-1}=\begin{pmatrix}F_{\mathsf{PP}}^{-1}+F_{\mathsf{PP}}^{-1}F_{\mathsf{PZ}}\Sigma^{-1}\eps^{-1}F_{\mathsf{ZP}}F_{\mathsf{PP}}^{-1}&-F_{\mathsf{PP}}^{-1}F_{\mathsf{PZ}}\Sigma^{-1}\\-\Sigma^{-1}\eps^{-1}F_{\mathsf{ZP}}F_{\mathsf{PP}}^{-1}&\Sigma^{-1}\end{pmatrix}
\]
has norm $\le C_{d,M}$, and $F_h^{-1}=(S_\eps F_h)^{-1}S_\eps$ gives the first claim since $\eps\ge c_{\mathrm t}h$. For the second, $hF_h^{-1}=(S_\eps F_h)^{-1}\operatorname{diag}(hI_{\mathsf P},(h/\eps)I_{\mathsf Z})$; for fixed $M$, as $h\to0$ the first block column tends to $0$, and in the second block column the upper entry is $O_{d,M}(h)$ while the lower one is $\Sigma^{-1}(h/\eps)$ with norm $\le2C_d/c_{\mathrm t}$. The diagonal exponential weights do not change diagonal blocks. This is a bound on the operator norm; the conditioning of the coordinate changes was established in \cref{lem:Fhstructure}.
\end{proof}

\subsection{Feedback from the high modes and the full border}
Write the full normalized system as $\begin{psmallmatrix}K_{\mathsf{HH}}&K_{\mathsf{HC}}\\K_{\mathsf{CH}}&F_h\end{psmallmatrix}$, with $K_{\mathsf{HH}}$ the high matrix of \cref{thm:highinverse} and $\mathsf C$ the finite center. The columns $K_{\mathsf{HC}}$ and rows $K_{\mathsf{CH}}$ vanish at the ball: for the amplitude column this uses the equation for $h_2$ (\cref{lem:scalar}); extra resonant contributions in higher formal terms are annihilated by the round multiplier; the projection rows vanish on high modes at the ball; and the frequency row carries $\eps$. In the weighted coordinates fixed above, the harmonic cross blocks have weighted row and column norms $O(h)$ by \eqref{eq:Aclass}, with constants independent of $M$; the high parts of the scalar columns are $O(h)$ analytic vectors with source radius above $r$; the projection rows are analytic dual vectors of size $O(h)$ on high modes; and the frequency row, whose coefficient on a normalized high mode is at most $C\eps(1+n)^{\lambda}$, is exponentially decaying after the weight $e^{-rn}$ and hence $O(h)$.

\begin{proposition}[Invertibility of the full border]\label{prop:border}
Let the profile lie in the tube \eqref{eq:tube} and let $\omega$ be retained. For $M$ fixed large enough and then $h$ small, the full normalized system is invertible. The Schur complement $\Sigma_h=F_h-K_{\mathsf{CH}}K_{\mathsf{HH}}^{-1}K_{\mathsf{HC}}$ satisfies $\norm{\Sigma_h^{-1}}\le2(C_d+1)h^{-1}$ on the fixed interior radius, and the inverse is given by
\[
x_{\mathsf C}=\Sigma_h^{-1}\bigl(f_{\mathsf C}-K_{\mathsf{CH}}K_{\mathsf{HH}}^{-1}f_{\mathsf H}\bigr),\qquad x_{\mathsf H}=K_{\mathsf{HH}}^{-1}\bigl(f_{\mathsf H}-K_{\mathsf{HC}}x_{\mathsf C}\bigr).
\]
For a width loss $\sigma$ its norm is at most $C_{d,M}h^{-1}\sigma^{-q}\,[1+B(\sigma/c,h)]^2$ for fixed $q,c$.
\end{proposition}

\begin{proof}
By \eqref{eq:JC}--\eqref{eq:middle}, $K_{\mathsf{HH}}^{-1}=\cJ[\text{inactive diagonal inverse}+\varpi(S+\varpi E_{\mathrm{off}}\varpi)^{-1}\varpi]$ up to the bounded inactive correction of \eqref{eq:active}. The inactive part contributes $O(h^2)$ to the feedback $K_{\mathsf{CH}}K_{\mathsf{HH}}^{-1}K_{\mathsf{HC}}$. For the active part, the cross blocks multiplied by one external weight $\varpi$ have norms at most $C_dh^{1/2}/M$ on $\mathsf Z$ (where $\varpi_n\le Ch^{-1/2}/n$ and $n\ge M$), $C_dh^{3/4}$ on $\mathsf P$, and $C_{d,M}h\,\eta_h^{-1/2}e^{-c(N_h-M)}=o_{d,M}(h^{1/2})$ on $\mathsf H$, using the exponential decay of the coupling between modes below $M$ and above $N_h$. With the middle inverse of norm $\le2$,
\[
\norm{K_{\mathsf{CH}}K_{\mathsf{HH}}^{-1}K_{\mathsf{HC}}}\le C_dh/M^2+o_{d,M}(h).
\]
By \cref{prop:sigmamin}, $\norm{F_h^{-1}}\le(C_d+o_{d,M}(1))h^{-1}$. Choose $M$ so that $C_d^2/M^2\le\frac14$ and then $h$ small; then $\norm{F_h^{-1}K_{\mathsf{CH}}K_{\mathsf{HH}}^{-1}K_{\mathsf{HC}}}\le\frac12$ and $\Sigma_h$ is invertible by a Neumann series with the stated bound. The displayed formulas give the inverse. For a target loss $\sigma$ both $K_{\mathsf{HH}}^{-1}$ factors cost $B(\sigma/c,h)$, the finite conversions cost fixed powers of $\sigma^{-1}$, and $\Sigma_h^{-1}$ costs $h^{-1}$; the inverse may contain two high inverse factors. The fixed-width response of \cref{thm:highinverse}(ii) gives the alternative bound $C_d(\epsilon_M+h^{1-\kappa_*})$ for the relative feedback when the center is measured with polynomial weights.
\end{proof}

\section{The native Newton iteration and proof of the theorem for \texorpdfstring{$d\ge5$}{d>=5}}\label{sec:newton}

\subsection{Order of choices}
Fix $d\ge5$, the base angle $\theta_0$, the exponent $\tau_0$ of \cref{lem:circle}, $\tau>\tau_0+1$, $\gamma\in(0,c_d/2]$ and $\delta=1/\tau$, and fix analytic radii $0<\rho_\infty<\rho_0\le\rho_{\max}$. Fix the finite list of coefficient fields of \cref{lem:ledger} below. Choose the high cutoff $M$ large enough for \cref{lem:effmono}, the middle-factor estimate of \cref{thm:highinverse} and \cref{prop:border}, and use this compression at every step of the iteration. Then fix the cost exponents of \cref{lem:costs}, the early-window constant $a_*$, the exponents $K\ge8$ and $d_*>(\varsigma+1)\log2$, a sufficiently high finite initializer order $N$ in \cref{prop:init}, and finally $h$, which may depend on $M$, $N$ and all finite coefficients.

\subsection{Spaces and the field ledger}
Let $\Wn^q_\rho$ be the space of the introduction (real, odd, $\pi$-periodic, norm $\norm\cdot_{\rho,q}$) and $\Wn^q_{\rho,\mathrm e}$ the analogous space of real, even, $\pi$-periodic functions. The tracked shape is $v=p-u\in\Wn^{q_d}_\rho$ ($p$ itself is not periodic); the residual is even and $\pi$-periodic. We use three norm levels:
\begin{gather*}
\text{residual: }\Wn^0_{\rho,\mathrm e};\qquad\text{shape and corrections: }\Wn^{q_d}_\rho;\\
\text{native tube complement: }\Wn^{t_d}_{\rho,\mathrm e},
\end{gather*}
with $t_d=8(d+1)$, $m_d=16(d+1)$, $q_d=32(d+1)$.

\begin{lemma}[Field ledger]\label{lem:ledger}
There is a finite list of coefficient fields (the weights, the coefficients of the normalized operators of \cref{sec:normalization}, the scalar columns and rows of \cref{sec:center}, the native projections and the volume map), each an analytic expression in $v$ and its angular derivatives of order at most $s_*(d)$, where
\[
s_*(d)=4r_d+2\lceil\lambda\rceil+24\quad(d\text{ odd},\ r_d=\lceil\lambda\rceil+8),\qquad s_*(d)=4d+8\quad(d\text{ even}),
\]
such that the following holds: the fields and their first explicit $\omega$-derivatives, and their shape derivatives up to order three, are bounded at the same radius by the norm $\norm v_{\rho,m_d}$, and $\max\{t_d+1,s_*(d)\}<m_d$.
\end{lemma}

\begin{proof}
\emph{Odd $d$.} By \cref{prop:Oclass} the coefficient norm with weight $s_d=\lceil\lambda\rceil+5$ of $\mathfrak E$ is controlled by the mixed angular derivatives of order $r_d$ in each variable. The construction of $\mathfrak E$ (\cref{prop:oddnorm}(c)) uses $G$ through order three; the mixed derivatives add at most $2r_d$ and one explicit frequency derivative one more, so derivatives of $G$ through order $2r_d+4$ suffice; an axial derivative of order $a$ costs at most $(1+n)^{2a}$ in the angular coefficient norm (repeated Chebyshev differentiation), so $4r_d+8$ angular derivatives of $v$ suffice for $\mathfrak E$. The weights, the projections, the volume amplitude $\Omega$ and $S_G$, and the scalar columns, including the noncommuting placement of the multipliers and the factor $\Lambda_0^{\lambda+1}$, use at most $2\lceil\lambda\rceil+16$ further angular powers. The tube complement \eqref{eq:unative} below is a separate field using at most $t_d$ derivatives, $t_d+1$ with an explicit frequency derivative. Since $\lceil\lambda\rceil\le d/2$, $s_*(d)=6\lceil\lambda\rceil+56\le3d+56<16(d+1)$ for $d\ge5$; for $d=5$, $s_*=68<96$.

\emph{Even $d$.} With $k=\lambda-1$: the $j$-th axial derivative of $r$ along the profile uses at most $j+1$ angular derivatives and $j$ removable divisions by $\sin\chi$, each costing one more derivative (\cref{lem:sinquot}); the endpoint coefficients of $E_\pm$ (\cref{lem:endpoint}, \cref{lem:orders}) and the residual coefficients $g_i$ involve $H^{(i)}$, $i\le k$, and $\partial_s^i\mathcal Q$, $i\le k$, hence at most $2k+2$ derivatives, and the removability of \cref{lem:meridian} at most $2k+1$ further removable divisions; the pull-through coefficients $M_J,Y_J$ add translations only. This gives order $\le4k+4=4\lambda$ for $V,C,W$. \Cref{lem:local} adds at most $\lambda+3$ powers for the order mismatch and the congruence by $\Lambda^{-1}$, and an explicit frequency derivative adds one (translations by $\pm\theta$). The scalar columns, applied through the same normalization to the primitives of $1$ and $\mathfrak j$, and the volume map \eqref{eq:slice} with its removable division by $\sin^{2\lambda+1}$, use fewer than $8\lambda+16=4d+8$ derivatives. Shape derivatives replace factors by variations and do not increase the maximal order. Finally $4d+8<16(d+1)$ and $t_d+1=8d+9<16(d+1)$.
\end{proof}

\emph{Norm conversions.} For odd $d$ the high system lives on the orthonormal Gegenbauer coefficients $\ell^2_{\lambda,\rho}$ in the variable $Y$. The map $\Wn^0_{\rho,\mathrm e}\to\ell^2_{\lambda,\rho}$ is bounded, because $T_k$ has degree $k$ and bounded $L^2_\lambda$ norm. Conversely an orthonormal Gegenbauer polynomial of degree $n$ has real supremum norm $O_d((1+n)^\lambda)$ and at most $n+1$ Chebyshev coefficients, each bounded by twice that norm, so Cauchy--Schwarz with a buffer $\sigma$ gives
\begin{equation}\label{eq:conv}
\norm f_{\rho-\sigma,q_d}\le C_d\bigl(1+\sigma^{-q_d-d-4}\bigr)\norm f_{\ell^2_{\lambda,\rho}} .
\end{equation}
For even $d$ the high system lives on the native sine coefficients, $\Wn^0_\rho\to\ell^2_\rho$ is bounded, and $\norm f_{\rho-\sigma,q_d}\le C\sigma^{-q_d-1}\norm f_{\ell^2_\rho}$. These are bounded linear maps; they act on first jets and Whitney remainders, and each use spends a positive buffer.

\emph{Same-width Lipschitz bound.} At fixed $\omega$ and for $\rho\ge\rho_\infty$,
\begin{equation}\label{eq:samewidth}
\begin{aligned}
N_\rho\bigl(A(v_{\mathrm{new}},\omega)-A(v_{\mathrm{old}},\omega)\bigr)&\le C_d\norm{v_{\mathrm{new}}-v_{\mathrm{old}}}_{\rho,m_d}\\
&\le C_d\norm{v_{\mathrm{new}}-v_{\mathrm{old}}}_{\rho,q_d},
\end{aligned}
\end{equation}
by \cref{prop:Oclass}, resp.\ \cref{lem:local} and \cref{prop:evenclass}; the frozen diagonal part of $A$ in the even case does not depend on $v$. This uses the coefficient norms at the same radius and not the frequency derivative of $A$, which grows with the mode in type (O).

\subsection{Native tube, residual and transport}
For each $p$ use the centered $\chi,\alpha$ of \cref{def:centered}. Define $R_*$ and $\eps$ by the physical degree-zero and degree-two projections of $r$, evaluated natively with the measure $(\sin\chi)^{d-2}\chi'\,dy$ and the fixed round normalizing constants, and the native complement
\begin{equation}\label{eq:unative}
u_{\mathrm n}(y)=\frac{\sec\alpha(y)-R_*-\eps Y_2(\cos\chi(y))-\eps^2h_2(\cos\chi(y))}{\eps}.
\end{equation}
Its $\Wn^{t_d}_{\rho,\mathrm e}$ norm is tracked directly, and $\Pi_0u_*=\Pi_2u_*=0$ holds exactly for $u_*=u_{\mathrm n}\circ\chi^{-1}$. With one fixed buffer for the identification of the charts, \eqref{eq:unative} gives the tube \eqref{eq:tube} needed in \cref{sec:center}. The coefficient fields and scalar columns are computed directly from $p$.

The residual is
\begin{equation}\label{eq:resid}
\cF(p,\omega)(u)=\frac{\cV_p(u)-\bar A_p}{\eps},\qquad
\bar A_p=\mathsf w_d\int_0^\pi\cV_p(u)(\sin p(u))^{d-2}p'(u)\,du,
\end{equation}
where $\mathsf w_d^{-1}=\int_0^\pi(\sin u)^{d-2}\,du$ and $\cV_p(u)=A_r(p(u))$ is the volume map: \eqref{eq:Qp} for odd $d$, and for even $d$ the slice formula \eqref{eq:slice} in native coordinates,
\begin{equation}\label{eq:evenQ}
A_r(\Phi(y))=-\kappa_{2\lambda}(\sin\Phi(y))^{-2\lambda-1}\int_y^{y+\omega}\Disc^\lambda\,\mathcal X'\,d\tilde y,
\end{equation}
with $\mathcal X=\sec\alpha\cos\chi$, whose singularity at $\sin\Phi=0$ is removable (\cref{sec:sections}). The mean $\bar A_p$ is the spherical mean of the physical output, pulled back by $p$. If $\mathsf k$ is a variation of $p$ at fixed $\omega$ and $\eta$ the induced physical radial variation, the chain rule gives
\begin{equation}\label{eq:transport}
D_p\cF[\mathsf k]=(\delta_{\mathrm{phys}}\cF)[\eta]\circ p+\frac{\partial_u\cF}{p'}\,\mathsf k ,
\end{equation}
where $\delta_{\mathrm{phys}}\cF$ is the physical normalized differential of \eqref{eq:border}. The second term, the \emph{transport term}, is proportional to the residual.

\subsection{Whitney norms and one global extension}
Let $\mathsf E\subset I_h$ be compact and $f$ a Banach-valued family on $\mathsf E$ with Whitney first jets $f^{[1]}$. Put
\begin{equation}\label{eq:Wnorm}
\norm f_{\mathrm{Wh}(\mathsf E)}=\sup_{\mathsf E}\norm f+h\sup_{\mathsf E}\norm{f^{[1]}}+h\sup_{a\ne b\in\mathsf E}\frac{\norm{f(b)-f(a)}}{\abs{b-a}} .
\end{equation}
The Whitney remainder property is an additional property, not implied by \eqref{eq:Wnorm}. Extend $f$ once: on each bounded gap $(a,b)$ of $\mathsf E$, with $\ell=b-a$ and $s=(\omega-a)/\ell$, put
\begin{multline*}
f^{\mathrm{ext}}(\omega)=(2s^3-3s^2+1)f(a)+(3s^2-2s^3)f(b)\\
+\ell(s^3-2s^2+s)f^{[1]}(a)+\ell(s^3-s^2)f^{[1]}(b),
\end{multline*}
and on the outer gaps use the affine extension from the endpoint. The derivative is a fixed polynomial combination of $f^{[1]}(a)$, $f^{[1]}(b)$ and $(f(b)-f(a))/\ell$, so
\begin{equation}\label{eq:hermite}
\sup_{I_h}\norm{f^{\mathrm{ext}}}+h\sup_{I_h}\norm{\partial_\omega f^{\mathrm{ext}}}\le C\norm f_{\mathrm{Wh}(\mathsf E)} .
\end{equation}
On shrinking gaps the derivative differs from $f^{[1]}(a)$ by a constant times $\norm{f(b)-f(a)-\ell f^{[1]}(a)}/\ell+\norm{f^{[1]}(b)-f^{[1]}(a)}$, which tends to zero by the Whitney property; hence $f^{\mathrm{ext}}$ is $C^1$ on $I_h$ in the Banach space, including at accumulation points of $\mathsf E$. Reality and odd $\pi$-periodicity are preserved by linearity.

\subsection{The correction step}
Let $\mathsf E_j$ be the retained set at stage $j$, $\rho_{j+1}=\rho_j-\sigma_j$ with $\sigma_j=c(j+1)^{-2}$ and $\rho_0-\sum_j\sigma_j=\rho_\infty$, and $e_j=\norm{\cF_j}_{\mathrm{Wh}(\mathsf E_j)}$ in $\Wn^0_{\rho_j,\mathrm e}$. Split $\sigma_j$ into twelve equal blocks. On $\mathsf E_j$:
\begin{enumerate}[label=(\arabic*)]
\item convert the native source $-\cF_j$ to normalized coordinates (one block);
\item solve the full normalized system of \cref{prop:border}, including the residual subtraction in the amplitude column and the factor $\eps$ in the frequency row, and form its first jet by the resolvent identity (two blocks, subdivided a fixed number of times for the at most two high inverse factors and their jets; in the even case this includes the cutoff losses of \cref{lem:cutoff});
\item return to physical radial data, by \eqref{eq:conv} and, for even $d$, by \eqref{eq:physrec} (one block); the frequency row gives $\cR_p(\eta_j)=0$ for $\eta_j=\eps z_j+\delta R_{*,j}+\mathfrak j\,\delta\eps_j$;
\item obtain the odd $\pi$-periodic correction $\mathsf k_j$ from \cref{lem:tangent} (one block, including the cutoff of its divisors for the jet);
\item upgrade $\mathsf k_j$ and its jet to $\Wn^{q_d}$ with a Cauchy buffer (one block); repeating this at later stages again produces $\Wn^{q_d}$;
\item use the remaining six blocks for the Cauchy derivative in the transport term, the Taylor and jet estimates, and the finitely many norm conversions.
\end{enumerate}
Then extend $\mathsf k_j$ once by \eqref{eq:hermite} (no loss of width), replace $p$ by $p+\mathsf k_j^{\mathrm{ext}}$ on all of $I_h$, and recompute every field, $R_*$, $\eps$, the physical mean, the normalization, the scalar columns and the projections from this single $p$. The correction need not solve the equation on the gaps; the centered construction gives an actual body at every $\omega\in I_h$.

\begin{proposition}[Quadratic recurrence]\label{prop:recurrence}
There are constants $\Gamma_j$, depending on $h$, $\sigma_j$ and the inverse bounds, such that
\begin{equation}\label{eq:rec}
e_{j+1}\le\Gamma_je_j^2\ \ \text{on }\mathsf E_j,\qquad\sup_{I_h}\norm{\Delta_j\mathsf f}+h\sup_{I_h}\norm{\partial_\omega\Delta_j\mathsf f}\le\Gamma_je_j
\end{equation}
for the increment $\Delta_j\mathsf f$ of every tracked field $\mathsf f$ of \cref{lem:ledger}, of the shape in $\Wn^{q_d}_{\rho_{j+1}}$, and of the tube complement.
\end{proposition}

\begin{proof}
On $\mathsf E_j$ the physical normalized differential applied to $\eta_j$ equals $-\cF_j$ by construction, so by \eqref{eq:transport} $D_p\cF[\mathsf k_j]=-\cF_j+(\partial_u\cF_j/p')\mathsf k_j$, and
\[
\cF(p+\mathsf k_j)=\frac{\partial_u\cF_j}{p'}\,\mathsf k_j+\mathcal T_2(p;\mathsf k_j),
\]
where $\mathcal T_2$ is the second-order Taylor remainder. By the Cauchy estimate on one block, $\norm{\partial_u\cF_j}\le Cb_j^{-1}e_j$ with $b_j=\sigma_j/12$, and $\norm{\mathsf k_j}\le\Gamma'_je_j$ with $\Gamma'_j$ the product of the inverse and conversion costs of steps (1)--(5); hence the transport term is quadratic in $e_j$. The exact nonlinear map has a common analytic domain (\cref{lem:fixedray}, \cref{lem:physmap}, \eqref{eq:evenQ}); products, nonzero scalar inverses, analytic powers, translations, fixed $\beta$-integrals, removable divided differences and removable divisions by $\sin$ give bounded second shape derivatives after positive buffers, with fixed powers of $\eps^{-1}\le C h^{-1}$; so $\norm{\mathcal T_2}\le Ch^{-P}b_j^{-P}\norm{\mathsf k_j}^2$. Three shape derivatives and one explicit frequency derivative suffice for the first-jet and divided-difference parts of \eqref{eq:Wnorm}, using products and the resolvent identities; the Whitney remainders follow from \eqref{eq:Rwhitney}, \cref{lem:cutoff} and \cref{lem:tangent}. The extension agrees with values and jets on $\mathsf E_j$, and restriction to a smaller set does not increase the norm. No smallness of the residual on the gaps is used. The increments of the fields are controlled by \eqref{eq:hermite} and the same finite calculus.
\end{proof}

\begin{lemma}[Cost classes]\label{lem:costs}
There are fixed $P_0,P$ such that, with $L=\log(1/h)$, for every $a_*>0$ and $j_*=\lceil a_*\log L\rceil$,
\begin{equation}\label{eq:costs}
\begin{gathered}
\log\Gamma_j\le P_0L+P_0\log(j+2)+C_{a_*}\quad(j\le j_*),\\
\log\Gamma_j\le C(L+j+1)^P\quad(\text{all }j\ge0).
\end{gathered}
\end{equation}
\end{lemma}

\begin{proof}
The inverse and border costs are those of \eqref{eq:Rcost}, \eqref{eq:Rjet}, \eqref{eq:Rwhitney} and \cref{prop:border}, evaluated at a fixed fraction of $\sigma_j$. For $j\le j_*$ we have $\sigma_j\ge c(a_*\log L+1)^{-2}$, so the early bound $B\le Ch^{-2}\sigma_j^{-\varsigma}$ of \cref{thm:highinverse}(i) applies. For all $j$, $\log B(\sigma_j/c,h)\le C(j+1)^{4\tau}(L+\log(j+2)+1)^{2\tau}\le C(L+j+1)^{6\tau}$. The factors $\mathfrak g_1,\mathfrak g_2$ of \eqref{eq:Rjet}--\eqref{eq:Rwhitney}, which in the even case contain $\gamma^{-4}\sigma_j^{-4\tau-4}$ from the cutoff of the cotangent diagonal, the tangent divisor costs $\gamma^{-2}\sigma_j^{-2\tau-1}$, the conversions \eqref{eq:conv}, the normalization order $\lambda+1$, the indices $m_d,q_d$, the powers of $\eps^{-1}$ and products of at most a fixed number of inverse bounds only enlarge $P_0$ and $P$; $\gamma$ is fixed.
\end{proof}

\begin{proposition}[Two-regime induction]\label{prop:induction}
Let $K\ge8$ and $d_*>0$, and choose $a_*$ with $a_*\log2>P-1$. There are a finite initializer order $N$ and $h_*>0$ such that for $h<h_*$, for all $j\ge0$,
\begin{equation}\label{eq:target}
\Gamma_je_j\le h^K\exp[-d_*(j+1)] .
\end{equation}
In particular $e_j\le h^{c_\star2^j}$ with $c_\star$ as large as required; if some $e_j$ vanishes, the subsequent corrections are taken to be zero.
\end{proposition}

\begin{proof}
Put $a_j=\log\Gamma_j$ and $b_j=-\log e_j$. From \eqref{eq:rec}, $b_{j+1}\ge2b_j-a_j$, whence $b_j\ge2^j\bigl[b_0-\sum_{i<j}2^{-i-1}a_i\bigr]$. By \eqref{eq:costs} the part of the sum with $i\le j_*$ is at most $P_0L+C$, and the part with $i>j_*$ is at most $C2^{-j_*}(L+j_*+1)^P\le CL^{P-a_*\log2}(1+o(1))=o(L)$. By \cref{prop:init} the initial residual satisfies $e_0\le C_Nh^N$, including its scaled jet and divided differences, so $b_0\ge NL-C_N'$. Choose $N\ge P_0+c_\star+1$; then $b_j\ge c_\star2^jL$ for all $j$ and small $h$. For $j\le j_*$, $a_j-b_j\le P_0L+P_0\log(j+2)+C-c_\star2^jL\le-KL-d_*(j+1)$ if $c_\star>P_0+K+1$. For $j>j_*$, the function $j\mapsto(L+j+1)^P/(2^jL)$ is decreasing once $L+j+1>P/\log2$, and at $j=j_*$ it is $O(L^{P-1-a_*\log2})\to0$; hence $a_j-b_j\le-\frac12c_\star2^jL\le-KL-d_*(j+1)$ for all $j>j_*$ and small $h$.
\end{proof}

We take $K\ge8$ and $d_*>(\varsigma+1)\log2$. The summed corrections are $O(h^K)$ in $C^1_h$ and preserve all native-field, branch and convexity margins and $\eps\asymp h$. The tube \eqref{eq:tube}, including its scaled jet, is preserved: the field $u_{\mathrm n}$ changes by $O(h^{K-1})$, the division by $\eps$ costing one power of $h$. The high scalar columns are initialized directly on the finite native initializer, where their high parts are $O(h)$ on a fixed source strip by \cref{lem:scalar} (with the cancellation coming from $h_2$); they belong to the field list, so their summed changes are $O(h^K)$. Hence the finite center and the feedback remain in the range of \cref{prop:sigmamin} and \cref{prop:border} at every step.

\subsection{Nested exclusions}
Let $\eta_h=h\exp[-L^\delta/(\log L)^4]$ and fix $\varsigma>0$. Start with $\mathsf E_0=\Dh\cap\{\abs{e_{n,0}}\ge2\eta_hn^{-\varsigma}\ \forall n\ge M\}$; by \eqref{eq:smallgap} and \cref{lem:circle}, $\Leb(I_h\setminus\mathsf E_0)=o_d(h)$. Put $\eta_{h,j}=\eta_h(1+2^{-j})$ and $N_j=N_\star2^j$ with a fixed $N_\star\ge M$. Given $\mathsf E_j$, update $p$ on $I_h$ and define
\[
\mathsf E_{j+1}=\mathsf E_j\cap\{\omega:\ \abs{e_{n,j+1}(\omega)}\ge\eta_{h,j+1}n^{-\varsigma}\ \text{for all }n\ge M\},
\]
where $e_{n,j+1}$ is formed from the new coefficient family with the same compression and the same smooth elimination \eqref{eq:JC}. By \cref{lem:samefreq}, \eqref{eq:samewidth}, \eqref{eq:rec} and \eqref{eq:target},
\begin{equation}\label{eq:stabj}
\sup_{n\ge M}\abs{e_{n,j+1}(\omega)-e_{n,j}(\omega)}\le C_dh^Ke^{-d_*(j+1)}\qquad(\omega\in\mathsf E_j),
\end{equation}
with no factor depending on $n$ (the frequency derivatives are not used here). For $n\le N_{j+1}$ the reserve $(\eta_{h,j}-\eta_{h,j+1})n^{-\varsigma}\ge\eta_hN_\star^{-\varsigma}2^{-(\varsigma+1)(j+1)}$ exceeds \eqref{eq:stabj} for every $j$, because $d_*>(\varsigma+1)\log2$ and $h^K/\eta_h\to0$. So only indices $n>N_{j+1}$ can be newly deleted. The new coefficients are defined by the new $p$ on all of $I_h$, including the old gaps, and satisfy the hypotheses of \cref{lem:effmono}; therefore the newly deleted part of $\mathsf E_j$ has measure at most
\begin{equation}\label{eq:newloss}
C_d\,\eta_h\bigl[N_{j+1}^{-\varsigma-1}+hN_{j+1}^{-\varsigma}\bigr],
\end{equation}
which is summable in $j$. The sets $\mathsf E_j$ are compact and nested (the conditions are closed relative to $\Dh$, which contains no pole). With $E_{d,h}=\bigcap_j\mathsf E_j$,
\begin{equation}\label{eq:measure}
\Leb(E_{d,h})\ge h-o_d(h).
\end{equation}
At every fixed mode $n$, the coefficients converge and $\eta_{h,j}\downarrow\eta_h$, so the limiting gap $\eta_hn^{-\varsigma}$ persists.

\subsection{The exact limit and the volume image}
The corrections $\mathsf k_j^{\mathrm{ext}}$ and their first frequency derivatives are summable in $\Wn^{q_d}_{\rho_\infty}$, the latter being bounded by $h^{-1}$ times \eqref{eq:target}. Hence $v_\omega=p_\infty-u$ is a $C^1$ map from $I_h$ to $\Wn^{q_d}_{\rho_\infty}$, small in $\norm\cdot_{\rho_\infty,1}$. By \cref{lem:centered}, for every $\omega\in I_h$ the body $K_\omega=K(v_\omega,\omega)$ is an origin-symmetric, strictly convex body of revolution with real-analytic boundary strictly containing $B^d$, whose radial function is $O(h)$-close to $\sqrt d$ on a fixed complex strip. Its physical degree-two coefficient stays comparable to $h$, so $K_\omega$ is not a ball. On $E_{d,h}$ the residual \eqref{eq:resid} tends to zero in $\Wn^0_{\rho_\infty,\mathrm e}$, and the volume map is continuous on the common domain, so $A_r-\Pi_0A_r=0$: all distance-one sections of $K_\omega$ have the volume $V(\omega)=\bar A_{p_\infty(\omega)}$.

The mean $V(\omega)=\bar A_{p_\infty(\omega)}$ is a bounded native integral of the $C^1$ family, so $V\in C^1(I_h)$. By \cref{prop:init}, the initializer has $V'(\omega)=-2\abs{S^{d-2}}(d-1)^{(d-2)/2}+O_d(h)$, and the summed derivative changes are $O_d(h^{K-1})$. For small $h$, $V'$ is negative and bounded away from zero on $I_h$, so $V$ is bi-Lipschitz there and $\Leb(V(E_{d,h}))\ge c\Leb(E_{d,h})>0$ by \eqref{eq:measure}.

\begin{proof}[Proof of \cref{thm:dge5}]
Items (i)--(iv) were established above; the measure statement follows from \eqref{eq:measure} and the bi-Lipschitz property of $V$.
\end{proof}

\begin{remark}[The case $d=5$]\label{rem:d5}
For odd $d$ the arguments of \cref{part:d5} use only $\lambda\ge\frac32$ and inequalities verified for every odd $d\ge5$. For $d=5$ the constants are
\begin{gather*}
\lambda=\tfrac32,\qquad \beta=\tfrac12,\qquad R_0=\sqrt5,\qquad r_1=-\tfrac16,\qquad c_{nn}=\frac{3n(n+3)}{2(4n^2+12n+5)},\\
a_{4,5}=-\frac{5\sqrt5}{72},\qquad Y_4(1/\sqrt5)=-\frac3{25},\qquad Y_6(1/\sqrt5)=\frac{11}{125},\\
L_2'=-4\pi^2,\qquad \omega_1=\frac{\sqrt5}{24},\qquad V'(\omega_0)=-32\pi^2;
\end{gather*}
in \cref{lem:integral}, $m_*=1$ and $c=2$; the remainder of \cref{prop:oddnorm} is built from $T_{7,R}$, with coarea exponent $\lambda+\frac12=2$ and $\Disc_0^{3/2}$; the constant of \cref{lem:physmap} is $\kappa_3/(\abs{S^4}R^5)=1/(2\pi R^5)$; and the ledger is $r_5=10$, $s_5=7$, $t_5=48$, $m_5=96$, $q_5=192$ with $s_*(5)=68<96$.
\end{remark}

\section*{Use of AI-assisted tools}
AI-assisted tools were used in the research and preparation of this manuscript.

\begin{thebibliography}{99}

\bibitem{BW} A. Baker and G. W\"ustholz, \emph{Logarithmic forms and group varieties}, J. reine angew. Math. \textbf{442} (1993), 19--62.

\bibitem{BL} J. A. Barker and D. G. Larman, \emph{Determination of convex bodies by certain sets of sectional volumes}, Discrete Math. \textbf{241} (2001), 79--96.

\bibitem{DLMF} NIST Digital Library of Mathematical Functions, Chapter~18: Orthogonal Polynomials, \url{https://dlmf.nist.gov/18}, accessed October 2026.

\bibitem{FHL} J.-Ll. Figueras, A. Haro, and A. Luque, \emph{Effective bounds for the measure of rotations}, Nonlinearity \textbf{33} (2020), 700--741. Equation numbers refer to the author manuscript dated 3 October 2019, \url{https://hdl.handle.net/2445/192675}.

\bibitem{FRSY} A. Fish, D. Ryabogin, W. R. Sun, and V. Yaskin, \emph{On determination of convex bodies by their section functions}, preprint (2026), arXiv:2610.10229v1.

\bibitem{Had} J. Haddad, \emph{An approximate counterexample to the Barker--Larman problem in dimension 4}, preprint (2026), arXiv:2609.10257v2.

\bibitem{HR} J. Haddad and D. Ryabogin, \emph{On convex bodies with constant non-central sections}, preprint (2026), arXiv:2605.00299v1.

\bibitem{KSdiv} Knowledge Sense, \emph{The formal Barker--Larman branch in dimension four has zero radius of convergence}, companion manuscript (2026).

\bibitem{Laz} V. F. Lazutkin, \emph{The existence of caustics for a billiard problem in a convex domain}, Math. USSR-Izv. \textbf{7} (1973), no.~1, 185--214.

\bibitem{OAI1} OpenAI, \emph{Continuous phase foliations create analytic caustic collars}, manuscript dated 24 September 2026, \url{https://github.com/openai/math/blob/fd4aeeb2ee4fc729c18d98444fed42fd0529eeeb/CONTENTS.md}.

\bibitem{OAI2} OpenAI, \emph{Rigidity of smooth billiards with a continuous caustic collar}, manuscript dated 24 September 2026, \url{https://github.com/openai/math/blob/fd4aeeb2ee4fc729c18d98444fed42fd0529eeeb/CONTENTS.md}.

\bibitem{YZ} V. Yaskin and N. Zhang, \emph{Non-central sections of convex bodies}, preprint (2015), arXiv:1509.08174v1.

\end{thebibliography}

\end{document}
